\documentclass[11pt,a4paper,openany]{book}
\usepackage[a4paper,margin=2.1cm,top=2.5cm,bottom=2.3cm,headheight=22pt]{geometry}
\usepackage{haese}
\renewcommand{\guidename}{Infi workbook Ch.\,0--7}
\renewcommand{\coursename}{Infinitesimaalrekening}
\renewenvironment{bookex}[2][from the dictaat]{\begin{bookexbox}[{#1}]{#2}}{\end{bookexbox}}
\hypersetup{pdftitle={Infinitesimaalrekening - Workbook Ch. 0-7},pdfauthor={Workbook}}
\begin{document}
\thispagestyle{empty}
\begin{tcolorbox}[enhanced, colback=hblue, colframe=hblue, boxrule=0pt, arc=4pt, left=16pt, right=16pt, top=22pt, bottom=22pt]
  \color{white}\sffamily
  {\large\bfseries Infinitesimaalrekening \textperiodcentered\ Blok 1 \textperiodcentered\ 2026--2027}\\[10pt]
  {\fontsize{34}{38}\selectfont\bfseries Workbook}\\[6pt]
  {\fontsize{20}{24}\selectfont Chapters 0--7}\\[14pt]
  {\normalsize Minimal theory, taken from the dictaat's own definitions, theorems and strategies, followed immediately by the dictaat's own exercises: every Oefening of Hoofdstuk 0--7 (1--30, 101--123, 201--220, 301--330, 401--420, 501--524, 601--635, 701--730) is either worked in the text or solved in Appendix S.}
\end{tcolorbox}

\vspace{10pt}
\begin{tcolorbox}[enhanced, colback=white, colframe=hblue2, boxrule=0.8pt, arc=2pt, left=10pt, right=10pt, top=8pt, bottom=8pt]
\sffamily\small\sloppy
\begin{itemize}[leftmargin=1.2em]
\item \textbf{\textcolor{hblue}{Source.}} S.~Wepster, \emph{Infinitesimaalrekening}, editie 2026, Mathematisch Instituut, Universiteit Utrecht: Blok 1, Hoofdstuk 0--7 (pages 17--113). Definition, Stelling, Voorbeeld and Oefening numbers are the dictaat's; nothing is used that the dictaat does not contain or take for granted (vwo algebra and, in the marked exercises, basic calculus).
\item \textbf{\textcolor{hblue}{Which exercises are worked where.}} Roughly half of the exercises, chosen to show each technique once, are worked in full in the chapter text. All others appear in NOW YOU boxes and are solved in Appendix S. Exercises the dictaat marks \emph{(Extra)}, and the (Extra) sections 3.5, 6.6 and 7.1, are outside the exam material but included.
\item \textbf{\textcolor{hblue}{How to use it.}} Read one theory box, cover the solution of the worked exercise that follows, attempt it, then compare. Do the NOW YOU exercises before opening Appendix S. Reflect after every exercise, as the dictaat asks: could you have seen the answer coming, would another approach work, what if the question were changed slightly?
\end{itemize}
\end{tcolorbox}

\vspace{8pt}
\hsub{How this workbook works}
Every section is a sequence of short cycles: \textbf{one piece of theory, then the exercises that use it.}

\begin{center}\sffamily\small
\begin{tabular}{@{}p{0.30\linewidth}p{0.65\linewidth}@{}}
\colorbox{hyellow}{\strut\textbf{Definitie n.m}} & The dictaat's definitions and notation, condensed; the things it lists under ``hoofdpunten: uit je hoofd weten''.\\[3pt]
\colorbox{hlight}{\strut\textbf{Stelling / Voorbeeld}} & A result you may use, with the idea behind it, or a worked example from the dictaat that serves as the model for the exercises.\\[3pt]
\colorbox{hgreen}{\strut\textbf{Strategy}} & The dictaat's methods (separating variables, integrating factor, substitution, partial fractions, function investigation) as step-by-step recipes.\\[3pt]
\colorbox{white}{\strut\textcolor{hblue}{\textbf{Exercise n}} \textcolor{horange}{\small\textbf{badge}}} & An Oefening worked in full, with the checks the dictaat insists on (differentiate your primitive, substitute your solution, sanity-check a value). Cover the solution and try first.\\[3pt]
\colorbox{hpurplebg}{\strut\textbf{NOW YOU}} & The remaining exercises for that piece of theory, all solved in Appendix S.\\[3pt]
\colorbox{hredbg}{\strut\textbf{Watch out}} & Where marks are lost: the dictaat's own warnings and the errors its author sees most in exams.
\end{tabular}
\end{center}

\begin{key}[{Conventions of the dictaat}]
\begin{itemize}
\item $\log$ is the \textbf{natural} logarithm (base $e$); $\ln$ is never used. $\N=\{0,1,2,\dots\}$.
\item The inverse function of $f$ is written $f_{\mathrm{inv}}$, not $f^{-1}$ (to avoid confusion with $1/f$); likewise $\arcsin,\arccos,\arctan$ and $\operatorname{arsinh},\operatorname{arcosh},\operatorname{artanh}$.
\item ``The function $f(x)$'' is sloppy but used; $f:A\to B$, $x\mapsto f(x)$ is correct. Intervals $(a,b)$ open, $[a,b]$ closed; $\infty$ is never a number.
\item Exercises are numbered per chapter: Oefening 203 is the third exercise of Hoofdstuk 2. The dictaat marks easy theory-checks and ``switching'' problems with icons that are not reproduced here; (Extra) marks exercises outside the exam material.
\item Exam: no calculator, no formula sheet; sum and double-angle formulas (Oefening 26) and the standard primitives (\S6.1) must be known by heart.
\end{itemize}
\end{key}

\clearpage
\thispagestyle{plain}
{\sffamily\bfseries\Large\color{hblue}Contents}\par\vspace{4pt}
\makeatletter\@starttoc{toc}\makeatother
\setcounter{chapter}{-1}
\chapter{Voorkennis: prior knowledge refreshed}
\chaptercontents{A & Algebraic skills (\S0.1) \\ B & Functions and inverse functions (\S0.2) \\ C & Polynomials and factorising (\S0.3) \\ D & Trigonometry (\S0.4) \\ E & The logarithm (\S0.5)}%
{Exercises 1--30 (all worked or solved).\\ The dictaat's own ``hoofdpunten'': when an inverse function exists and how to find it; $|x|$ and $\sqrt{x^2}=|x|$; composition $f\circ g$; even and odd; the factor theorem; the unit circle, exact values at ``nice'' angles, symmetries, sum formulas (needed hard in Chapter 6); basic properties of $\log$ and $\exp$.}

\hsection{Algebraic skills}

\begin{strategy}[{\S0.1 \textperiodcentered\ the point of these exercises}]
These exercises refresh vwo algebra. The dictaat insists that you \textbf{find your own way to check an answer}: substitute back, multiply out again, try a special value (``sanity check''). The one identity to recognise instantly: $a^2-b^2=(a-b)(a+b)$.
\end{strategy}

\begin{bookex}{2}
Simplify: (a) $(2-\sqrt3)^2$; (b) $(2-\sqrt3)(2+\sqrt3)$; (c) $\dfrac{2025^2-1}{2024}$; (d) $\dfrac{(2+\sqrt\pi-y)^2-(2-\sqrt\pi+y)^2}{\sqrt\pi-y}$ (choose a handy approach!).
\solution
(a) $(2-\sqrt3)^2=4-4\sqrt3+3=7-4\sqrt3$.\quad (b) $(2-\sqrt3)(2+\sqrt3)=4-3=1$.\\
(c) $2025^2-1=(2025-1)(2025+1)=2024\cdot2026$, so the quotient is $2026$.\\
(d) Put $b=\sqrt\pi-y$; then the numerator is $(2+b)^2-(2-b)^2=\big((2+b)-(2-b)\big)\big((2+b)+(2-b)\big)=2b\cdot4=8b$, and dividing by $b$ gives $8$ (for $y\ne\sqrt\pi$). Expanding both squares would give the same, with ten times the work.\qed
\end{bookex}

\begin{bookex}{4}
(a) For which $r$ does $2x^2+rx+6=0$ have exactly one solution $x$? Answer in the form $r=n\sqrt m$, $n\in\Z$, $m\in\N$ minimal. (b) $\dfrac{3t}{t+\frac{1}{1+1/t}}=1$ has exactly one solution; which? (c) Find all $u$ with $\dfrac{1}{1+\frac1{1+u}}<2$. (d) Solve $\sqrt v+1=\sqrt{v+1}$ for $v\in\R$.
\solution
(a) Exactly one solution iff the discriminant is $0$: $r^2-4\cdot2\cdot6=r^2-48=0$, so $r=\pm\sqrt{48}=\pm4\sqrt3$. Both $r=4\sqrt3$ and $r=-4\sqrt3$ work ($n=\pm4$, $m=3$).\\
(b) The fraction only makes sense for $t\ne0$ and $t\ne-1$. Simplify from the inside: $\dfrac1{1+1/t}=\dfrac{t}{t+1}$, so the denominator is $t+\dfrac t{t+1}=\dfrac{t(t+1)+t}{t+1}=\dfrac{t(t+2)}{t+1}$. The equation becomes $\dfrac{3t(t+1)}{t(t+2)}=1$, i.e.\ (as $t\ne0$) $3(t+1)=t+2$, so $t=-\tfrac12$. Check: $1+1/t=-1$, $\frac1{-1}=-1$, $t+(-1)=-\frac32$, and $\frac{-3/2}{-3/2}=1$. \checkmark\\
(c) Need $u\ne-1$ and $1+\frac1{1+u}\ne0$, i.e.\ $u\ne-2$. Then $1+\dfrac1{1+u}=\dfrac{u+2}{u+1}$ and the inequality is $\dfrac{u+1}{u+2}<2$. Never multiply by $u+2$ without knowing its sign; instead bring everything to one side: $\dfrac{u+1}{u+2}-2=\dfrac{u+1-2u-4}{u+2}=\dfrac{-(u+3)}{u+2}<0$, i.e.\ $\dfrac{u+3}{u+2}>0$: both factors positive ($u>-2$) or both negative ($u<-3$). Removing $u=-1$: $u\in(-\infty,-3)\cup(-2,-1)\cup(-1,\infty)$.\\
(d) Both sides are defined for $v\ge0$ and non-negative, so squaring is reversible: $v+2\sqrt v+1=v+1$, so $\sqrt v=0$, $v=0$. Check: $\sqrt0+1=1=\sqrt1$. \checkmark\qed
\end{bookex}

\begin{bookex}{5}
Find all $c$ with: (a) $(3+c)(c^2+2)=(2c+2)(3+c)$; (b) $1-c=\sqrt{1-2c+c^2}$; (c) $2^{c^2}\cdot2^{2c+1}=16$. (d) Check that $c=-3$ is a solution of all three and $c=3$ of none.
\solution
(a) Do \emph{not} divide by $3+c$ (you would lose a solution). Bring to one side and factor: $(3+c)\big(c^2+2-2c-2\big)=(3+c)(c^2-2c)=c(c-2)(c+3)=0$, so $c\in\{-3,0,2\}$.\\
(b) $1-2c+c^2=(1-c)^2$, so the right side is $\sqrt{(1-c)^2}=|1-c|$. The equation $1-c=|1-c|$ holds exactly when $1-c\ge0$: all $c\le1$.\\
(c) $2^{c^2+2c+1}=2^4$, so $c^2+2c+1=(c+1)^2=4$, $c+1=\pm2$: $c=1$ or $c=-3$.\\
(d) $-3$ appears in (a), satisfies $-3\le1$ in (b), and appears in (c); $3$ is not in (a), fails $3\le1$, and $3^2+6+1=16\ne4$. \checkmark\qed
\end{bookex}

\begin{trybox}[{1, 3}]
\textbf{1} All positive divisors of $60,61,\dots,65$; which are squares?\quad \textbf{3} Check $\frac16-\frac17=\frac1{6\cdot7}$; for which $p,q$ is $\frac1p-\frac1q=\frac1{pq}$?
\end{trybox}

\hsection{Functions and inverse functions}

\begin{key}[{Definition 0.1 \textperiodcentered\ function, domain, codomain, range}]
A \term{function} $f:A\to B$ is a rule assigning to every $x\in A$ exactly one $y=f(x)\in B$, written $x\mapsto f(x)$; $y$ is the \term{image} (beeld) of $x$ and $x$ an \term{original} of $y$. $A$ is the \term{domain}, $B$ the \term{codomain}, and the \term{range} (bereik) is the set of $y\in B$ that actually occur as images. If no domain is given, take the largest subset of $\R$ on which the formula makes sense. Voorbeeld 0.2: $x\mapsto x^2$ on $\R$ has range $[0,\infty)$; $5$ has the two originals $\pm\sqrt5$. Voorbeeld 0.3: ``$f:\R\to\R$, $x\mapsto\sqrt x$'' is wrong (negative $x$); $f:[0,\infty)\to\R$ is right. Strictly, ``the function $f(x)$'' is wrong ($f(x)$ is a number), but the dictaat uses it anyway. \textbf{Piecewise} definitions (Voorbeeld 0.4): $\operatorname{sgn}x=-1,0,+1$ for $x<0$, $x=0$, $x>0$; in graphs a closed dot is on the graph, an open dot is not.
\end{key}
\begin{key}[{Definition 0.5 and Voorbeeld 0.7 \textperiodcentered\ inverse function}]
If $f:A\to B$ is \term{bijective} (a one-to-one pairing of $A$ with $B$), its \term{inverse} $f_{\mathrm{inv}}:B\to A$ is the unique function with $f_{\mathrm{inv}}(f(x))=x$. The dictaat writes $f_{\mathrm{inv}}$, not $f^{-1}$, to avoid confusion with $1/f$. $y=f(x)\Leftrightarrow x=f_{\mathrm{inv}}(y)$; the graphs are mirror images in the line $y=x$. Voorbeeld 0.6: $x\mapsto(x-1)^2$ is not bijective on $\R\to\R$ ($f(-1)=f(3)=4$), nor on $[1,\infty)\to\R$ (negative numbers are never images), but it is bijective $[1,\infty)\to[0,\infty)$ with inverse $y\mapsto1+\sqrt y$. \textbf{To find the inverse:} write $y=f(x)$ and solve for $x$; then swap the letters.
\end{key}

\begin{bookex}{9}
Find the inverse of the bijective functions (a) $p:x\mapsto1+\sqrt[3]{2x}$; (b) $q:x\mapsto\dfrac{x^2-4}{x+2}$; (c) $r:x\mapsto\dfrac{3x}{4+5x}$, and state on which domain and range each is bijective.
\solution
(a) $y=1+\sqrt[3]{2x}\Leftrightarrow(y-1)^3=2x\Leftrightarrow x=\tfrac12(y-1)^3$. So $p_{\mathrm{inv}}(x)=\tfrac12(x-1)^3$; the cube root is defined for all reals, so $p:\R\to\R$ is bijective (domain $\R$, range $\R$).\\
(b) For $x\ne-2$: $q(x)=\dfrac{(x-2)(x+2)}{x+2}=x-2$. Domain $\R\setminus\{-2\}$; the value $-2-2=-4$ is never taken, so the range is $\R\setminus\{-4\}$. From $y=x-2$: $q_{\mathrm{inv}}(x)=x+2$ on $\R\setminus\{-4\}$, with range $\R\setminus\{-2\}$.\\
(c) Domain $x\ne-\frac45$. $y=\dfrac{3x}{4+5x}\Leftrightarrow4y+5xy=3x\Leftrightarrow x(3-5y)=4y\Leftrightarrow x=\dfrac{4y}{3-5y}$, provided $y\ne\frac35$ (and $y=\frac35$ would need $3x=\frac35(4+5x)$, i.e.\ $0=\frac{12}5$: impossible, so $\frac35$ is not in the range). Hence $r:\R\setminus\{-\frac45\}\to\R\setminus\{\frac35\}$ is bijective with $r_{\mathrm{inv}}(x)=\dfrac{4x}{3-5x}$.\qed
\end{bookex}

\begin{bookex}{10}
Let $f$ be bijective with inverse $f_{\mathrm{inv}}$ and let $s(x)=1-2f(3-4x)$. Show how to compute $s_{\mathrm{inv}}$ from $f_{\mathrm{inv}}$.
\solution
Put $y=s(x)=1-2f(3-4x)$ and peel the layers off from the outside: $f(3-4x)=\dfrac{1-y}2$, so $3-4x=f_{\mathrm{inv}}\!\Big(\dfrac{1-y}2\Big)$, so $x=\dfrac{3-f_{\mathrm{inv}}\big(\frac{1-y}2\big)}4$. Hence $s_{\mathrm{inv}}(y)=\dfrac14\Big(3-f_{\mathrm{inv}}\Big(\dfrac{1-y}2\Big)\Big)$: the inverse undoes the steps $x\mapsto3-4x\mapsto f(\cdot)\mapsto1-2(\cdot)$ in reverse order.\qed
\end{bookex}

\begin{trybox}[{6, 7, 8}]
\textbf{6} Sketch the graphs of $x\mapsto x^2$ and $x\mapsto\sqrt x$ by hand.\quad \textbf{7} Give a piecewise formula for the Heaviside function.\quad \textbf{8} Fill in the steps of Voorbeeld 0.7 (inverse of $f(x)=3-\frac1{2+x}$).
\end{trybox}

\hsub{Absolute value}
\begin{key}[{\S0.2.1}]
$|x|=x$ for $x\ge0$ and $-x$ for $x<0$; equivalently $|x|=\sqrt{x^2}$. \textbf{Distance:} $|x-y|$ is the distance between $x$ and $y$ on the number line, so $|x-1|<3$ means ``less than $3$ away from $1$''. Recognise $|x|$ when it is hidden, as in $\sqrt{x^2}$ or $\sqrt{1-\cos^2x}$.
\end{key}

\begin{bookex}{11}
Solve $|x-1|\le|x-3|$.
\solution
\emph{Distance:} the points no further from $1$ than from $3$ are those on the side of the midpoint $2$ nearer to $1$: $x\le2$.\\
\emph{Algebra:} both sides are $\ge0$, so squaring is allowed: $x^2-2x+1\le x^2-6x+9\Leftrightarrow4x\le8\Leftrightarrow x\le2$.\qed
\end{bookex}

\begin{trybox}[{12}]
\textbf{12} Simplify $\sqrt{1-\cos^2x}$ as far as possible; sanity check at $x=-\frac\pi6$.
\end{trybox}

\hsub{Compositions, even and odd functions}
\begin{key}[{\S0.2.2 and Definition 0.8}]
$(f\circ g)(x)=f(g(x))$ (``$f$ after $g$'': apply $g$ first). $f$ is \term{even} if $f(x)=f(-x)$ for all $x$ in its domain, \term{odd} if $f(x)=-f(-x)$. Even: symmetric in the $y$-axis ($x^2$, $\cos$, $|x|$); odd: symmetric about the origin ($x$, $x^3$, $\sin$).
\end{key}

\begin{bookex}{15}
$f(x)=\sqrt x$, $g(x)=x^2$. Give domain and range of each; simplify $f\circ g$ and $g\circ f$; what is the difference?
\solution
$f$: domain $[0,\infty)$, range $[0,\infty)$. $g$: domain $\R$, range $[0,\infty)$.\\
$(f\circ g)(x)=\sqrt{x^2}=|x|$, defined for all $x\in\R$ (since $x^2\ge0$ is always in the domain of $f$).\\
$(g\circ f)(x)=(\sqrt x)^2=x$, but only for $x\ge0$ (the domain of $f$).\\
Difference: different domains ($\R$ versus $[0,\infty)$) and different formulas on the negative axis: $f\circ g$ is $|x|$, not $x$. On $[0,\infty)$ they agree.\qed
\end{bookex}

\begin{bookex}{17}
Are there values of $p$ for which $f:x\mapsto|x+p|+|x-4|$ is even? Odd?
\solution
\emph{Even.} For $x$ larger than both $4$ and $|p|$ all absolute values open with a plus sign: $f(x)=(x+p)+(x-4)=2x+p-4$ and $f(-x)=|{-x}+p|+|{-x}-4|=(x-p)+(x+4)=2x-p+4$. Evenness forces $p-4=-p+4$, so $p=4$. Conversely for $p=4$: $f(-x)=|{-x}+4|+|{-x}-4|=|x-4|+|x+4|=f(x)$ for all $x$. So $f$ is even exactly for $p=4$.\\
\emph{Odd.} An odd function defined at $0$ has $f(0)=-f(0)$, so $f(0)=0$. Here $f(0)=|p|+4\ge4>0$. So $f$ is odd for no $p$.\qed
\end{bookex}

\begin{bookex}{19}
$v$ even and $w$ odd, both defined on all of $\R$. Decide whether $vw$, $w^2$, $v\circ w$, $w\circ v$, $w\circ w$ are even, odd, or neither.
\solution
Replace $x$ by $-x$ and use $v(-x)=v(x)$, $w(-x)=-w(x)$:\\
$(vw)(-x)=v(x)\cdot(-w(x))=-(vw)(x)$: \textbf{odd}.\quad $w^2(-x)=(-w(x))^2=w(x)^2$: \textbf{even}.\\
$(v\circ w)(-x)=v(-w(x))=v(w(x))$: \textbf{even}.\quad $(w\circ v)(-x)=w(v(x))$: \textbf{even}.\\
$(w\circ w)(-x)=w(-w(x))=-w(w(x))$: \textbf{odd}.\qed
\end{bookex}

\begin{trybox}[{13, 14, 16, 18}]
\textbf{13} $p(x)=1+x$, $q(x)=x^2$: find $p\circ p$, $p\circ q$, $q\circ p$, $q\circ q$.\quad \textbf{14} $m(x)=-x$ satisfies $(m\circ m)(x)=x$; find more functions $f$ with $(f\circ f)(x)=x$.\quad \textbf{16} Examples of even, odd, neither, both.\quad \textbf{18} ``Every odd function has $f(0)=0$'': true? Counterexample? Something similar for even functions?
\end{trybox}

\hsection{Polynomials, divisibility and factorising}

\begin{key}[{\S0.3 \textperiodcentered\ polynomials, long division, the factor theorem}]
A \term{polynomial} (veelterm) of degree $n$ is $a_0+a_1x+\dots+a_nx^n$ with $a_n\ne0$. \term{Factorising} is the reverse of expanding: $4-16x^2=4(1-2x)(1+2x)$, $1-x^4=(1-x)(1+x)(1+x^2)$. \textbf{Long division (staartdeling)} of a polynomial by one of equal or lower degree: divide the highest-power terms, subtract, repeat; the dictaat's example gives
\[
\frac{8x^3-12x^2+6x+3}{2x^2-x}=4x-4+\frac{2x+3}{2x^2-x},\qquad\text{remainder }2x+3\text{ (lower degree than the divisor).}
\]
\textbf{Stelling 0.9 (factorstelling):} $p(x_0)=0$ iff $x-x_0$ is a factor of $p$. (Proof: divide, $p(x)=(x-x_0)q(x)+c$; substituting $x_0$ gives $c=p(x_0)$.) \textbf{Strategy:} to factorise, try $x=0,\pm1,\pm2,\dots$ to find a zero, divide out the factor, repeat; a polynomial of degree $n$ has at most $n$ zeros. \textbf{Vocabulary:} factors are multiplied, terms are added.
\end{key}

\begin{bookex}{20(a)}
Divide $\dfrac{6x^4-4x^3-x^2+6x-12}{2x^2-3}$ and check by multiplying back.
\solution
$6x^4\div2x^2=3x^2$; subtract $3x^2(2x^2-3)=6x^4-9x^2$: remainder $-4x^3+8x^2+6x-12$.\\
$-4x^3\div2x^2=-2x$; subtract $-2x(2x^2-3)=-4x^3+6x$: remainder $8x^2-12$.\\
$8x^2\div2x^2=4$; subtract $4(2x^2-3)=8x^2-12$: remainder $0$.\\
So the quotient is $3x^2-2x+4$ with remainder $0$: $2x^2-3$ is a factor. Check: $(2x^2-3)(3x^2-2x+4)=6x^4-4x^3+8x^2-9x^2+6x-12=6x^4-4x^3-x^2+6x-12$. \checkmark\qed
\end{bookex}

\begin{bookex}{21(a)}
Factorise $x^3+2x^2-5x-6$ and check by expanding.
\solution
Try small integers: $x=1$: $1+2-5-6\ne0$; $x=-1$: $-1+2+5-6=0$. So $x+1$ is a factor (Stelling 0.9). Dividing: $x^3+2x^2-5x-6=(x+1)(x^2+x-6)$, and $x^2+x-6=(x+3)(x-2)$. Result: $(x+1)(x+3)(x-2)$. Check: $(x+1)(x^2+x-6)=x^3+x^2-6x+x^2+x-6=x^3+2x^2-5x-6$. \checkmark\qed
\end{bookex}

\begin{trybox}[{20(b), 21(b), 21(c)}]
\textbf{20(b)} $\dfrac{6x^4-4x^3-x^2+6x-12}{x+2x^3}$.\quad \textbf{21(b)} $45x^4(x-14)-50x^3(x-14)+30x^2(x-14)$.\quad \textbf{21(c)} $x-x^3+x^5-x^7$.
\end{trybox}

\hsection{Trigonometry}

\begin{key}[{\S0.4 \textperiodcentered\ the unit circle}]
Angles in radians. The point of the unit circle at angle $t$ (counterclockwise from the positive $x$-axis) is $(\cos t,\sin t)$; $\tan t=\sin t/\cos t$ is the length cut off on the vertical tangent line at $(1,0)$ (similar triangles). Symmetries read off from the circle: $\sin(-t)=-\sin t$, $\cos(-t)=\cos t$, $\sin(\frac\pi2-t)=\cos t$, $\sin^2t+\cos^2t=1$ (Pythagoras). The sum rule $\sin(x+y)=\sin x\cos y+\cos x\sin y$ is given; everything else in Oefening 26 follows from it. Exact values come from the square with diagonal $1$ (angles $\frac\pi4$, sides $\frac12\sqrt2$) and the equilateral triangle with side $1$ (angles $\frac\pi3$; half of it has angles $\frac\pi6,\frac\pi3$ and sides $\frac12,\frac12\sqrt3$).
\end{key}

\begin{bookex}{23}
Fill in the table of $\sin t$, $\cos t$, $\tan t$ for $t=0,\frac\pi6,\frac\pi4,\frac\pi3,\frac\pi2$.
\solution
\begin{center}\renewcommand{\arraystretch}{1.3}
\begin{tabular}{c|ccccc}
$t$ & $0$ & $\pi/6$ & $\pi/4$ & $\pi/3$ & $\pi/2$\\\hline
$\sin t$ & $0$ & $\frac12$ & $\frac12\sqrt2$ & $\frac12\sqrt3$ & $1$\\
$\cos t$ & $1$ & $\frac12\sqrt3$ & $\frac12\sqrt2$ & $\frac12$ & $0$\\
$\tan t$ & $0$ & $\frac13\sqrt3$ & $1$ & $\sqrt3$ & undefined
\end{tabular}
\end{center}
Memory aid: the sines are $\frac12\sqrt0,\frac12\sqrt1,\frac12\sqrt2,\frac12\sqrt3,\frac12\sqrt4$; the cosines are the same row reversed.\qed
\end{bookex}

\begin{bookex}{25}
Use the table and the symmetries of the circle to find: $\sin\frac{2\pi}3$, $\cos\frac{2\pi}3$, $\tan\frac{2\pi}3$, $\sin\frac{5\pi}6$, $\cos\frac{5\pi}6$, $\tan\frac{5\pi}6$, $\sin\pi$, $\cos\pi$, $\tan\pi$, $\sin(-\frac\pi4)$, $\cos(-\frac{3\pi}4)$, $\tan\frac{3\pi}4$, $\sin\frac{4\pi}3$, $\cos\frac{5\pi}3$, $\tan\frac{6\pi}3$, $\sin(-\frac\pi6)$, $\cos(-\frac{2\pi}3)$, $\tan(-\frac{5\pi}6)$.
\solution
Reduce each angle to a ``nice'' angle in the first quadrant and fix the sign from the quadrant ($\sin$ positive above the axis, $\cos$ positive to the right, $\tan$ positive in quadrants I and III).\\
$\frac{2\pi}3=\pi-\frac\pi3$ (II): $\sin=\frac12\sqrt3$, $\cos=-\frac12$, $\tan=-\sqrt3$.\quad $\frac{5\pi}6=\pi-\frac\pi6$ (II): $\sin=\frac12$, $\cos=-\frac12\sqrt3$, $\tan=-\frac13\sqrt3$.\\
$\pi$: $\sin\pi=0$, $\cos\pi=-1$, $\tan\pi=0$.\quad $\sin(-\frac\pi4)=-\frac12\sqrt2$; $\cos(-\frac{3\pi}4)=\cos\frac{3\pi}4=-\frac12\sqrt2$; $\tan\frac{3\pi}4=-1$.\\
$\frac{4\pi}3=\pi+\frac\pi3$ (III): $\sin=-\frac12\sqrt3$.\quad $\frac{5\pi}3=2\pi-\frac\pi3$ (IV): $\cos=\frac12$.\quad $\tan\frac{6\pi}3=\tan2\pi=0$.\\
$\sin(-\frac\pi6)=-\frac12$; $\cos(-\frac{2\pi}3)=\cos\frac{2\pi}3=-\frac12$; $\tan(-\frac{5\pi}6)=-\tan\frac{5\pi}6=\frac13\sqrt3$.\qed
\end{bookex}

\begin{bookex}{26}
From $\sin(x+y)=\sin x\cos y+\cos x\sin y$ derive: (a) $\sin(x-y)$; (b) $\cos(x+y)$ using $\cos(x+y)=\sin((\frac\pi2-x)-y)$; (c) $\tan(x+y)$ in terms of $\tan$ only; (d) the double-angle formulas.
\solution
(a) Replace $y$ by $-y$ and use $\cos(-y)=\cos y$, $\sin(-y)=-\sin y$: $\sin(x-y)=\sin x\cos y-\cos x\sin y$.\\
(b) $\cos(x+y)=\sin\big((\tfrac\pi2-x)-y\big)=\sin(\tfrac\pi2-x)\cos y-\cos(\tfrac\pi2-x)\sin y=\cos x\cos y-\sin x\sin y$, using $\sin(\frac\pi2-x)=\cos x$ and $\cos(\frac\pi2-x)=\sin x$.\\
(c) $\tan(x+y)=\dfrac{\sin x\cos y+\cos x\sin y}{\cos x\cos y-\sin x\sin y}$; divide numerator and denominator by $\cos x\cos y$: $\tan(x+y)=\dfrac{\tan x+\tan y}{1-\tan x\tan y}$.\\
(d) Put $y=x$: $\sin2x=2\sin x\cos x$; $\cos2x=\cos^2x-\sin^2x=2\cos^2x-1=1-2\sin^2x$ (using $\sin^2x+\cos^2x=1$); $\tan2x=\dfrac{2\tan x}{1-\tan^2x}$.\qed
\end{bookex}

\begin{trybox}[{22, 24, 27 (Extra)}]
\textbf{22} Angles (radians) and sides of the square with diagonal $1$ and the equilateral triangle with side $1$.\quad \textbf{24} (a) Why are the labels $\sin,\cos,\tan$ placed where they are in the unit circle? (b) Signs in $\sin(\pi\pm t)$, $\cos(\pi\pm t)$, $\tan(\pi\pm t)$. (c) Check the four basic identities in the figure.\quad \textbf{27 (Extra)} A point moves along the unit circle at speed $1$; conclude $\sin'=\cos$, $\cos'=-\sin$.
\end{trybox}

\hsection{The logarithm}

\begin{key}[{Definitions 0.10, 0.11 \textperiodcentered\ logarithm, base, change of base}]
A \term{logarithm} is a function $\log:(0,\infty)\to\R$ with $\log(ab)=\log a+\log b$ for all $a,b>0$ (it turns multiplication into addition; that is why Napier and Briggs invented it). Its \term{base} (grondtal) is the number $g$ with $\log g=1$; notation ${}^g\!\log x$ (or $\log_gx$). All rules follow from the definition (Oefening 28). \textbf{Change of base:} ${}^h\!\log x=\dfrac{{}^g\!\log x}{{}^g\!\log h}$ (the right side satisfies the functional equation and takes the value $1$ at $h$). \textbf{From now on $\log$ means the natural logarithm}, base $e=2.71828\dots$ (the primitive of $1/x$, the inverse of $x\mapsto e^x$).
\end{key}

\begin{bookex}{28}
For $a,b>0$ derive from the definition, one by one: (a) $\log1=0$; (b) $\log a^2=2\log a$; (c) $\log\frac1a=-\log a$; (d) $\log a^n=n\log a$ for $n\in\Z$; (e) $\log a^{p/q}=\frac pq\log a$ for $\frac pq\in\Q$.
\solution
(a) $\log1=\log(1\cdot1)=\log1+\log1$, so $\log1=0$.\\
(b) $\log a^2=\log(a\cdot a)=\log a+\log a=2\log a$.\\
(c) $0=\log1=\log\big(a\cdot\tfrac1a\big)=\log a+\log\tfrac1a$, so $\log\tfrac1a=-\log a$.\\
(d) For $n>0$ repeat (b): $\log a^n=\log(a\cdots a)=n\log a$. For $n=0$: $\log a^0=\log1=0=0\cdot\log a$. For $n<0$: $\log a^n=\log\frac1{a^{-n}}=-\log a^{-n}=-(-n)\log a=n\log a$ by (c).\\
(e) First $p=1$: $a=(a^{1/q})^q$, so by (d) $\log a=q\log a^{1/q}$, i.e.\ $\log a^{1/q}=\frac1q\log a$. Then $\log a^{p/q}=\log(a^{1/q})^p=p\log a^{1/q}=\frac pq\log a$.\\
(The dictaat remarks that $\log a^c=c\log a$ also holds for irrational $c$, but that needs \S6.6.)\qed
\end{bookex}

\begin{bookex}{30}
Without a calculator: is $(\log3)^2-\log(3^2)$ negative, zero or positive?
\solution
$(\log3)^2-\log(3^2)=(\log3)^2-2\log3=\log3\,(\log3-2)$. Now $\log3>\log1=0$, and $\log3<2$ because $3<e^2$ (as $e>1.7$ gives $e^2>2.89$; more crudely $e^2>2^2=4>3$) and $\log$ is increasing. So the product is \textbf{negative}. (Indeed $\log3\approx1.10$, giving $\approx-0.99$.)\qed
\end{bookex}

\begin{trybox}[{29}]
\textbf{29} (a) Why do logarithms with bases $g$ and $h$ differ everywhere by a constant factor? (b) Which factor? (c) What can you say about ${}^g\!\log h\cdot{}^h\!\log g$?
\end{trybox}

\begin{warn}
\begin{itemize}
\item $\sqrt{x^2}=|x|$, never $x$; $\sqrt{1-\cos^2x}=|\sin x|$ (Oefening 12).
\item $f_{\mathrm{inv}}$ is the inverse under composition, not $1/f$. Before inverting, check bijectivity and state domain and range.
\item Do not divide an equation by an expression that can be $0$ (Oefening 5a); do not multiply an inequality by an expression of unknown sign (Oefening 4c).
\item $\log$ is the natural logarithm in this course; ``$\log$'' on a calculator is often base $10$.
\end{itemize}
\end{warn}

\chapter{Complex numbers}
\chaptercontents{A & Why complex numbers (\S1.1) \\ B & The set $\C$ (\S1.2) \\ C & Modulus, argument, conjugate (\S1.3) \\ D & Calculations in $\C$ (\S1.4) \\ E & The complex exponential (\S1.5)}%
{Exercises 101--123 (all worked or solved).\\ Hoofdpunten: use the complex plane; real and imaginary part, modulus, argument, conjugate, $e^{i\varphi}$; elementary calculations in $\C$; solving simple polynomial equations in $\C$.}

\hsection{Why complex numbers are unavoidable}
\begin{key}[{\S1.1 \textperiodcentered\ Cardano and Bombelli}]
Cardano (1545) solved $x^3+px=q$ by the formula $x=\sqrt[3]{\frac q2+\sqrt{(\frac q2)^2+(\frac p3)^3}}+\sqrt[3]{\frac q2-\sqrt{(\frac q2)^2+(\frac p3)^3}}$. It works for $x^3+6x=20$ ($x=2$), but for $x^3-15x=4$ it needs $\sqrt{-121}$, although $x=4$ is visibly a solution. Bombelli (1572) just computed on with $\sqrt{-121}$: since $(2\pm\sqrt{-1})^3=2\pm\sqrt{-121}$, the formula gives $x=(2+\sqrt{-1})+(2-\sqrt{-1})=4$. Moral: the complex numbers forced themselves upon us. (No exercises; not exam material.)
\end{key}

\hsection{The set of complex numbers}
\begin{key}[{Definition 1.1 \textperiodcentered\ $\C$, real and imaginary part}]
$i$ is a new number with $i^2=-1$ ($i=\sqrt{-1}$, and $-i$ is the other square root of $-1$). $\C=\{a+bi\mid a,b\in\R\}$; for $z=a+bi$, $\operatorname{Re}z=a$ and $\operatorname{Im}z=b$ (both real). $\R\subseteq\C$ via $x=x+0i$. \textbf{Fundamental theorem of algebra:} every non-constant polynomial has a zero in $\C$, so all zeros of all polynomials lie in $\C$ (no further extension needed). For $b\ge0$: $\sqrt{-b}=i\sqrt b$.
\end{key}

\begin{bookex}{101}
Compute the complex zeros of $P(x)=x^2-4x+13$ with the abc-formula, in the form $a+bi$.
\solution
$x=\dfrac{4\pm\sqrt{16-52}}2=\dfrac{4\pm\sqrt{-36}}2=\dfrac{4\pm6i}2=2\pm3i$. Check: $(2+3i)^2-4(2+3i)+13=4+12i-9-8-12i+13=0$. \checkmark\qed
\end{bookex}

\hsection{The complex plane: modulus, argument, conjugate}
\begin{key}[{Definitions 1.2, 1.3 \textperiodcentered\ modulus, argument, polar form, conjugate}]
Plot $a+bi$ as the point $(a,b)$: real axis horizontal, imaginary axis vertical (the Argand plane). The \term{modulus} $|z|=\sqrt{a^2+b^2}$ is the distance to $0$; the \term{argument} $\arg z$ is the angle from the positive real axis, counterclockwise, with the convention $-\pi<\arg z\le\pi$ ($\arg0$ is undefined). $\tan(\arg z)=\frac ba$ when $a\ne0$, but two angles in $(-\pi,\pi]$ have the same tangent: \textbf{always sketch $z$} to pick the right one. Positive reals have argument $0$, negative reals $\pi$, $bi$ has argument $\pm\frac\pi2$. With $r=|z|$, $\varphi=\arg z$: $z=r(\cos\varphi+i\sin\varphi)$ (\term{polar coordinates}). The \term{conjugate} $\bar z=a-bi$ is the mirror image of $z$ in the real axis.
\end{key}

\begin{bookex}{103}
Modulus and argument of $1+i$, $3i$, $-20$, $-1-i$; draw them.
\solution
$1+i$: $|1+i|=\sqrt{1+1}=\sqrt2$, in quadrant I on the diagonal: $\arg=\frac\pi4$.\quad $3i$: $|3i|=3$, on the positive imaginary axis: $\arg=\frac\pi2$.\\
$-20$: $|{-20}|=20$, negative real: $\arg=\pi$.\quad $-1-i$: $|{-1-i}|=\sqrt2$, quadrant III on the diagonal: $\arg=-\frac{3\pi}4$ (not $\frac{5\pi}4$: the argument must lie in $(-\pi,\pi]$).\qed
\end{bookex}

\begin{bookex}{104}
$z=a+bi$. Give $\arg z$ in terms of $\arctan$ for: (a) $a>0$, $b\ge0$; (b) $a>0$, $b<0$; (c) $a<0$, $b\ge0$; (d) $a<0$, $b<0$; (e) $a=0$, $b>0$; (f) $a=0$, $b<0$.
\solution
$\arctan$ takes values in $(-\frac\pi2,\frac\pi2)$, i.e.\ it returns the angle of a point in the \emph{right} half-plane.\\
(a) and (b): $z$ is in the right half-plane, so $\arg z=\arctan\frac ba$ (in $[0,\frac\pi2)$ for (a), in $(-\frac\pi2,0)$ for (b)).\\
(c) $z$ in quadrant II: $\arctan\frac ba$ is the angle of $-z$ (which is in quadrant IV), so $\arg z=\arctan\frac ba+\pi$.\\
(d) $z$ in quadrant III: $\arg z=\arctan\frac ba-\pi$ (adding $\pi$ would leave $(-\pi,\pi]$).\\
(e) $\arg z=\frac\pi2$; (f) $\arg z=-\frac\pi2$ ($\tan$ undefined).\qed
\end{bookex}

\begin{bookex}{105}
Determine $\arg z$ for $z=-\frac12\sqrt3-\frac12i$.
\solution
$|z|=\sqrt{\frac34+\frac14}=1$ and $z$ is in quadrant III (both parts negative). The reference angle has $\tan=\frac{1/2}{\sqrt3/2}=\frac1{\sqrt3}$, i.e.\ $\frac\pi6$. So $\arg z=-\pi+\frac\pi6=-\frac{5\pi}6$ (case (d) of 104: $\arctan\frac{-1/2}{-\sqrt3/2}-\pi=\frac\pi6-\pi$). Indeed $z=\cos(-\frac{5\pi}6)+i\sin(-\frac{5\pi}6)$.\qed
\end{bookex}

\begin{trybox}[{102, 106, 107, 108, 109, 110}]
\textbf{102} Draw $1+2i$, $3i$, $-1$, $-4-i$, $2-i$, $4$.\quad \textbf{106} Why is $\arg z$ undefined for $z=0$?\quad \textbf{107} Check $\operatorname{Re}\bar z=\operatorname{Re}z$, $\operatorname{Im}\bar z=-\operatorname{Im}z$, $|\bar z|=|z|$, $\arg\bar z=-\arg z$, $\bar{\bar z}=z$.\quad \textbf{108} $\bar z$ for $1+i$, $3i$, $-20$, $-1-i$.\quad \textbf{109} Solve $z^2+2z+2=0$.\quad \textbf{110} If $z$ solves $az^2+bz+c=0$ with $a,b,c\in\R$, so does $\bar z$.
\end{trybox}

\hsection{Calculations in $\C$}
\begin{key}[{Definition 1.4 and Stelling 1.5 \textperiodcentered\ sum, product, conjugation rules, division}]
$(a+bi)+(c+di)=(a+c)+(b+d)i$ and $(a+bi)(c+di)=(ac-bd)+(ad+bc)i$: this is ordinary expanding of brackets with $i^2=-1$. \textbf{Stelling 1.5:} $\overline{z+w}=\bar z+\bar w$, $\overline{zw}=\bar z\bar w$, $z\bar z=|z|^2$, $z+\bar z=2\operatorname{Re}z$. \textbf{Division:} for $w\ne0$, $\dfrac zw=\dfrac{z\bar w}{w\bar w}=\dfrac{z\bar w}{|w|^2}$: multiply numerator and denominator by the conjugate of the denominator, which makes the denominator real.
\end{key}

\begin{bookex}{111(a)}
$z=3-2i$, $w=1+i$: compute $z+w$ and $zw$.
\solution
$z+w=(3+1)+(-2+1)i=4-i$.\quad $zw=(3-2i)(1+i)=3+3i-2i-2i^2=3+i+2=5+i$. (With the formula: $ac-bd=3\cdot1-(-2)\cdot1=5$, $ad+bc=3\cdot1+(-2)\cdot1=1$.)\qed
\end{bookex}

\begin{bookex}{114}
Write $\dfrac zw$ in the form $a+bi$: (a) $z=1$, $w=2+3i$; (b) $z=1+i$, $w=i$; (c) $z=-2-i$, $w=3-4i$.
\solution
(a) $\dfrac1{2+3i}=\dfrac{2-3i}{(2+3i)(2-3i)}=\dfrac{2-3i}{4+9}=\dfrac2{13}-\dfrac3{13}i$.\\
(b) $\dfrac{1+i}i=\dfrac{(1+i)(-i)}{i(-i)}=\dfrac{-i-i^2}{1}=1-i$. (Check: $i(1-i)=i+1$. \checkmark)\\
(c) $\dfrac{-2-i}{3-4i}=\dfrac{(-2-i)(3+4i)}{9+16}=\dfrac{-6-8i-3i-4i^2}{25}=\dfrac{-2-11i}{25}=-\dfrac2{25}-\dfrac{11}{25}i$.\qed
\end{bookex}

\begin{trybox}[{111(b),(c), 112, 113, 115}]
\textbf{111} (b) $z=i$, $w=1-i$; (c) $z=w=-4+6i$.\quad \textbf{112} Check that Definition 1.4 is ``expanding brackets with $i^2=-1$''.\quad \textbf{113} Prove parts (a), (c), (d) of Stelling 1.5 by writing out.\quad \textbf{115} Commutativity, associativity and distributivity in $\C$ by writing out.
\end{trybox}

\hsection{The complex exponential function}
\begin{result}[{Stelling 1.6, Definition 1.8, Stelling 1.11 \textperiodcentered\ multiply moduli, add arguments}]
$|zw|=|z|\,|w|$ and $\arg(zw)=\arg z+\arg w$ (modulo $2\pi$); proof: multiply the polar forms and use the sum formulas. So powers are easy in polar form: Voorbeeld 1.7, $(\sqrt3+i)^3$ has modulus $2^3=8$ and argument $3\cdot\frac\pi6=\frac\pi2$, hence equals $8i$. \textbf{Definition 1.8:} $e^{it}=\cos t+i\sin t$ for real $t$ (consistent: $e^{0}=1$, $\frac d{dt}e^{it}=ie^{it}$, $e^{is}e^{it}=e^{i(s+t)}$, Voorbeeld 1.9). Hence $z=re^{i\varphi}$ with $r=|z|$, $\varphi=\arg z$, and $re^{i\varphi}\cdot se^{i\theta}=rse^{i(\varphi+\theta)}$. \textbf{De Moivre:} $(\cos\varphi+i\sin\varphi)^n=\cos n\varphi+i\sin n\varphi$. For $z=a+ib$: $e^z=e^a(\cos b+i\sin b)$.
\end{result}
\begin{strategy}[{Voorbeeld 1.10 \textperiodcentered\ solving $z^n=w$}]
Write $z=re^{i\varphi}$ and $w=se^{i\theta}$. Then $z^n=r^ne^{in\varphi}$, so $r^n=s$ gives $r=s^{1/n}$ (moduli are $\ge0$), and $n\varphi=\theta+2k\pi$ gives $\varphi=\dfrac{\theta+2k\pi}n$, $k=0,1,\dots,n-1$ ($n$ different solutions, equally spaced on a circle; other $k$ repeat them). Example: $z^6=1$ has $r=1$, $\varphi=\frac{k\pi}3$: $1,\ \frac12\pm\frac12\sqrt3i,\ -\frac12\pm\frac12\sqrt3i,\ -1$. Convert back to $a+bi$ with $\cos,\sin$ of nice angles.
\end{strategy}

\begin{bookex}{117}
Write $(1+i)^7$ in the form $a+bi$.
\solution
$1+i=\sqrt2\,e^{i\pi/4}$, so $(1+i)^7=(\sqrt2)^7e^{7i\pi/4}=8\sqrt2\,e^{-i\pi/4}$ (since $\frac{7\pi}4=2\pi-\frac\pi4$). Hence $(1+i)^7=8\sqrt2\big(\cos\tfrac\pi4-i\sin\tfrac\pi4\big)=8\sqrt2\big(\tfrac12\sqrt2-\tfrac12\sqrt2\,i\big)=8-8i$. (Check: $(1+i)^2=2i$, $(1+i)^4=-4$, $(1+i)^7=(-4)(2i)(1+i)=-8i+8$. \checkmark)\qed
\end{bookex}

\begin{bookex}{120}
Write $i$ and $1+i$ in the form $re^{i\varphi}$.
\solution
$|i|=1$, $\arg i=\frac\pi2$: $i=e^{i\pi/2}$. $|1+i|=\sqrt2$, $\arg(1+i)=\frac\pi4$: $1+i=\sqrt2\,e^{i\pi/4}$.\qed
\end{bookex}

\begin{bookex}{123}
Solve: (a) $z^4=16$; (b) $z^2=i$; (c) $z^6-2z^3+1=49$.
\solution
(a) $z=re^{i\varphi}$: $r^4=16$ gives $r=2$; $4\varphi=0+2k\pi$ gives $\varphi=\frac{k\pi}2$, $k=0,1,2,3$: $z=2,\ 2i,\ -2,\ -2i$.\\
(b) $|i|=1$, $\arg i=\frac\pi2$: $r=1$ and $2\varphi=\frac\pi2+2k\pi$, $\varphi=\frac\pi4$ or $\frac{5\pi}4\equiv-\frac{3\pi}4$: $z=\pm e^{i\pi/4}=\pm\big(\tfrac12\sqrt2+\tfrac12\sqrt2\,i\big)$. So yes, $\sqrt i$ is a complex number: $\big(\frac{1+i}{\sqrt2}\big)^2=\frac{2i}2=i$.\\
(c) Substitute $w=z^3$: $w^2-2w-48=0$, $w=\frac{2\pm\sqrt{4+192}}2=\frac{2\pm14}2$, so $w=8$ or $w=-6$.\\
$z^3=8$: $r=2$, $\varphi=\frac{2k\pi}3$: $z=2,\ 2e^{2\pi i/3}=-1+\sqrt3\,i,\ 2e^{-2\pi i/3}=-1-\sqrt3\,i$.\\
$z^3=-6=6e^{i\pi}$: $r=\sqrt[3]6$, $3\varphi=\pi+2k\pi$, $\varphi=\frac\pi3,\pi,-\frac\pi3$: $z=\sqrt[3]6\big(\tfrac12+\tfrac12\sqrt3\,i\big),\ -\sqrt[3]6,\ \sqrt[3]6\big(\tfrac12-\tfrac12\sqrt3\,i\big)$.\\
Six solutions in total, as expected for degree $6$.\qed
\end{bookex}

\begin{trybox}[{116, 118, 119, 121, 122}]
\textbf{116} Adding $a+bi$ and $c+di$ is adding the vectors $(a,b)$ and $(c,d)$.\quad \textbf{118} Geometric meaning of multiplication by $i$.\quad \textbf{119} $e^{i\pi}$ and $4e^{i\pi/3}$ in the form $a+bi$.\quad \textbf{121} $|e^{i\varphi}|=1$: all $e^{i\varphi}$ lie on the unit circle.\quad \textbf{122} Draw the six solutions of $z^6=1$; the conjugate of each solution is again a solution.
\end{trybox}

\begin{warn}
\begin{itemize}
\item $\arg z\in(-\pi,\pi]$: $-1-i$ has argument $-\frac{3\pi}4$, not $\frac{5\pi}4$. $\tan$ alone does not determine the quadrant: sketch.
\item $|z|^2=z\bar z$, not $z^2$. $\overline{zw}=\bar z\bar w$ is the conjugate of the product; do not confuse with $\bar z\,\bar w$ written out.
\item An equation $z^n=w$ has $n$ solutions; after finding the modulus, list $k=0,\dots,n-1$ and then stop ($k=n$ repeats $k=0$).
\end{itemize}
\end{warn}

\chapter{Limits, differentiation and integration}
\chaptercontents{A & Limits (\S2.1) \\ B & Continuity (\S2.2) \\ C & Differentiation (\S2.3) \\ D & Integration (\S2.4) \\ E & Integration with common sense (\S2.5)}%
{Exercises 201--220 (all worked or solved).\\ Hoofdpunten: the limit definition of the derivative and its picture; the rules for limits; the rules for differentiating, above all the chain rule; integration as ``sum of strips'', and the elementary rules that follow from that picture; some definite integrals without (much) computation.}

\hsection{Limits}

\begin{key}[{Definition 2.1 (informal), one-sided limits, Stelling 2.3 \textperiodcentered\ limits and their rules}]
$\lim_{x\to a}f(x)=b$ means: $f(x)$ can be made as close to $b$ as we like by taking $x$ close enough to $a$. Two things to keep sharp: the limit says nothing about $x=a$ itself ($f(a)$ need not equal $b$, $f$ need not even be defined at $a$); and it does not claim that $f(x)=b$ for some $x$. \textbf{Left limit} $\lim_{x\to a^-}$ ($x\uparrow a$), \textbf{right limit} $\lim_{x\to a^+}$ ($x\downarrow a$): if they differ, the limit does not exist; if they agree, that common value is the limit. $\lim f(x)=+\infty$ means $f$ takes arbitrarily large values near $a$ (strictly the limit then does not exist, $\infty$ is not a number, but we distinguish this from the case of different one-sided limits).\\
\textbf{Rules (Stelling 2.3):} if $f(x)\to\varphi$ and $g(x)\to\gamma$ as $x\to a$, then $f\pm g\to\varphi\pm\gamma$, $fg\to\varphi\gamma$, $f/g\to\varphi/\gamma$ if $\gamma\ne0$, and $\lim_{x\to a}f(g(x))=\lim_{x\to\gamma}f(x)$ (substitution).
\end{key}
\begin{result}[{Stelling 2.4 \textperiodcentered\ squeeze theorem (insluitstelling)}]
If $g(x)\le f(x)\le h(x)$ on an open interval around $a$ (except possibly at $a$) and $\lim_{x\to a}g(x)=\lim_{x\to a}h(x)=\lambda$, then $\lim_{x\to a}f(x)=\lambda$. The creative step is choosing $g$ and $h$. Voorbeeld 2.5: for $x\sin\frac1x$ use $-1\le\sin\frac1x\le1$, but since $x$ can be negative, squeeze between $g(x)=-|x|$ and $h(x)=|x|$ (Oefening 204). Standard use: ``bounded times something tending to $0$'' tends to $0$.
\end{result}

\begin{bookex}{201}
Determine the limits, if they exist: (a) $\lim_{k\to1}\dfrac{k^5-k^2}{k-1}$; (b) $\lim_{p\to\infty}e^p$; (c) $\lim_{q\to0^-}e^{1/q}$; (d) $\lim_{t\to\infty}\dfrac{\sin t}t$; (e) $\lim_{z\to0^+}\big(x-\frac1z\big)$; (f) $\lim_{x\to\infty}\log\dfrac1{1+x^2}$.
\solution
(a) Numerator and denominator both tend to $0$; factor: $k^5-k^2=k^2(k^3-1)=k^2(k-1)(k^2+k+1)$, so for $k\ne1$ the quotient equals $k^2(k^2+k+1)$, a polynomial, with limit $1\cdot3=3$ (Oefening 203b).\\
(b) $e^p$ grows without bound: $+\infty$ (no real limit).\\
(c) As $q\to0^-$, $1/q\to-\infty$ and $e^{1/q}\to0$. (From the right, $1/q\to+\infty$ and $e^{1/q}\to+\infty$: the two-sided limit does not exist.)\\
(d) $-\frac1t\le\frac{\sin t}t\le\frac1t$ for $t>0$, and $\pm\frac1t\to0$: by the squeeze theorem the limit is $0$.\\
(e) $\frac1z\to+\infty$ as $z\downarrow0$, so $x-\frac1z\to-\infty$ for any fixed $x$: no real limit.\\
(f) $\frac1{1+x^2}\to0^+$ and $\log y\to-\infty$ as $y\downarrow0$, so the limit is $-\infty$.\qed
\end{bookex}

\begin{bookex}{205}
Determine: (a) $\lim_{x\to0}\sin(2x)\cos(x^{-2})$; (b) $\lim_{x\to\infty}\sin(5x)e^{-5x}$; (c) $\lim_{x\to\infty}\dfrac{4x-3\cos(2\sqrt x+1)}{5-6x}$.
\solution
(a) $|\sin(2x)\cos(x^{-2})|\le|\sin2x|$ because $|\cos|\le1$, i.e.\ $-|\sin2x|\le f(x)\le|\sin2x|$, and $|\sin2x|\to0$. Squeeze: the limit is $0$. (The factor $\cos(x^{-2})$ oscillates wildly, but it is bounded.)\\
(b) $-e^{-5x}\le\sin(5x)e^{-5x}\le e^{-5x}$ and $e^{-5x}\to0$: limit $0$.\\
(c) Divide numerator and denominator by $x$: $\dfrac{4-3\cos(2\sqrt x+1)/x}{5/x-6}$. Here $|3\cos(\cdot)/x|\le3/x\to0$ (squeeze) and $5/x\to0$, so by the quotient rule the limit is $\dfrac4{-6}=-\dfrac23$.\qed
\end{bookex}

\begin{trybox}[{202, 203, 204}]
\textbf{202} (a) $\lim_{x\to0}\frac{e^x}{1+x}$ and $\lim_{x\to1}\frac{e^{x-1}}x$. (b) Why is $\lim_{x\to a-b}f(x)=\lim_{x\to a}f(x-b)$, and how does that apply to (a)?\quad \textbf{203} From the rules: $\lim kf(x)=k\varphi$; $\lim_{x\to a}P(x)=P(a)$ for polynomials; $\lim f(x)^k=\varphi^k$ for $k\in\Q$, $k>0$.\quad \textbf{204} Finish Voorbeeld 2.5: $\lim_{x\to0}|x|=0$, $\lim_{x\to0}(-|x|)=0$, squeeze $x\sin\frac1x$, sketch.
\end{trybox}

\hsection{Continuity}
\begin{key}[{Definition 2.6 and Voorbeeld 2.7}]
$f$ is \term{continuous at $a$} if $a$ is in the domain and $\lim_{x\to a}f(x)=f(a)$; continuous if this holds at every point of the domain (no jumps). Watch fractions: $f(x)=\frac{x^2-9}{x-3}$ is not defined at $3$, but for $x\ne3$ it equals $x+3$, so $\lim_{x\to3}f(x)=6$ (cancel the factor $x-3$ only for $x\ne3$). Continuity plays a small role in this course; differentiability is what matters.
\end{key}

\hsection{Differentiation}
\begin{key}[{Definition 2.8 and Oefening 206 \textperiodcentered\ the derivative as a limit}]
\[
f'(a)=\lim_{h\to0}\frac{f(a+h)-f(a)}h=\lim_{x\to a}\frac{f(x)-f(a)}{x-a},
\]
provided the limit exists as a real number: the slope of the chord through $(a,f(a))$ and $(x,f(x))$ ``just before'' $x-a$ becomes $0$ (figure 2.3: the small right triangle with legs $x-a$ and $f(x)-f(a)$). Remember the formula together with the picture. Notations: $f'=\frac{df}{dx}=\frac d{dx}f$, and $\dot x$ for $\frac{dx}{dt}$ in physics.
\end{key}
\begin{result}[{Stelling 2.9 \textperiodcentered\ rules}]
$\frac d{dx}x^r=rx^{r-1}$, $\sin'=\cos$, $\cos'=-\sin$, $\log'x=\frac1x$, $(e^x)'=e^x$; sum rule; product rule $(fg)'=f'g+fg'$; quotient rule $\big(\frac fg\big)'=\frac{f'g-fg'}{g^2}$; \textbf{chain rule} $f(g(x))'=f'(g(x))\,g'(x)$, i.e.\ $(f\circ g)'=(f'\circ g)\cdot g'$ (Oefening 208). The quotient rule follows from the product and chain rules applied to $f\cdot g^{-1}$ (Oefening 207).
\end{result}

\begin{bookex}{209}
Differentiate with the rules: (a) $f(x)=e^{x^3}(1-x^2)$; (b) $g(x)=\sin(x^2+\sqrt{3x})$; (c) $h(x)=\dfrac{\log(4x+1)}{4x}$; (d) $j(x)=2^x$ (rewrite as an $e$-power); (e) $k(x)=\log(\log(\pi x))$; (f) $m(x)=\sqrt[3]{x^3+1}\,e^{\sin x}$.
\solution
(a) Product and chain rule: $f'(x)=3x^2e^{x^3}(1-x^2)+e^{x^3}(-2x)=e^{x^3}\big(3x^2-3x^4-2x\big)$. (If the exponent is meant to be $x^3(1-x^2)=x^3-x^5$, then $f'(x)=(3x^2-5x^4)e^{x^3-x^5}$.)\\
(b) Chain rule, inner derivative $2x+\dfrac3{2\sqrt{3x}}$: $g'(x)=\cos\big(x^2+\sqrt{3x}\big)\Big(2x+\dfrac3{2\sqrt{3x}}\Big)$. (If the inner function is $x^2+\sqrt3\,x$, the inner derivative is $2x+\sqrt3$.)\\
(c) Quotient rule with $(\log(4x+1))'=\frac4{4x+1}$: $h'(x)=\dfrac{\frac4{4x+1}\cdot4x-4\log(4x+1)}{16x^2}=\dfrac{\frac{4x}{4x+1}-\log(4x+1)}{4x^2}$.\\
(d) $2^x=e^{x\log2}$, so $j'(x)=\log2\cdot e^{x\log2}=2^x\log2$.\\
(e) Chain rule twice: $k'(x)=\dfrac1{\log(\pi x)}\cdot\dfrac1{\pi x}\cdot\pi=\dfrac1{x\log(\pi x)}$.\\
(f) $\big((x^3+1)^{1/3}\big)'=\frac13(x^3+1)^{-2/3}\cdot3x^2=x^2(x^3+1)^{-2/3}$ and $(e^{\sin x})'=\cos x\,e^{\sin x}$, so $m'(x)=e^{\sin x}\Big(\dfrac{x^2}{(x^3+1)^{2/3}}+\cos x\sqrt[3]{x^3+1}\Big)$.\qed
\end{bookex}

\begin{bookex}{210}
(a) Show $1+\tan^2x=\dfrac1{\cos^2x}$. (b) Differentiate $\tan x$. (c) Conclude that $\tan'$ has two very different-looking forms.
\solution
(a) $1+\tan^2x=\dfrac{\cos^2x+\sin^2x}{\cos^2x}=\dfrac1{\cos^2x}$.\quad (b) Quotient rule: $\tan'x=\dfrac{\cos x\cos x-\sin x(-\sin x)}{\cos^2x}=\dfrac1{\cos^2x}$.\quad (c) $\tan'x=\dfrac1{\cos^2x}=1+\tan^2x$; the second form is what makes Oefening 317 work.\qed
\end{bookex}

\begin{strategy}[{\S2.3.1 \textperiodcentered\ recognise a limit as a derivative}]
$\lim_{x\to0}\dfrac{\sin x}x=\lim_{x\to0}\dfrac{\sin x-\sin0}{x-0}=\sin'(0)=\cos0=1$: a \textbf{standard limit}. Whenever a limit has the shape $\frac{f(x)-f(a)}{x-a}$ with $x\to a$, it equals $f'(a)$.
\end{strategy}

\begin{bookex}{211}
Compute in the same way: (a) $\lim_{x\to0}\dfrac{\cos x-1}x$; (b) $\lim_{x\to1}\dfrac{e^x-e}{x-1}$.
\solution
(a) $\cos0=1$, so this is $\lim_{x\to0}\dfrac{\cos x-\cos0}{x-0}=\cos'(0)=-\sin0=0$.\quad (b) $e=e^1$, so this is $\dfrac d{dx}e^x\big|_{x=1}=e^1=e$.\qed
\end{bookex}

\begin{bookex}{212}
Which is true: (A) every continuous function is differentiable; (B) every differentiable function is continuous? Why?
\solution
\emph{(B) is true.} If $f'(a)$ exists then $f(x)-f(a)=\dfrac{f(x)-f(a)}{x-a}\cdot(x-a)\to f'(a)\cdot0=0$ as $x\to a$ (product rule for limits), i.e.\ $\lim_{x\to a}f(x)=f(a)$.\\
\emph{(A) is false.} $f(x)=|x|$ is continuous at $0$, but $\dfrac{f(x)-f(0)}{x-0}=\dfrac{|x|}x$ is $-1$ for $x<0$ and $+1$ for $x>0$: the one-sided limits differ, so $f'(0)$ does not exist (the graph has a corner).\qed
\end{bookex}

\begin{trybox}[{206, 207, 208}]
\textbf{206} Why can Definition 2.8 be written with $x\to a$ instead of $h\to0$?\quad \textbf{207} Derive the quotient rule from $f\cdot g^{-1}$.\quad \textbf{208} Complete $(\dots)'=\dots\cdot g'$ in composition notation.
\end{trybox}

\hsection{Integration}
\begin{key}[{Definitions 2.10 and Stelling 2.11 \textperiodcentered\ primitives, the fundamental theorem}]
$g$ is a \term{primitive} (antiderivative) of $f$ if $g'=f$. Primitives are unique up to a constant: the \term{indefinite integral} $\int f(x)\,dx=g(x)+c$ is the whole set of primitives ($c$: integration constant; easily forgotten, and the dictaat lets it pop up unexpectedly). The \term{definite integral} is a number: $\int_a^bf(x)\,dx=g(b)-g(a)$ (the constant cancels, Oefening 213). \textbf{Picture (figure 2.4):} $\int_a^bf(x)\,dx$ is the sum ($\int$ is a long s) of strips of height $f(x)$ and infinitely small width $dx$ between $a$ and $b$; heights can be negative. \textbf{Hoofdstelling 2.11:} for continuous $f$ on $[a,b]$, $F(x)=\int_a^xf(t)\,dt$ is a primitive of $f$, and $\int_a^bf(x)\,dx=F(b)-F(a)$ for any primitive $F$. The integration variable is a dummy: $\int f(x)dx=\int f(t)dt$; never reuse it outside the integral ($x^2+\int_0^xe^t\,dt$ is right, $x^2+\int_0^xe^x\,dx$ is wrong).
\end{key}

\begin{trybox}[{213, 214 (Extra)}]
\textbf{213} If $\tilde g=g+c$ then $\tilde g(b)-\tilde g(a)=g(b)-g(a)$: no integration constant in definite integrals.\quad \textbf{214 (Extra)} With $g(x)=\int_0^xf(t)\,dt$: write $g(x+dx)$ and $dg=g(x+dx)-g(x)$ as integrals and argue $\frac{dg}{dx}=f(x)$.
\end{trybox}

\hsection{Integration with common sense}
\begin{result}[{Stelling 2.12 \textperiodcentered\ rules for definite integrals}]
(a) for $f\ge0$ continuous, $\int_a^bf$ is the area under the graph; (b) $\int_a^af=0$; (c) $\int_a^bf=-\int_b^af$; (d) $\int_a^bf=\int_a^cf+\int_c^bf$ (cut up); (e) $\int_a^bf(x)\,dx=\int_{a+c}^{b+c}f(x-c)\,dx$ (translate); (f) linearity; (g) $f$ odd $\Rightarrow\int_{-a}^af=0$; (h) $f$ even $\Rightarrow\int_{-a}^af=2\int_0^af$; (i) changing $f$ at a single point does not change the integral; (j) $f\le g$ on $(a,b)\Rightarrow\int_a^bf\le\int_a^bg$; (k) $\big|\int_a^bf\big|\le\int_a^b|f|$. Do not memorise; recognise when they help. Voorbeeld 2.13: $\int_{-1}^1(x-4)\,dx=0-4\cdot2=-8$ (odd part vanishes, rectangle $2\times1$ four times); or translate to $\int_{-5}^{-3}x\,dx=-8$ (signed area of a trapezium); $\int_{-2}^2\sqrt{4-x^2}\,dx=2\pi$ (half disc of radius $2$).
\end{result}

\begin{bookex}{215}
Compare, with a sketch: (a) the area enclosed by $y=x-1$, $x=0$, $y=0$; (b) the integral $\int_0^1(x-1)\,dx$.
\solution
(a) The region is the triangle with vertices $(0,0)$, $(1,0)$, $(0,-1)$: area $\frac12\cdot1\cdot1=\frac12$.\\
(b) $\int_0^1(x-1)\,dx=\big[\tfrac12x^2-x\big]_0^1=-\tfrac12$. The integral is minus the area, because the strips have negative height: the triangle lies below the $x$-axis. Area is non-negative; a definite integral is signed.\qed
\end{bookex}

\begin{bookex}{218}
Compute with the basic rules but without primitives: (a) $\int_0^{2\pi}(\pi x+e^\pi\sin x)\,dx$; (b) $\int_{-2}^0(x+1)\cos(x+1)\,dx$; (c) $\int_0^6\sqrt{9-(x-3)^2}\,dx$; (d) $\int_{-1}^3\big(2(x-1)^3-4\sin(x-1)+\operatorname{sgn}x+5|x|\big)dx$; (e) $\int_0^{\pi/2}\sin^2t\,dt$.
\solution
(a) Linearity: $\pi\int_0^{2\pi}x\,dx+e^\pi\int_0^{2\pi}\sin x\,dx$. The first integral is the area of the triangle with base $2\pi$ and height $2\pi$: $2\pi^2$. The second is $0$: translating by $\pi$, $\int_0^{2\pi}\sin x\,dx=\int_{-\pi}^{\pi}\sin(x+\pi)\,dx=-\int_{-\pi}^\pi\sin x\,dx=0$ by oddness (or: the positive hump on $[0,\pi]$ cancels the negative one on $[\pi,2\pi]$). Total: $2\pi^3$.\\
(b) Translate by $1$: $\int_{-2}^0(x+1)\cos(x+1)\,dx=\int_{-1}^1x\cos x\,dx=0$, since $x\cos x$ is odd (odd times even).\\
(c) $y=\sqrt{9-(x-3)^2}$ is the upper half of the circle with centre $(3,0)$ and radius $3$: the integral is the area of a half disc, $\frac12\pi\cdot3^2=\frac92\pi$.\\
(d) Split by linearity. Translating by $1$: $\int_{-1}^3\big(2(x-1)^3-4\sin(x-1)\big)dx=\int_{-2}^2(2x^3-4\sin x)\,dx=0$ (odd integrand). $\int_{-1}^3\operatorname{sgn}x\,dx=\int_{-1}^0(-1)\,dx+\int_0^31\,dx=-1+3=2$ (the value at $0$ is irrelevant, rule (i)). $\int_{-1}^35|x|\,dx=5\big(\tfrac12\cdot1\cdot1+\tfrac12\cdot3\cdot3\big)=25$ (two triangles). Total: $27$.\\
(e) The substitution $t\mapsto\frac\pi2-t$ reflects $[0,\frac\pi2]$ onto itself and turns $\sin$ into $\cos$, so $\int_0^{\pi/2}\sin^2t\,dt=\int_0^{\pi/2}\cos^2t\,dt$. Their sum is $\int_0^{\pi/2}(\sin^2t+\cos^2t)\,dt=\int_0^{\pi/2}1\,dt=\frac\pi2$, so each equals $\frac\pi4$.\qed
\end{bookex}

\begin{bookex}{219}
Show with Stelling 2.12, without a primitive, that $\int_a^bx\,dx=\frac12(b^2-a^2)$.
\solution
First $\int_0^cx\,dx=\frac12c^2$ for every $c$: for $c\ge0$ it is the area of the triangle with legs $c$ and $c$; for $c<0$, by the reversal rule $\int_0^cx\,dx=-\int_c^0x\,dx$, and $\int_c^0x\,dx$ is minus the area $\frac12c^2$ of the triangle below the axis, so again $\frac12c^2$. Now cut up: $\int_a^bx\,dx=\int_a^0x\,dx+\int_0^bx\,dx=-\int_0^ax\,dx+\int_0^bx\,dx=\frac12b^2-\frac12a^2$.\qed
\end{bookex}

\begin{trybox}[{216, 217, 220}]
\textbf{216} Sketch the rules of Stelling 2.12 that surprise you.\quad \textbf{217} Match the graphs of $f$ with those of $\int_a^xf(t)\,dt$.\quad \textbf{220} The Gini index $G=2\int_0^1(x-v(x))\,dx$: sketch $v$ for a socialist and a capitalist society, interpret $G$ graphically, find its extreme values.
\end{trybox}

\begin{warn}
\begin{itemize}
\item A limit at $a$ ignores $f(a)$; one-sided limits that differ mean ``does not exist''; $\pm\infty$ is a description of behaviour, not a number.
\item The chain rule: differentiate the outer function at the inner one, then multiply by the inner derivative. ``Kettingregel vergeten?'' is the dictaat's first question when you are stuck.
\item A definite integral is a signed quantity; area is not. Odd/even symmetry and translation remove most of the work in \S2.5 problems.
\item Never write the integration variable outside its integral.
\end{itemize}
\end{warn}

\chapter{Approximating functions}
\chaptercontents{A & Linearisation (\S3.1) \\ B & Second-order approximation (\S3.2) \\ C & Taylor approximation of order $n$ (\S3.3) \\ D & l'H\^opital, standard limits (\S3.4) \\ E & (Extra) Infinite sums (\S3.5)}%
{Exercises 301--330 (all worked or solved).\\ Hoofdpunten: notation $f^{(n)}$ and $n!$; the Taylor formula (with steunpunt and order); using Taylor approximations for values and for limits; l'H\^opital; the standard limits $\lim_{x\to0}\frac{\sin x}x=1$, $\lim_{x\to\infty}\frac{\log x}x=0$, $\lim_{t\to\infty}\frac{t^x}{e^t}=0$.}

\hsection{Linearisation}
\begin{key}[{\S3.1 \textperiodcentered\ the tangent line as approximation}]
The \term{linearisation} of $f$ at the \term{steunpunt} (base point) $a$ is the tangent line, written in powers of $x-a$:
\[
P_1(x)=f(a)+f'(a)(x-a),\qquad P_1(a)=f(a),\ P_1'(a)=f'(a).
\]
Voorbeeld 3.1: $\sqrt5$ with $f(x)=\sqrt x$, $a=4$: $f(4)=2$, $f'(4)=\frac1{2\sqrt4}=\frac14$, so $\sqrt5\approx2+\frac14(5-4)=2\frac14$; self-check: $(2\frac14)^2=5\frac1{16}$, close to $5$.
\end{key}

\begin{bookex}{301}
Linear approximations of: (a) $\sqrt e$ with $f(x)=e^x$, $a=0$; (b) $\log3$ with $f(x)=\log(1+x)$, $a=0$; (c) $\sqrt[3]4$ with $f(x)=\sqrt[3]x$, $a=1$.
\solution
(a) $f(0)=1$, $f'(0)=1$: $P_1(x)=1+x$, and $\sqrt e=e^{1/2}\approx1+\frac12=1.5$.\\
(b) $f(0)=0$, $f'(x)=\frac1{1+x}$, $f'(0)=1$: $P_1(x)=x$, and $\log3=f(2)\approx2$.\\
(c) $f(1)=1$, $f'(x)=\frac13x^{-2/3}$, $f'(1)=\frac13$: $P_1(x)=1+\frac13(x-1)$, and $\sqrt[3]4\approx1+\frac13\cdot3=2$.\\
(True values: $1.6487$, $1.0986$, $1.5874$: (b) and (c) are far from the steunpunt, so the approximations are poor.)\qed
\end{bookex}

\begin{bookex}{302}
The approximation of $\sqrt5$ was slightly too large. Predict, without a calculator, whether the answers in 301 are too large or too small. Hint: the second derivative.
\solution
The tangent line lies below the graph where $f''>0$ (convex) and above it where $f''<0$ (concave), so $P_1$ is too small for convex $f$ and too large for concave $f$ (near $a$; the second-order term $\frac12f''(a)(x-a)^2$ has the sign of $f''$).\\
(a) $(e^x)''=e^x>0$: $1.5$ is too small. (b) $\log(1+x)''=-\frac1{(1+x)^2}<0$: $2$ is too large. (c) $(x^{1/3})''=-\frac29x^{-5/3}<0$ for $x>0$: $2$ is too large. ($\sqrt x$ is concave, which explains Voorbeeld 3.1.)\qed
\end{bookex}

\hsection{Second-order approximation}
\begin{key}[{\S3.2 \textperiodcentered\ matching the curvature}]
$P_2(x)=f(a)+f'(a)(x-a)+c_2(x-a)^2$ with $c_2=\frac12f''(a)$ satisfies $P_2(a)=f(a)$, $P_2'(a)=f'(a)$, $P_2''(a)=f''(a)$ (Oefening 303). More work does not always give a better approximation, especially far from the steunpunt; the cure is even more work (next section).
\end{key}

\begin{bookex}{303}
(a) Show that $P_2(x)=f(a)+f'(a)(x-a)+c_2(x-a)^2$ satisfies $P_2(a)=f(a)$, $P_2'(a)=f'(a)$ for every $c_2$. (b) Show that $c_2=\frac12f''(a)$ gives $P_2''(a)=f''(a)$. (c) Redo Voorbeeld 3.1 with a second-order approximation; check that it is sometimes better and sometimes not.
\solution
(a) $P_2(a)=f(a)+0+0$. $P_2'(x)=f'(a)+2c_2(x-a)$, so $P_2'(a)=f'(a)$.\\
(b) $P_2''(x)=2c_2$, so $P_2''(a)=f''(a)$ iff $c_2=\frac12f''(a)$.\\
(c) $f(x)=\sqrt x$, $a=4$: $f''(x)=-\frac14x^{-3/2}$, $f''(4)=-\frac1{32}$, $c_2=-\frac1{64}$: $P_2(x)=2+\frac14(x-4)-\frac1{64}(x-4)^2$. At $x=5$: $P_2(5)=2.25-0.015625=2.234375$ versus $\sqrt5=2.23607$: error $0.0017$ instead of $0.014$ for $P_1$. Better. At $x=20$: $P_1(20)=6$, $P_2(20)=6-\frac{256}{64}=2$, while $\sqrt{20}=4.47$: $P_2$ is now worse than $P_1$ (far from the steunpunt the quadratic term dominates in the wrong direction).\qed
\end{bookex}

\begin{trybox}[{304}]
\textbf{304} Look at 302 again, now that you know $P_2$: the correction term $\frac12f''(a)(x-a)^2$ has the sign of $f''(a)$.
\end{trybox}

\hsection{Taylor approximation of order $n$}
\begin{key}[{Definitions 3.2--3.4 \textperiodcentered\ $f^{(k)}$, $k!$, the Taylor polynomial}]
$f^{(k)}$ is the $k$-th derivative, $f^{(0)}=f$; $k!=k(k-1)\cdots2\cdot1$, $0!=1$. The \term{Taylor approximation of order $n$ with steunpunt $a$} is
\[
P_n(x)=f(a)+f'(a)(x-a)+\frac{f''(a)}{2!}(x-a)^2+\dots+\frac{f^{(n)}(a)}{n!}(x-a)^n=\sum_{k=0}^n\frac{f^{(k)}(a)}{k!}(x-a)^k,
\]
the unique polynomial of degree $\le n$ with $P_n^{(k)}(a)=f^{(k)}(a)$ for $k=0,\dots,n$ (Oefening 306). Raising the order adds terms without redoing earlier work. \textbf{The standard ones at $a=0$} (Oefening 307): $e^x=1+x+\frac{x^2}2+\frac{x^3}{3!}+\dots$; $\sin x=x-\frac{x^3}{3!}+\frac{x^5}{5!}-\dots$; $\cos x=1-\frac{x^2}2+\frac{x^4}{4!}-\dots$; $\log(1+x)=x-\frac{x^2}2+\frac{x^3}3-\frac{x^4}4+\dots$; also $\frac1{1-x}=1+x+x^2+\dots$ and $\sqrt{1+x}=1+\frac x2-\frac{x^2}8+\dots$ (Oefening 308).
\end{key}
\begin{strategy}[{Voorbeelden 3.5, 3.6 and Oefening 316 \textperiodcentered\ three shortcuts}]
\textbf{Substitution:} if $u\to0$ when $x\to0$, substitute $u$ into a known expansion: $\log(1+x^2)=x^2-\frac12x^4+\dots$; $\log(\cos x)=\log(1+u)$ with $u=\cos x-1=-\frac12x^2+\frac1{24}x^4+\text{h.o.t.}$, so $\log\cos x=u-\frac12u^2+\dots=-\frac12x^2+\frac1{24}x^4-\frac12\cdot\frac14x^4+\dots=-\frac12x^2-\frac1{12}x^4+\text{h.o.t.}$ Keep track of which terms are higher order than you need. \textbf{Algebra:} $\log\frac{1+x}{1-x}=\log(1+x)-\log(1-x)$ (Oefening 314). \textbf{A differential equation:} if $f'=(2x-2)f$ (Oefening 316) or $\tan'=1+\tan^2$ (317), differentiate the relation repeatedly and substitute $x=0$ instead of computing messy quotients.
\end{strategy}

\begin{bookex}{308}
Fourth-order Taylor approximations at $0$ of: (a) $\dfrac1{x-1}$; (b) $\sqrt{1+x}$; (c) $\sqrt{1+x^2}$.
\solution
(a) $f(x)=(x-1)^{-1}$: $f'=-(x-1)^{-2}$, $f''=2(x-1)^{-3}$, $f'''=-6(x-1)^{-4}$, $f''''=24(x-1)^{-5}$; at $0$: $-1,-1,-2,-6,-24$, so $c_k=f^{(k)}(0)/k!=-1$: $P_4(x)=-1-x-x^2-x^3-x^4$. (Check: $\frac1{x-1}=-\frac1{1-x}=-(1+x+x^2+\dots)$.)\\
(b) $f=(1+x)^{1/2}$: $f'=\frac12(1+x)^{-1/2}$, $f''=-\frac14(1+x)^{-3/2}$, $f'''=\frac38(1+x)^{-5/2}$, $f''''=-\frac{15}{16}(1+x)^{-7/2}$; at $0$: $1,\frac12,-\frac14,\frac38,-\frac{15}{16}$. Divide by $k!$: $P_4(x)=1+\frac12x-\frac18x^2+\frac1{16}x^3-\frac5{128}x^4$.\\
(c) Substitute $u=x^2$ in (b): $\sqrt{1+x^2}=1+\frac12x^2-\frac18x^4+\text{h.o.t.}$, so $P_4(x)=1+\frac12x^2-\frac18x^4$ (the $x^3$ term of (b) would give $x^6$).\qed
\end{bookex}

\begin{bookex}{309}
Compute $\sin37^\circ$ to three decimals with a Taylor approximation, and check.
\solution
Radians first: $37^\circ=\frac{37\pi}{180}\approx0.64577$. With $P_5(x)=x-\frac{x^3}6+\frac{x^5}{120}$: $x^3\approx0.26930$, $x^5\approx0.11230$, so $\sin37^\circ\approx0.64577-0.04488+0.00094=0.60182$. Three decimals: $0.602$. Check: $\sin37^\circ\approx0.6018$ (and $\sin30^\circ=0.5<0.602<0.707=\sin45^\circ$, plausible). Using $37$ instead of $0.6458$ in the polynomial gives nonsense: degrees must be converted.\qed
\end{bookex}

\begin{bookex}{314}
Find a fifth-degree Taylor approximation of $\log\dfrac{1+x}{1-x}$ at $0$, the smart way.
\solution
$\log\frac{1+x}{1-x}=\log(1+x)-\log(1-x)$. From Oefening 307/313: $\log(1+x)=x-\frac{x^2}2+\frac{x^3}3-\frac{x^4}4+\frac{x^5}5$ and $\log(1-x)=-x-\frac{x^2}2-\frac{x^3}3-\frac{x^4}4-\frac{x^5}5$ (substitute $-x$). Subtracting, the even powers cancel and the odd ones double:
\[
\log\frac{1+x}{1-x}\approx2x+\frac23x^3+\frac25x^5 .
\]
(The function is odd, so only odd powers can occur; Oefening 311.)\qed
\end{bookex}

\begin{bookex}{316}
Taylor approximation of $f(x)=e^{x^2-2x}$ at $0$, reusing results: (a) $f(0)$; (b) show $f'(x)=(2x-2)f(x)$, find $f'(0)$; (c) differentiate again, find $f''(0)$; (d) as many $f^{(n)}(0)$ as you can.
\solution
(a) $f(0)=e^0=1$.\quad (b) Chain rule: $f'(x)=(2x-2)e^{x^2-2x}=(2x-2)f(x)$, so $f'(0)=-2\cdot1=-2$.\\
(c) Product rule on $f'=(2x-2)f$: $f''=2f+(2x-2)f'$, so $f''(0)=2\cdot1+(-2)(-2)=6$.\\
(d) Keep differentiating the relation: $f'''=2f'+2f'+(2x-2)f''=4f'+(2x-2)f''$, so $f'''(0)=4(-2)-2\cdot6=-20$; $f''''=4f''+2f''+(2x-2)f'''=6f''+(2x-2)f'''$, so $f''''(0)=36+40=76$; in general $f^{(n+1)}=2nf^{(n-1)}+(2x-2)f^{(n)}$, giving $f^{(n+1)}(0)=2nf^{(n-1)}(0)-2f^{(n)}(0)$: $f^{(5)}(0)=8\cdot6-2\cdot76=-104$, $f^{(6)}(0)=10(-20)-2(-104)=8$, \dots\ Hence $P_4(x)=1-2x+3x^2-\frac{10}3x^3+\frac{19}6x^4$.\qed
\end{bookex}

\begin{bookex}{317}
Redo Oefening 312 ($P_5$ of $\tan x$ at $0$) efficiently, using $\tan'x=1+\tan^2x$.
\solution
Write $T=\tan x$, so $T'=1+T^2$ and $T(0)=0$. Differentiate the relation repeatedly, always substituting $T'=1+T^2$:
$T''=2TT'=2T+2T^3$; $T'''=(2+6T^2)T'=(2+6T^2)(1+T^2)=2+8T^2+6T^4$; $T''''=(16T+24T^3)(1+T^2)=16T+40T^3+24T^5$; $T^{(5)}=(16+120T^2+120T^4)(1+T^2)$.\\
At $x=0$ ($T=0$): $T=0$, $T'=1$, $T''=0$, $T'''=2$, $T''''=0$, $T^{(5)}=16$. So
\[
\tan x\approx x+\frac2{3!}x^3+\frac{16}{5!}x^5=x+\frac13x^3+\frac2{15}x^5 .
\]
(Odd function, odd powers only. Doing this with the quotient rule on $\sin/\cos$ is the long way of Oefening 312.)\qed
\end{bookex}

\begin{bookex}{319}
With the substitution technique, order-4 approximations at $0$ of: (a) $\log(2+x)$ (hint: $2+x=2(1+\frac x2)$); (b) $\log\big((1+x)^2\big)$, in a very short and in a longer way.
\solution
(a) $\log(2+x)=\log2+\log(1+\tfrac x2)$; substitute $u=\frac x2$ in $\log(1+u)=u-\frac{u^2}2+\frac{u^3}3-\frac{u^4}4$: $\log(2+x)\approx\log2+\frac x2-\frac{x^2}8+\frac{x^3}{24}-\frac{x^4}{64}$.\\
(b) Short: $\log(1+x)^2=2\log(1+x)\approx2x-x^2+\frac23x^3-\frac12x^4$. Longer: $(1+x)^2=1+u$ with $u=2x+x^2$; then $u^2=4x^2+4x^3+x^4$, $u^3=8x^3+12x^4+\dots$, $u^4=16x^4+\dots$, and $u-\frac{u^2}2+\frac{u^3}3-\frac{u^4}4=2x+x^2-2x^2-2x^3-\frac12x^4+\frac83x^3+4x^4-4x^4+\text{h.o.t.}=2x-x^2+\frac23x^3-\frac12x^4$. Same answer. \checkmark\qed
\end{bookex}

\begin{bookex}{320}
Order-4 approximations at $0$ of (a) $e^{\sin x}$ and (b) $e^{\cos2x}$; (c) check.
\solution
(a) $u=\sin x=x-\frac16x^3+\text{h.o.t.}$ (the $x^5$ term is beyond order 4). $u^2=x^2-\frac13x^4+\dots$, $u^3=x^3+\dots$, $u^4=x^4+\dots$. Then $e^u=1+u+\frac{u^2}2+\frac{u^3}6+\frac{u^4}{24}=1+x-\frac16x^3+\frac12x^2-\frac16x^4+\frac16x^3+\frac1{24}x^4+\dots=1+x+\frac12x^2-\frac18x^4$. (The $x^3$ terms cancel.)\\
(b) $\cos2x=1-2x^2+\frac23x^4+\dots$ (substitute $2x$ in the cosine series: $\frac{(2x)^2}2=2x^2$, $\frac{(2x)^4}{24}=\frac23x^4$). Since $\cos2x$ does not tend to $0$, first split off the constant: $e^{\cos2x}=e\cdot e^{u}$ with $u=-2x^2+\frac23x^4+\dots$; $u^2=4x^4+\dots$; so $e^u=1+u+\frac{u^2}2+\dots=1-2x^2+\frac23x^4+2x^4=1-2x^2+\frac83x^4$, and $e^{\cos2x}\approx e\big(1-2x^2+\frac83x^4\big)$.\\
(c) Direct differentiation (or Wolfram Alpha) confirms both; the direct route for (a) needs $f'=\cos x\,e^{\sin x}$, $f''=(\cos^2x-\sin x)e^{\sin x}$, \dots, much messier.\qed
\end{bookex}

\begin{trybox}[{305, 306, 307, 310, 311, 312, 313, 315, 318, 321, 322 (Extra)}]
\textbf{305} $P_3$ matches $f,f',f'',f'''$ at $a$.\quad \textbf{306} $P_n^{(k)}(a)=f^{(k)}(a)$.\quad \textbf{307} Complete the degree-9 expansions of $e^x$, $\sin x$, $\cos x$, $\log(1+x)$.\quad \textbf{310} All Taylor approximations of $2x^3-3x^2+4$ at $0$; is a polynomial equal to its own $P_{n,a}$?\quad \textbf{311} Taylor approximations at $0$ of even/odd functions.\quad \textbf{312} $P_5$ of $\tan x$ by organised differentiation.\quad \textbf{313} $k$-th derivatives of $\log(1\pm x)$ at $0$; $P_5$ of $\log(1-x)$; substitution $-x$.\quad \textbf{315} From $\sqrt{1-x^2}$: $\frac12\sqrt3\approx\frac{111}{128}$.\quad \textbf{318} No order-$\le4$ terms were dropped in Voorbeeld 3.6.\quad \textbf{321} $P_5$ of $e^{x-\frac16x^3}$.\quad \textbf{322 (Extra)} Approximating $\log\frac12$, $\log0$, $\log1.9$, $\log2.1$, $\log2$ by Taylor polynomials.
\end{trybox}

\hsection{l'H\^opital's rule and some important limits}
\begin{result}[{Voorbeeld 3.7, Stelling 3.8 \textperiodcentered\ Taylor for limits; l'H\^opital}]
\textbf{Taylor:} near $0$, $\dfrac{\sin x}x=\dfrac{x+c_3x^3+\dots}x=1+c_3x^2+\dots\to1$; the exact coefficients do not matter, only that the remaining terms tend to $0$. \textbf{l'H\^opital:} if $f(x)\to0$ and $g(x)\to0$ as $x\to a$ (or both $\to\pm\infty$; $a=\pm\infty$ allowed) and $\lim\frac{f'(x)}{g'(x)}$ exists, then $\lim\frac{f(x)}{g(x)}=\lim\frac{f'(x)}{g'(x)}$. Voorbeeld 3.9: $\lim_{x\to0}\frac x{\tan x}=\lim\frac1{1/\cos^2x}=\lim\cos^2x=1$. \textbf{Check the hypotheses first}: l'H\^opital on a quotient that is not $\frac00$ or $\frac\infty\infty$ gives wrong answers (325d).
\end{result}

\begin{bookex}{323}
Repeat the Taylor argument for $\lim_{x\to0}\dfrac{\sin3x}{2x}$ and $\lim_{x\to0}\dfrac{\sin x}{9x^3}$.
\solution
$\sin3x=3x-\frac{(3x)^3}6+\dots=3x-\frac92x^3+\dots$, so $\dfrac{\sin3x}{2x}=\dfrac32-\dfrac94x^2+\dots\to\dfrac32$.\\
$\dfrac{\sin x}{9x^3}=\dfrac{x-\frac16x^3+\dots}{9x^3}=\dfrac1{9x^2}-\dfrac1{54}+\dots$; the first term tends to $+\infty$, so the limit does not exist as a real number ($+\infty$).\qed
\end{bookex}

\begin{bookex}{325}
Determine in two ways (with and without l'H\^opital): (a) $\lim_{x\to0}\dfrac{\log(1+x)}x$; (b) $\lim_{x\to1}\dfrac{\log x}{1-x}$; (c) $\lim_{x\to0}\dfrac{1-\cos x}{\sin x}$; (d) $\lim_{x\to\pi/2}\dfrac{1-\cos x}{\sin x}$; (e) $\lim_{x\to0}\Big(\dfrac1x-\dfrac1{\sin x}\Big)$.
\solution
(a) $\frac00$. l'H: $\dfrac{1/(1+x)}1\to1$. Taylor: $\dfrac{x-\frac12x^2+\dots}x=1-\frac12x+\dots\to1$. (Also: it is $\log'(1)$.)\\
(b) $\frac00$. l'H: $\dfrac{1/x}{-1}\to-1$. Taylor at $1$: $\log x=\log(1+(x-1))=(x-1)-\frac12(x-1)^2+\dots$, so the quotient is $\dfrac{(x-1)(1-\frac12(x-1)+\dots)}{-(x-1)}\to-1$.\\
(c) $\frac00$. l'H: $\dfrac{\sin x}{\cos x}\to0$. Taylor: $\dfrac{\frac12x^2-\dots}{x-\dots}=\dfrac{\frac12x-\dots}{1-\dots}\to0$.\\
(d) Not indeterminate: numerator $\to1-0=1$, denominator $\to1$; the limit is $1$ by the quotient rule. l'H\^opital would wrongly give $\dfrac{\sin x}{\cos x}\to\infty$.\\
(e) $\dfrac1x-\dfrac1{\sin x}=\dfrac{\sin x-x}{x\sin x}$, a $\frac00$ form. Taylor: $\dfrac{-\frac16x^3+\dots}{x^2-\dots}=\dfrac{-\frac16x+\dots}{1-\dots}\to0$. l'H twice: $\dfrac{\cos x-1}{\sin x+x\cos x}$ (still $\frac00$) $\to\dfrac{-\sin x}{2\cos x-x\sin x}\to\dfrac02=0$.\qed
\end{bookex}

\begin{bookex}{326}
Prove in order: (a) $\lim_{x\to\infty}\dfrac{\log x}x=0$; (b) $\lim_{x\to\infty}\dfrac{(\log x)^n}x=0$ for $n\in\N$; (c) $\lim_{t\to\infty}\dfrac{t^n}{e^t}=0$; (d) $\lim_{t\to\infty}\dfrac{t^x}{e^t}=0$ for all real $x$.
\solution
(a) $\frac\infty\infty$; l'H\^opital: $\dfrac{1/x}1=\dfrac1x\to0$.\\
(b) Induction on $n$; (a) is $n=1$. If $\frac{(\log x)^{n-1}}x\to0$, then l'H\^opital on $\frac{(\log x)^n}x$ gives $\dfrac{n(\log x)^{n-1}\cdot\frac1x}1=n\dfrac{(\log x)^{n-1}}x\to0$.\\
(c) Substitute $t=\log x$, i.e.\ $x=e^t$; as $t\to\infty$, $x\to\infty$, and $\dfrac{t^n}{e^t}=\dfrac{(\log x)^n}x\to0$ by (b).\\
(d) Choose $n\in\N$ with $n\ge x$. For $t\ge1$: $0\le t^x\le t^n$, so $0\le\dfrac{t^x}{e^t}\le\dfrac{t^n}{e^t}\to0$; squeeze.\\
Meaning: $\log$ grows slower than every power, every power grows slower than $e^t$.\qed
\end{bookex}

\begin{trybox}[{324, 327, 328, 329, 330 (Extra)}]
\textbf{324} $\lim_{x\to0}\frac x{\tan x}$ without l'H\^opital.\quad \textbf{327} $\lim_{x\to0}x\log x=0$ from 326(a).\quad \textbf{328, 329 (Extra)} The general terms of the series for $e^x,\sin,\cos,\log(1+x)$; $e^{ix}=\cos x+i\sin x$ from the series.\quad \textbf{330 (Extra)} $\frac1{1+x}=1-x+x^2-\dots$: meaning for $x=0,-\frac12,-1,1$; its Taylor series; differentiating the series of $\log(1+x)$.
\end{trybox}

\begin{warn}
\begin{itemize}
\item Always say which steunpunt and which order. Before substituting a known expansion, check that the inner expression tends to $0$ (split off $e^{1}$ in $e^{\cos2x}$).
\item Keep terms up to the required order and \emph{all} of them: in $\log\cos x$ the $u^2$ term contributes at order $4$ (Oefening 318).
\item l'H\^opital only for $\frac00$ or $\frac\infty\infty$; differentiate numerator and denominator separately (no quotient rule).
\item Degrees to radians before using a Taylor polynomial (Oefening 309).
\end{itemize}
\end{warn}

\chapter{Tapas: four short topics}
\chaptercontents{A & Pascal's triangle (\S4.1) \\ B & Logarithmic differentiation (\S4.2) \\ C & arcsin, arccos, arctan (\S4.3) \\ D & Hyperbolic functions (\S4.4)}%
{Exercises 401--420 (all worked or solved).\\ Hoofdpunten: expanding $(a+b)^n$ with Pascal's triangle; $\arcsin$, $\arccos$, $\arctan$ with domain, graph and derivative; logarithmic differentiation; simplifying expressions with inverse trig functions. The hyperbolic functions need not be memorised: the skill is deriving an inverse and its derivative from a given formula.}

\hsection{Pascal's triangle}
\begin{key}[{\S4.1}]
Rows $1$; $1\,1$; $1\,2\,1$; $1\,3\,3\,1$; $1\,4\,6\,4\,1$; \dots, each number the sum of the two above it. Row $n$ (counting from $0$) gives the coefficients of $(x+y)^n$: $(x+y)^3=x^3+3x^2y+3xy^2+y^3$.
\end{key}

\begin{bookex}{401}
Check this for $n=0,1,2,3$ and predict $(x+y)^7$.
\solution
$(x+y)^0=1$; $(x+y)^1=x+y$; $(x+y)^2=x^2+2xy+y^2$; $(x+y)^3=(x+y)(x^2+2xy+y^2)=x^3+3x^2y+3xy^2+y^3$: coefficients $1$; $1\,1$; $1\,2\,1$; $1\,3\,3\,1$. Continuing the triangle: $1\,5\,10\,10\,5\,1$, $1\,6\,15\,20\,15\,6\,1$, $1\,7\,21\,35\,35\,21\,7\,1$, so
\[
(x+y)^7=x^7+7x^6y+21x^5y^2+35x^4y^3+35x^3y^4+21x^2y^5+7xy^6+y^7 .
\]\qed
\end{bookex}

\begin{trybox}[{402, 403}]
\textbf{402} Why is the sum of row $n$ equal to $2^n$?\quad \textbf{403} How does Pascal's triangle give the $n$-th derivative of $f(x)g(x)$? Try $n=0,1,2,\dots$
\end{trybox}

\hsection{Logarithmic differentiation}
\begin{strategy}[{\S4.2 and Voorbeeld 4.1 \textperiodcentered\ take the log first}]
$\dfrac d{dx}\log f(x)=\dfrac{f'(x)}{f(x)}$, so $f'(x)=f(x)\cdot\big(\log f(x)\big)'$. Useful when $\log f$ is easier to differentiate than $f$: products, quotients, powers and variable exponents become sums, differences and multiples. Voorbeeld 4.1: $f(x)=x^x$, $\log f=x\log x$, $(\log f)'=\log x+1$, so $f'(x)=(\log x+1)x^x$. (Alternative for exponents: $a^{b}=e^{b\log a}$.)
\end{strategy}

\begin{bookex}{404}
Differentiate in two ways (``normally'' and logarithmically): (a) $(6x+1)^5(x^4-3)^6$; (b) $\sqrt[5]{\dfrac{x^2+2}{x^2+3}}$.
\solution
(a) \emph{Logarithmic:} $\log f=5\log(6x+1)+6\log(x^4-3)$, so $\dfrac{f'}f=\dfrac{30}{6x+1}+\dfrac{24x^3}{x^4-3}$ and $f'=(6x+1)^5(x^4-3)^6\Big(\dfrac{30}{6x+1}+\dfrac{24x^3}{x^4-3}\Big)$. \emph{Normal:} product and chain rule, $f'=30(6x+1)^4(x^4-3)^6+24x^3(6x+1)^5(x^4-3)^5$. Both equal $6(6x+1)^4(x^4-3)^5\big(5(x^4-3)+4x^3(6x+1)\big)=6(6x+1)^4(x^4-3)^5(29x^4+4x^3-15)$.\\
(b) \emph{Logarithmic:} $\log g=\frac15\big(\log(x^2+2)-\log(x^2+3)\big)$, $\dfrac{g'}g=\dfrac15\Big(\dfrac{2x}{x^2+2}-\dfrac{2x}{x^2+3}\Big)=\dfrac{2x}{5(x^2+2)(x^2+3)}$, so $g'=\sqrt[5]{\dfrac{x^2+2}{x^2+3}}\cdot\dfrac{2x}{5(x^2+2)(x^2+3)}$. \emph{Normal:} $g'=\dfrac15\Big(\dfrac{x^2+2}{x^2+3}\Big)^{-4/5}\cdot\dfrac{2x(x^2+3)-2x(x^2+2)}{(x^2+3)^2}=\dfrac15\Big(\dfrac{x^2+2}{x^2+3}\Big)^{-4/5}\dfrac{2x}{(x^2+3)^2}$, the same after writing $(\cdot)^{-4/5}=(\cdot)^{1/5}\cdot\frac{x^2+3}{x^2+2}$.\qed
\end{bookex}

\begin{bookex}{405}
Find $f'$ for (a) $f:x\mapsto a^{\sin(bx)}$ ($a>0$, $b\ne0$); (b) $f:x\mapsto\big(\log(4x+3)\big)^{2x}$.
\solution
(a) $\log f=\sin(bx)\log a$, so $f'=f\cdot b\cos(bx)\log a=b\log a\cos(bx)\,a^{\sin(bx)}$.\\
(b) $\log f=2x\log\big(\log(4x+3)\big)$. Product and chain rule: $(\log f)'=2\log\big(\log(4x+3)\big)+2x\cdot\dfrac1{\log(4x+3)}\cdot\dfrac4{4x+3}$, so
\[
f'(x)=\big(\log(4x+3)\big)^{2x}\Big(2\log\big(\log(4x+3)\big)+\frac{8x}{(4x+3)\log(4x+3)}\Big).
\]\qed
\end{bookex}

\begin{bookex}{407}
Find $f'(0)$ if $f(x)=(1-x)(1-2x)^2(1-3x)^3(1-4x)^4(1-5x)^5$.
\solution
$\log f=\sum_{k=1}^5k\log(1-kx)$, so $\dfrac{f'}f=\sum_{k=1}^5\dfrac{-k^2}{1-kx}$. At $x=0$: $f(0)=1$ and $f'(0)=-(1+4+9+16+25)=-55$. (Expanding the product would be madness.)\qed
\end{bookex}

\begin{trybox}[{406, 408}]
\textbf{406} Differentiate $x^{(e^x)}$ and $e^{(x^e)}$.\quad \textbf{408} What do the examples have in common that makes logarithmic differentiation work so well, and how would you do them without it?
\end{trybox}

\hsection{Circle functions: arcsin, arccos, arctan}
\begin{key}[{Tabel 4.1 \textperiodcentered\ inverse trigonometric functions}]
$\sin,\cos,\tan$ are periodic, hence not bijective; restrict the domain to make them bijective and invert:
\begin{center}\small
\begin{tabular}{l|c|l|c}
function & restricted domain & inverse & domain of the inverse\\\hline
$\sin$ & $[-\frac\pi2,\frac\pi2]$ & $\arcsin$ & $[-1,1]$, values in $[-\frac\pi2,\frac\pi2]$\\
$\cos$ & $[0,\pi]$ & $\arccos$ & $[-1,1]$, values in $[0,\pi]$\\
$\tan$ & $(-\frac\pi2,\frac\pi2)$ & $\arctan$ & $\R$, values in $(-\frac\pi2,\frac\pi2)$
\end{tabular}
\end{center}
Graphs: mirror the restricted graph in the line $y=x$. The American notation $\sin^{-1}$ is avoided (confusion with $1/\sin$). \textbf{Derivatives} (from $\sin(\arcsin x)=x$ by the chain rule): $\arcsin'x=\dfrac1{\sqrt{1-x^2}}$, $\arccos'x=\dfrac{-1}{\sqrt{1-x^2}}$, $\arctan'x=\dfrac1{1+x^2}$ (Oefening 412).
\end{key}
\begin{strategy}[{Oefeningen 410, 411 \textperiodcentered\ simplifying}]
$\arcsin(\sin\alpha)=\alpha$ \emph{only} for $\alpha\in[-\frac\pi2,\frac\pi2]$; otherwise first rewrite $\sin\alpha$ as the sine of an angle in that interval using the unit circle. For $\sin(\arccos x)$ and the like, draw a right triangle with hypotenuse $1$ and an angle $\varphi=\arccos x$: adjacent side $x$, opposite side $\sqrt{1-x^2}$; for $\arctan x$ take legs $1$ and $x$, hypotenuse $\sqrt{1+x^2}$. Then check the sign for negative $x$.
\end{strategy}

\begin{bookex}{410}
(a) Why is $\arcsin(\sin\alpha)=\alpha$ for $-\frac\pi2\le\alpha\le\frac\pi2$? (b) Simplify $\arcsin(\sin(\pi+\alpha))$ for such $\alpha$. (c) Largest interval $I$ on which $\arccos(\cos\beta)=\beta$. (d) Simplify $\arccos(\cos(\zeta-\pi))$ for $0\le\zeta\le\pi$. (e) Simplify $\arctan(-\tan(\frac\pi2+\eta))$ for $\frac\pi2\le\eta\le\pi$. (f) An angle $\lambda\in[-\pi,\pi]$ with $\tan\lambda=-\frac13\sqrt3$ but $\arctan(\tan\lambda)\ne\lambda$.
\solution
(a) On $[-\frac\pi2,\frac\pi2]$ the sine is bijective and $\arcsin$ is defined as its inverse there, so $\arcsin$ undoes $\sin$: $\arcsin(\sin\alpha)=\alpha$.\\
(b) The point at angle $\pi+\alpha$ is the point at angle $\alpha$ rotated by a half turn, so its height is $\sin(\pi+\alpha)=-\sin\alpha=\sin(-\alpha)$. Since $-\alpha\in[-\frac\pi2,\frac\pi2]$, $\arcsin$ returns $-\alpha$: $\arcsin(\sin(\pi+\alpha))=-\alpha$.\\
(c) $I=[0,\pi]$, the restricted domain of $\cos$ (for $\beta$ outside it, $\arccos(\cos\beta)$ is the angle in $[0,\pi]$ with the same cosine, not $\beta$).\\
(d) $\cos(\zeta-\pi)=-\cos\zeta=\cos(\pi-\zeta)$ (reflection in the $y$-axis), and $\pi-\zeta\in[0,\pi]$. So $\arccos(\cos(\zeta-\pi))=\pi-\zeta$.\\
(e) $-\tan(\frac\pi2+\eta)=\tan(-\frac\pi2-\eta)=\tan(\frac\pi2-\eta)$ ($\tan$ is odd and has period $\pi$). For $\frac\pi2\le\eta<\pi$, $\frac\pi2-\eta\in(-\frac\pi2,0]$, inside the restricted domain, so the expression equals $\frac\pi2-\eta$. (At $\eta=\pi$, $\tan\frac{3\pi}2$ is undefined.)\\
(f) $\tan\lambda=-\frac13\sqrt3$ for $\lambda=-\frac\pi6+k\pi$. In $[-\pi,\pi]$: $\lambda=-\frac\pi6$ (then $\arctan(\tan\lambda)=\lambda$) or $\lambda=\frac{5\pi}6$, for which $\arctan(\tan\frac{5\pi}6)=\arctan(-\frac13\sqrt3)=-\frac\pi6\ne\frac{5\pi}6$. So $\lambda=\frac{5\pi}6$.\qed
\end{bookex}

\begin{bookex}{411}
(a) With $\arccos x=\varphi$ and hypotenuse $1$, express the legs in $x$. (b) Simplify $\sin(\arccos x)$. (c) Check the simplification for obtuse $\varphi$ ($x<0$). (d) Simplify $\cos(\arcsin x)$; what changes, and does it hold for $x<0$? (e) Simplify $\cos(\arctan x)$. (f) Simplify $\tan(\arccos x)$ and check for $x<0$.
\solution
(a) $\cos\varphi=x$, so the adjacent leg is $x$ and the opposite leg is $\sqrt{1-x^2}$ (Pythagoras).\\
(b) $\sin(\arccos x)=\sin\varphi=\sqrt{1-x^2}$.\\
(c) For $\frac\pi2<\varphi<\pi$, $x=\cos\varphi<0$; the triangle in the right figure has adjacent leg $|x|=-x$, but its height is still $\sqrt{1-x^2}$, and $\sin\varphi>0$ for $\varphi\in[0,\pi]$. So $\sin(\arccos x)=\sqrt{1-x^2}$ for all $x\in[-1,1]$.\\
(d) Now $\varphi=\arcsin x$: opposite leg $x$, adjacent leg $\sqrt{1-x^2}$: $\cos(\arcsin x)=\sqrt{1-x^2}$. For $x<0$, $\varphi\in[-\frac\pi2,0)$ and $\cos\varphi\ge0$, so the positive root is right.\\
(e) $\varphi=\arctan x$: take the adjacent leg $1$ and the opposite leg $x$; hypotenuse $\sqrt{1+x^2}$; $\cos(\arctan x)=\dfrac1{\sqrt{1+x^2}}$ (positive, as $\varphi\in(-\frac\pi2,\frac\pi2)$).\\
(f) $\tan(\arccos x)=\dfrac{\sin\varphi}{\cos\varphi}=\dfrac{\sqrt{1-x^2}}x$ ($x\ne0$). For $x<0$, $\varphi$ is obtuse and $\tan\varphi<0$; the formula is negative too. \checkmark\qed
\end{bookex}

\begin{bookex}{412}
Show in the same way that $\arccos'x=\dfrac{-1}{\sqrt{1-x^2}}$ and $\arctan'x=\dfrac1{1+x^2}$.
\solution
Differentiate $\cos(\arccos x)=x$: $-\sin(\arccos x)\cdot\arccos'x=1$, so $\arccos'x=\dfrac{-1}{\sin(\arccos x)}=\dfrac{-1}{\sqrt{1-x^2}}$ by 411(b).\\
Differentiate $\tan(\arctan x)=x$ with $\tan'=1+\tan^2$: $\big(1+\tan^2(\arctan x)\big)\arctan'x=1$, i.e.\ $(1+x^2)\arctan'x=1$, so $\arctan'x=\dfrac1{1+x^2}$.\qed
\end{bookex}

\begin{trybox}[{409, 413, 414 (Extra)}]
\textbf{409} Fill in Tabel 4.2 ($\arcsin t$, $\arccos t$, $\arctan t$ for the standard values $t$); mark ``not nice'' versus ``not defined''.\quad \textbf{413} Check that the derivatives match the shapes of the graphs.\quad \textbf{414 (Extra)} Locate $\sec t$, $\csc t$, $\cot t$ in the unit circle.
\end{trybox}

\hsection{Hyperbolic functions}
\begin{key}[{\S4.4}]
$\sinh t=\dfrac{e^t-e^{-t}}2$, $\cosh t=\dfrac{e^t+e^{-t}}2$. Properties to \emph{derive}, not memorise: $\sinh$ is odd with only zero $0$, $\cosh$ is even and $\ge1$; $\sinh'=\cosh$, $\cosh'=\sinh$; $\cosh^2t-\sinh^2t=1$ (Oefeningen 415--418). $\sinh$ is bijective $\R\to\R$; its inverse $\operatorname{arsinh}x=\log\big(x+\sqrt{x^2+1}\big)$ has derivative $\dfrac1{\sqrt{x^2+1}}$ (Oefeningen 419, 420; compare $\arcsin'$). Chapter 7 has the same for $\cosh$ and $\tanh$.
\end{key}

\begin{bookex}{419}
Let $\operatorname{arsinh}x=t$. (a) Check $x=\frac12(e^t-e^{-t})$. (b) Substitute $e^t=u$ and reduce to an expression in $x$ and $u$. (c) Solve the resulting quadratic equation for $u$. (d) Conclude $\operatorname{arsinh}x=\log(x+\sqrt{x^2+1})$ and explain why $\log(x-\sqrt{x^2+1})$ is not an option.
\solution
(a) $t=\operatorname{arsinh}x$ means $x=\sinh t=\frac12(e^t-e^{-t})$.\quad (b) With $u=e^t$ ($u>0$) and $e^{-t}=\frac1u$: $x=\frac12\big(u-\frac1u\big)$.\\
(c) Multiply by $2u$: $2xu=u^2-1$, i.e.\ $u^2-2xu-1=0$, so $u=\dfrac{2x\pm\sqrt{4x^2+4}}2=x\pm\sqrt{x^2+1}$.\\
(d) $u=e^t$ must be positive. Since $\sqrt{x^2+1}>|x|$, the root $x-\sqrt{x^2+1}$ is negative for every $x$ and $\log$ of it does not exist; the other root $x+\sqrt{x^2+1}$ is positive. Hence $t=\log\big(x+\sqrt{x^2+1}\big)$.\qed
\end{bookex}

\begin{bookex}{420}
Differentiate the expression found in 419; the result should resemble $\arcsin'$.
\solution
Chain rule with $\big(\sqrt{x^2+1}\big)'=\dfrac{x}{\sqrt{x^2+1}}$:
\[
\operatorname{arsinh}'x=\frac{1+\frac{x}{\sqrt{x^2+1}}}{x+\sqrt{x^2+1}}=\frac{\frac{\sqrt{x^2+1}+x}{\sqrt{x^2+1}}}{x+\sqrt{x^2+1}}=\frac1{\sqrt{x^2+1}}.
\]
Compare $\arcsin'x=\dfrac1{\sqrt{1-x^2}}$: only the sign inside the root differs.\qed
\end{bookex}

\begin{trybox}[{415, 416, 417, 418}]
\textbf{415} Even/odd and zeros of $\sinh$, $\cosh$.\quad \textbf{416} Sketch $e^x$, $e^{-x}$, $\sinh x$, $\cosh x$ in one figure.\quad \textbf{417} Derivatives of $\sinh$, $\cosh$; compare with $\sin$, $\cos$.\quad \textbf{418} $\cosh t-\sinh t$, $\cosh t+\sinh t$, and hence $\cosh^2t-\sinh^2t=1$.
\end{trybox}

\begin{warn}
\begin{itemize}
\item $\arcsin(\sin\alpha)=\alpha$ only on $[-\frac\pi2,\frac\pi2]$, $\arccos(\cos\beta)=\beta$ only on $[0,\pi]$, $\arctan(\tan\lambda)=\lambda$ only on $(-\frac\pi2,\frac\pi2)$. Draw the circle.
\item $\arcsin$ and $\arccos$ take only $t\in[-1,1]$; $\arctan(\frac12)$ exists but has no nice value.
\item Logarithmic differentiation needs $f>0$ (or use $\log|f|$); the final answer must be multiplied back by $f$.
\end{itemize}
\end{warn}

\chapter{Differential equations}
\chaptercontents{A & First acquaintance (\S5.1) \\ B & Homogeneous linear equations (\S5.2.1) \\ C & Inhomogeneous equations (\S5.2.2) \\ D & Initial value problems (\S5.2.3)}%
{Exercises 501--524 (all worked or solved).\\ Hoofdpunten: what a d.v.\ is; order, (in)homogeneous, initial value, integration constant, homogeneous/particular/general solution; recognising linearity; solving a homogeneous d.v.\ by separating variables; solving an inhomogeneous first-order linear initial value problem; checking that a function solves a d.v.}

\hsection{First acquaintance}
\begin{key}[{\S5.1, Definition 5.3 \textperiodcentered\ what a differential equation is}]
A \term{differential equation} (d.v.) relates a function to one or more of its derivatives. Voorbeeld 5.1 (radioactive decay): $y'=ky$, solved by $y(t)=ce^{kt}$ for every $c$ (a whole family; $c$ is fixed by the initial amount). Voorbeeld 5.2 (falling with air resistance): $m\,\dfrac{dv}{dt}=mg-kv$; without solving: $v$ increases, ever more slowly, and cannot exceed $\frac{mg}k$; the solution is $v(t)=\frac{mg}k\big(1-e^{-kt/m}\big)$. The \term{order} is the highest derivative that occurs. Solving a d.v.\ means finding \emph{all} functions that satisfy it; often impossible in closed form, but \textbf{checking} a proposed solution is always possible: substitute it. When the variable is clear one writes $u''+6u'+5u=10x+2$ for $u''(x)+\dots$.
\end{key}

\begin{bookex}{502}
For $y(t)=ce^{kt}$: (a) find $c$ (grams) if you start at $t=0$ with $3$ g of Radium-226; (b) find $k$ if after $100$ years $2.87281$ g is left; (c) find the half-life; (d) what do $c=0$ and $k=0$ mean physically?
\solution
(a) $y(0)=ce^0=c=3$.\\
(b) $3e^{100k}=2.87281$, so $k=\dfrac1{100}\log\dfrac{2.87281}3\approx-4.33\cdot10^{-4}$ per year (negative: decay, Oefening 501).\\
(c) The half-life $T$ satisfies $e^{kT}=\frac12$, so $T=\dfrac{\log\frac12}k=\dfrac{\log2}{4.33\cdot10^{-4}}\approx1600$ years (independent of the starting amount).\\
(d) $c=0$: there is no radium at all ($y\equiv0$). $k=0$: $y$ is constant, the substance does not decay (a stable isotope).\qed
\end{bookex}

\begin{bookex}{506}
Check that the following solve $h'=h^2$: (a) $t\mapsto0$ on $\R$; (b) $t\mapsto-\frac1t$ on $(0,\infty)$; (c) $t\mapsto\frac1{1-t}$ on $(-\infty,1)$.
\solution
(a) $h'=0=0^2$. (b) $h'=\frac1{t^2}=\big(-\frac1t\big)^2=h^2$. (c) $h'=\frac{1}{(1-t)^2}=h^2$ (chain rule: $\frac d{dt}(1-t)^{-1}=-(1-t)^{-2}\cdot(-1)$). Note that one d.v.\ has very different solutions, each on its own domain.\qed
\end{bookex}

\begin{trybox}[{501, 503, 504, 505, 507}]
\textbf{501} What do you know about $k$ in $y'=ky$ for decay?\quad \textbf{503} Orders of (5.2)--(5.7).\quad \textbf{504} Check $v(t)=\frac{mg}k(1-e^{-kt/m})$ in Voorbeeld 5.2.\quad \textbf{505} Solutions of $g'^2+g^2=1$: the constant $-1$, and more.\quad \textbf{507} Check $2x-2$ and $c_1e^{-x}+c_2e^{-5x}+2(x-1)$ in $u''+6u'+5u=10x+2$.
\end{trybox}

\hsection{Linear first-order equations: the homogeneous case}
\begin{key}[{Definitions 5.4, 5.5 \textperiodcentered\ linear, homogeneous}]
A \term{linear first-order d.v.} can be written as $y'+a(x)y=f(x)$ with known $a,f$; it is \term{homogeneous} if $f\equiv0$. (General $n$-th order linear: $y^{(n)}+a_{n-1}(x)y^{(n-1)}+\dots+a_0(x)y=f(x)$.) $g'^2+g^2=1$ and $h'=h^2$ are not linear; $j'-xj=0$ is linear, homogeneous, first order; $u''+6u'+5u=10x+2$ is linear, inhomogeneous, second order. The zero function always solves the homogeneous equation (Oefening 508).
\end{key}
\begin{strategy}[{\S5.2.1 \textperiodcentered\ separating variables}]
For $\dfrac{dy}{dx}+2(x+1)y=0$: put everything with $y$ on one side and everything with $x$ on the other, $\dfrac{dy}y=-2(x+1)\,dx$; integrate both sides, $\log|y|=-x^2-2x+C_1$; exponentiate, $y=C_2e^{-x^2-2x}$. Remarks: one integration constant suffices (a constant on each side is just their difference); $C_2=e^{C_1}$ is renamed to a new constant; the absolute value disappears because $C_2$ may be negative, and $C_2=0$ gives the zero solution. \textbf{General rule (Oefening 512):} $y'+a(x)y=0$ has general solution $y=Ce^{-A(x)}$ with $A$ a primitive of $a$. Whatever manipulations you use, \textbf{justify the result by substituting it into the d.v.}
\end{strategy}

\begin{bookex}{511}
Solve the homogeneous equations: (a) $y'+3y=0$; (b) $y'-xy=0$; (c) $y'=-\dfrac{2y}x$; (d) $4y=x^2y'$; (e) $y'+\sin(x)\,y=0$.
\solution
(a) $\frac{dy}y=-3\,dx$, $\log|y|=-3x+C_1$: $y=Ce^{-3x}$. Check: $-3Ce^{-3x}+3Ce^{-3x}=0$.\\
(b) $\frac{dy}y=x\,dx$, $\log|y|=\frac12x^2+C_1$: $y=Ce^{x^2/2}$. Check: $y'=xCe^{x^2/2}=xy$.\\
(c) $\frac{dy}y=-\frac{2\,dx}x$, $\log|y|=-2\log|x|+C_1$: $y=Cx^{-2}=\dfrac C{x^2}$. Check: $y'=-2Cx^{-3}=-\frac2x\cdot\frac C{x^2}$.\\
(d) $y'=\frac{4y}{x^2}$, $\frac{dy}y=\frac4{x^2}dx$, $\log|y|=-\frac4x+C_1$: $y=Ce^{-4/x}$. Check: $y'=\frac4{x^2}Ce^{-4/x}$, and $x^2y'=4Ce^{-4/x}=4y$.\\
(e) $\frac{dy}y=-\sin x\,dx$, $\log|y|=\cos x+C_1$: $y=Ce^{\cos x}$. Check: $y'=-\sin x\,Ce^{\cos x}=-\sin(x)\,y$.\qed
\end{bookex}

\begin{trybox}[{508, 509, 510, 512, 513}]
\textbf{508} The zero function solves $y'+a(x)y=0$.\quad \textbf{509} Check the manipulation that separates the variables.\quad \textbf{510} Substitute $y=Ce^{-x^2-2x}$ (all $C$, also $0$).\quad \textbf{512} General rule for $y'+a(x)y=0$.\quad \textbf{513} Check your solutions of 511 (done above).
\end{trybox}

\hsection{The inhomogeneous equation}
\begin{result}[{Definition 5.6, Stelling 5.7 \textperiodcentered\ general = homogeneous + particular}]
A \term{particular solution} $y_p$ is any one solution of $y'+a(x)y=f(x)$. If $y_h$ is the general solution of the homogeneous equation and $y_p$ a particular solution, then $y=y_h+y_p$ is the general solution of the inhomogeneous equation (Oefening 514: the sum solves it, and two particular solutions differ by a homogeneous solution). So: first solve the homogeneous equation, then find \emph{one} $y_p$.
\end{result}
\begin{strategy}[{\S5.2.2 \textperiodcentered\ three ways to a particular solution}]
\textbf{Guessing (gokken met voorkennis):} for $y'+3y=x+2$ try $y_p=ax+b$: $3ax+3b+a=x+2$ gives $a=\frac13$, $b=\frac59$; general solution $y=Ce^{-3x}+\frac13x+\frac59$. Try the same type as $f$: polynomial, $Ae^{5x}$ for $e^{5x}$, constants for constants.\\
\textbf{Integrating factor:} multiply $y'+ay=f$ by $u=e^{\int a\,dx}$; since $u'=au$, the left side becomes $(uy)'$, so $uy=\int uf\,dx$ and $y=\dfrac1u\int uf\,dx$. Voorbeeld 5.8: $y'+\frac2xy=2x+1$, $u=e^{2\log x}=x^2$, $(x^2y)'=2x^3+x^2$, $y=\frac12x^2+\frac13x+\frac C{x^2}$.\\
\textbf{Variation of constants:} replace $C$ in $y_h=Ce^{-A(x)}$ by $C(x)$ and substitute: the $C(x)$-terms cancel and $C'(x)e^{-A(x)}=f(x)$ remains, so $C(x)=\int f(x)e^{A(x)}dx$ (Oefening 518). Example: $xy'+4\frac yx=\frac1x$ has $y_h=Ce^{4/x}$; $C'(x)=\frac1{x^2}e^{-4/x}$, $C(x)=\frac14e^{-4/x}+c$, $y=\frac14+ce^{4/x}$.\\
Any method is allowed on the exam; \textbf{always check the result in the equation.}
\end{strategy}

\begin{bookex}{516}
Use an integrating factor to solve: (a) $y'+2y=3$; (b) $y'+2xy=3e^{-x^2}$.
\solution
(a) $a=2$, $u=e^{2x}$: $(e^{2x}y)'=3e^{2x}$, so $e^{2x}y=\frac32e^{2x}+C$ and $y=\frac32+Ce^{-2x}$. Check: $y'+2y=-2Ce^{-2x}+3+2Ce^{-2x}=3$.\\
(b) $a=2x$, $u=e^{x^2}$: $(e^{x^2}y)'=e^{x^2}\cdot3e^{-x^2}=3$, so $e^{x^2}y=3x+C$ and $y=(3x+C)e^{-x^2}$. Check: $y'=3e^{-x^2}-2x(3x+C)e^{-x^2}$ and $2xy=2x(3x+C)e^{-x^2}$; sum $3e^{-x^2}$.\qed
\end{bookex}

\begin{bookex}{517}
Solve by a method of your choice: (a) $y'+xy=x$; (b) $y'-3y=e^{5x}$; (c) $y'+\dfrac{2xy}{1+x^2}=2x$; (d) $xy'+(1+x^2)y=x$.
\solution
(a) Homogeneous: $y_h=Ce^{-x^2/2}$ (511b with the sign changed). Guess a constant: $y_p=1$ gives $0+x=x$. General solution $y=1+Ce^{-x^2/2}$.\\
(b) $y_h=Ce^{3x}$. Guess $y_p=Ae^{5x}$: $5Ae^{5x}-3Ae^{5x}=e^{5x}$ gives $A=\frac12$. So $y=Ce^{3x}+\frac12e^{5x}$.\\
(c) $a=\dfrac{2x}{1+x^2}$ has primitive $\log(1+x^2)$, so the integrating factor is $u=1+x^2$: $\big((1+x^2)y\big)'=2x(1+x^2)=2x+2x^3$, hence $(1+x^2)y=x^2+\frac12x^4+C$ and $y=\dfrac{x^2+\frac12x^4+C}{1+x^2}$.\\
(d) Divide by $x$: $y'+\big(\frac1x+x\big)y=1$. Primitive of $a$: $\log|x|+\frac12x^2$, so $u=xe^{x^2/2}$ (for $x>0$): $\big(xe^{x^2/2}y\big)'=xe^{x^2/2}$, hence $xe^{x^2/2}y=e^{x^2/2}+C$ and $y=\dfrac1x+\dfrac{C}{x}e^{-x^2/2}=\dfrac{1+Ce^{-x^2/2}}x$. Check: with $y=\frac1x+\frac Cxe^{-x^2/2}$ one finds $y'=-\frac1{x^2}-\frac C{x^2}e^{-x^2/2}-Ce^{-x^2/2}$, so $xy'+(1+x^2)y=\big(-\frac1x-\frac Cxe^{-x^2/2}-Cxe^{-x^2/2}\big)+\big(\frac1x+x+\frac Cxe^{-x^2/2}+Cxe^{-x^2/2}\big)=x$. \checkmark\qed
\end{bookex}

\begin{trybox}[{514, 515, 518}]
\textbf{514} Prove Stelling 5.7: (a) $y_h+y_p$ solves the inhomogeneous equation; (b) $\tilde y_p-y_p$ solves the homogeneous one; (c) why does this prove the theorem?\quad \textbf{515} Find the integrating factor $x^2$ of Voorbeeld 5.8 from $e^{\int a}$.\quad \textbf{518} Variation of constants for general $y'+a(x)y=f(x)$.
\end{trybox}

\hsection{Initial value problems}
\begin{key}[{Definition 5.9, Voorbeeld 5.10 \textperiodcentered\ fixing the constant}]
The general solution is a family (one integration constant for a first-order equation). An \term{initial value} $y(a)=b$ picks one member: substitute $x=a$ into the \emph{general} solution $y=y_h+y_p$ and solve for $C$ (Oefening 521). Voorbeeld 5.10: $y'+\frac2xy=2x+1$, $y(-1)=\frac23$: from $\frac23=\frac12-\frac13+C$ we get $C=\frac12$.
\end{key}

\begin{bookex}{519}
$x^2y'+4y=1$ has $y_h=Ce^{4/x}$ and $y_p=\frac14$. Solve the initial value problem with $y(4)=1$.
\solution
General solution $y=\frac14+Ce^{4/x}$. At $x=4$: $1=\frac14+Ce^{1}$, so $C=\dfrac3{4e}$ and $y=\dfrac14+\dfrac3{4e}e^{4/x}=\dfrac14+\dfrac34e^{4/x-1}$. Check: $y(4)=\frac14+\frac34=1$ \checkmark; and $x^2y'=x^2\cdot\frac34e^{4/x-1}\cdot(-\frac4{x^2})=-3e^{4/x-1}$, $4y=1+3e^{4/x-1}$, sum $1$. \checkmark\qed
\end{bookex}

\begin{bookex}{522}
A pizza leaves the freezer at 8 o'clock at $-18^\circ$C; kitchen temperature $T_k=20^\circ$C; model $\dfrac{dT}{dt}=-p(T-T_k)$, $p>0$. (a) Argue that $T$ approaches $T_k$ quickly at first, then ever more slowly. (b) At 9 o'clock the pizza is $-9^\circ$C; solve the d.v.\ and find $p$. (c) When is the pizza thawed ($0^\circ$C)? (d) How warm must the kitchen be for the pizza to be thawed at 9 o'clock?
\solution
(a) While $T<T_k$ the right side $-p(T-T_k)$ is positive, so $T$ rises; as $T$ gets closer to $T_k$ the difference $T-T_k$ shrinks, so $T'$ shrinks: slower and slower, never passing $T_k$.\\
(b) Linear: $T'+pT=pT_k$; $T_h=Ce^{-pt}$, $T_p=T_k$ (constant guess). So $T(t)=20+Ce^{-pt}$ with $t$ in hours after 8. $T(0)=-18$ gives $C=-38$: $T(t)=20-38e^{-pt}$. $T(1)=-9$: $38e^{-p}=29$, $e^{-p}=\frac{29}{38}$, $p=\log\frac{38}{29}\approx0.270$ per hour.\\
(c) $T=0$: $38e^{-pt}=20$, $t=\dfrac{\log(38/20)}{p}=\dfrac{\log1.9}{\log(38/29)}\approx\dfrac{0.642}{0.270}\approx2.37$ hours: about 10:22.\\
(d) With kitchen temperature $T_k$: $T(t)=T_k+(-18-T_k)e^{-pt}$ (same $p$, a property of the pizza). $T(1)=0$: $T_k+(-18-T_k)\cdot\frac{29}{38}=0$, so $T_k\big(1-\frac{29}{38}\big)=18\cdot\frac{29}{38}$, $T_k=\frac{18\cdot29}{9}=58^\circ$C.\qed
\end{bookex}

\begin{bookex}{523}
Solve by separation: $y'=-2y^3$, $y(1)=\frac13$.
\solution
Not linear, but separable: $\dfrac{dy}{y^3}=-2\,dx$, so $-\dfrac1{2y^2}=-2x+C_1$, i.e.\ $\dfrac1{y^2}=4x+c$. The initial value: $9=4+c$, $c=5$, so $y^2=\dfrac1{4x+5}$ and, since $y(1)>0$, $y=\dfrac1{\sqrt{4x+5}}$ (for $x>-\frac54$). Check: $y'=-\frac12(4x+5)^{-3/2}\cdot4=-2(4x+5)^{-3/2}=-2y^3$ and $y(1)=\frac13$. \checkmark\qed
\end{bookex}

\begin{trybox}[{520, 521, 524}]
\textbf{520} Derive the solution of Voorbeeld 5.2 ($mv'=mg-kv$, $v(0)=0$).\quad \textbf{521} Which equation determines $C$ in an initial value problem?\quad \textbf{524} Logistic equation $y'=\alpha y(\beta-y)$: constant solutions; $\lim_{t\to\infty}y(t)$ for $y(0)>0$ without solving.
\end{trybox}

\begin{warn}
\begin{itemize}
\item Bring the equation to the form $y'+a(x)y=f(x)$ first (divide by the coefficient of $y'$, as in 517d); the integrating factor is $e^{\int a}$, not $e^{\int f}$.
\item The initial condition is applied to the \emph{general} solution $y_h+y_p$, not to $y_h$ alone (Oefening 521).
\item Every solution must be checked by substitution; the symbolic manipulations with $dy$ and $dx$ are justified only by that check.
\item Keep the integration constant until the end, and keep it even when it turns out to be $0$ (the zero solution).
\end{itemize}
\end{warn}

\chapter{Finding primitives}
\chaptercontents{A & Primitives, improper integrals (\S6.1) \\ B & Integration by parts (\S6.2) \\ C & Integration by substitution (\S6.3) \\ D & Partial fractions (\S6.4) \\ E & Mixed exercises (\S6.5) \\ F & (Extra) $\log$ and $\exp$ (\S6.6)}%
{Exercises 601--635 (all worked or solved).\\ Hoofdpunten: lots of practice with parts, substitution, partial fractions and rewriting integrands, also in combination; willingness to try something else when an approach fails; the double-angle formulas as a tool; what an improper integral is. \textbf{Always check a primitive by differentiating.}}

\hsection{Standard primitives and improper integrals}
\begin{key}[{\S6.1 \textperiodcentered\ know these by heart (constants omitted)}]
\[
\int x^r\,dx=\frac{x^{r+1}}{r+1}\ (r\ne-1),\quad\int\frac{dx}x=\log|x|,\quad\int e^x\,dx=e^x,\quad\int\sin x\,dx=-\cos x,\quad\int\cos x\,dx=\sin x,
\]
\[
\int\frac{dx}{\sqrt{1-x^2}}=\arcsin x,\qquad\int\frac{dx}{1+x^2}=\arctan x.
\]
Differentiating is automatic; primitiving needs creativity, which you get by trying, failing and learning (``goede keuzes maak je met ervaring, ervaring maak je met foute keuzes'').
\end{key}
\begin{key}[{\S6.1.1 and Voorbeelden 6.1--6.4 \textperiodcentered\ improper integrals}]
An integral is \term{improper} if the interval or the integrand is unbounded; evaluate it with a limit: $\int_1^\infty x^{-2}dx=\lim_{b\to\infty}\big[-x^{-1}\big]_1^b=1$ (\term{convergent}); $\int_0^1x^{-2}dx=\big[-x^{-1}\big]_0^1=-1-\lim_{x\to0^+}(-\tfrac1x)$ does not exist (\term{divergent}). Two problems (e.g.\ $\int_0^\infty x^{-2}dx$, or $\int_{-\infty}^\infty$): split the interval and treat each limit separately; the whole integral converges only if every piece does. $\int_{-\infty}^\infty\frac{dx}{1+x^2}=\frac\pi2+\frac\pi2=\pi$.
\end{key}

\begin{bookex}{605}
Convergent? If so, compute: (a) $\int_0^1\dfrac{dx}{\sqrt x}$; (b) $\int_1^\infty\dfrac{dx}{\sqrt x}$.
\solution
A primitive of $x^{-1/2}$ is $2\sqrt x$.\\
(a) Improper at $0$ (integrand unbounded): $\int_0^1x^{-1/2}dx=2\sqrt1-\lim_{a\to0^+}2\sqrt a=2-0=2$. Convergent, value $2$.\\
(b) Improper at $\infty$: $\int_1^bx^{-1/2}dx=2\sqrt b-2\to\infty$ as $b\to\infty$. Divergent. (Compare $x^{-2}$: at infinity, a higher power in the denominator is needed; near $0$, a lower one.)\qed
\end{bookex}

\begin{trybox}[{601, 602, 603, 604, 606}]
\textbf{601} Why does $\int c\,dx=cx$ follow from the list?\quad \textbf{602} $\frac1x$ and $\log|x|$ on both sides of $0$; $\int_a^b\frac{dx}x=\log|b|-\log|a|$ for $a<b<0$, sign check.\quad \textbf{603} $\int\frac{p}{qx+r}\,dx$.\quad \textbf{604} Find $f,g\not\equiv0$ with $\int f\cdot\int g=\int fg$.\quad \textbf{606} $\lim_{a\to\infty}\int_{-a}^a\sin x\,dx$; yet $\int_{-\infty}^\infty\sin x\,dx$ diverges.
\end{trybox}

\hsection{Integration by parts}
\begin{key}[{\S6.2 \textperiodcentered\ partieel integreren}]
From the product rule: $\displaystyle\int f'(x)g(x)\,dx=f(x)g(x)-\int f(x)g'(x)\,dx$, also written $\int g\,df=fg-\int f\,dg$; with limits, $\int_a^bf'g\,dx=\big[fg\big]_a^b-\int_a^bfg'\,dx$. Split the integrand sensibly into a factor $f'$ you can primitive and a factor $g$ that becomes simpler when differentiated. Voorbeeld 6.5: $\int xe^x\,dx$ with $f'=e^x$, $g=x$: $xe^x-\int e^x\,dx=(x-1)e^x+C$. Voorbeeld 6.6: $\int e^x\sin x\,dx$: twice by parts brings back the original integral: $I=-e^x\cos x+e^x\sin x-I$, so $I=\frac12e^x(\sin x-\cos x)+C$. The wrong choice shows itself by a more complicated new integral (Oefening 609). ``Multiply by $1$'' (Oefening 610): $\int\log x\,dx=\int1\cdot\log x\,dx$.
\end{key}

\begin{bookex}{607}
Find $\int x\sin x\,dx$.
\solution
$f'=\sin x$ (so $f=-\cos x$), $g=x$ ($g'=1$): $\int x\sin x\,dx=-x\cos x+\int\cos x\,dx=-x\cos x+\sin x+C$. Check: $(-x\cos x+\sin x)'=-\cos x+x\sin x+\cos x=x\sin x$. \checkmark\qed
\end{bookex}

\begin{bookex}{610(c)}
Determine $\int\log x\,dx$ and $\int\arctan x\,dx$ by ``multiplying by $1$''.
\solution
$f'=1$, $f=x$, $g=\log x$: $\int\log x\,dx=x\log x-\int x\cdot\frac1x\,dx=x\log x-x+C$.\\
$f'=1$, $g=\arctan x$: $\int\arctan x\,dx=x\arctan x-\int\dfrac x{1+x^2}\,dx=x\arctan x-\tfrac12\log(1+x^2)+C$, since $\big(\log(1+x^2)\big)'=\frac{2x}{1+x^2}$.\qed
\end{bookex}

\begin{bookex}{611}
Compute $\int_1^ex^3\log^2x\,dx$.
\solution
Parts with $f'=x^3$, $g=\log^2x$ ($g'=\frac{2\log x}x$): $\Big[\dfrac{x^4}4\log^2x\Big]_1^e-\int_1^e\dfrac{x^4}4\cdot\dfrac{2\log x}x\,dx=\dfrac{e^4}4-\dfrac12\int_1^ex^3\log x\,dx$. Again by parts: $\int_1^ex^3\log x\,dx=\Big[\dfrac{x^4}4\log x\Big]_1^e-\int_1^e\dfrac{x^3}4\,dx=\dfrac{e^4}4-\dfrac{e^4-1}{16}=\dfrac{3e^4+1}{16}$. Hence the integral is $\dfrac{e^4}4-\dfrac{3e^4+1}{32}=\dfrac{5e^4-1}{32}$.\qed
\end{bookex}

\begin{bookex}{612}
The gamma function is $\Gamma(x)=\int_0^\infty t^{x-1}e^{-t}\,dt$ ($x>0$; assume convergence). (a) $\Gamma(1)$. (b) Show $\Gamma(x+1)=x\Gamma(x)$. (c) $\Gamma(6)$, and the pattern.
\solution
(a) $\Gamma(1)=\int_0^\infty e^{-t}dt=\lim_{b\to\infty}\big[-e^{-t}\big]_0^b=0+1=1$.\\
(b) Parts with $f'=e^{-t}$ ($f=-e^{-t}$), $g=t^x$: $\Gamma(x+1)=\big[-t^xe^{-t}\big]_0^\infty+x\int_0^\infty t^{x-1}e^{-t}dt$. The boundary term vanishes: at $t=0$ because $t^x\to0$ ($x>0$), at $\infty$ because $t^xe^{-t}\to0$ (Oefening 326d). So $\Gamma(x+1)=x\Gamma(x)$.\\
(c) $\Gamma(6)=5\Gamma(5)=5\cdot4\Gamma(4)=5\cdot4\cdot3\cdot2\cdot1\cdot\Gamma(1)=120=5!$; likewise $\Gamma(7)=720=6!$, and in general $\Gamma(n)=(n-1)!$: the gamma function extends the factorial to all positive reals.\qed
\end{bookex}

\begin{trybox}[{608, 609, 610(a),(b),(d)}]
\textbf{608} Check $\frac12e^x(\sin x-\cos x)$ by differentiating.\quad \textbf{609} Swap the roles of $f'$ and $g$ in Voorbeelden 6.5 and 6.6: which one goes wrong, and what is the signal?\quad \textbf{610} (a) $\int x\,dx$ and (b) $\int x^2dx$ by parts with $f'=1$; (d) try the same for $\int\frac1x\,dx$ and explain the paradox.
\end{trybox}

\hsection{Integration by substitution}
\begin{key}[{\S6.3 \textperiodcentered\ the chain rule backwards}]
$\displaystyle\int f(u(x))\,u'(x)\,dx=\int f(u)\,du$. Choose $u$, compute $du=u'(x)\,dx$, rewrite the integrand entirely in $u$, integrate, substitute back. Voorbeeld 6.7: $\int\frac{x}{\sqrt{1+x^2}}dx$, $u=1+x^2$, $du=2x\,dx$: $\int\frac1{\sqrt u}\cdot\frac12du=\sqrt u+C=\sqrt{1+x^2}+C$. Questions to ask: does the integrand remind you of a standard integral? which part would you like to get rid of? is one part the derivative of another? \textbf{Definite integrals:} either substitute back and use the original limits, or convert the limits ($u(a)$, $u(b)$) and never substitute back (Voorbeeld 6.8: $\int_0^8\frac{\cos\sqrt{1+x}}{\sqrt{1+x}}dx=2\int_1^3\cos u\,du=2(\sin3-\sin1)$).
\end{key}
\begin{strategy}[{\S6.3 \textperiodcentered\ trigonometric substitutions}]
Three variants of $\sin^2u+\cos^2u=1$: for $a^2-x^2$ put $x=a\sin u$ ($a^2-x^2=a^2\cos^2u$); for $x^2-a^2$ put $x=\dfrac a{\cos u}$ ($x^2-a^2=a^2\tan^2u$); for $a^2+x^2$ put $x=a\tan u$ ($a^2+x^2=\dfrac{a^2}{\cos^2u}$). Voorbeeld 6.9: $\int\sqrt{25-x^2}\,dx$ with $x=5\sin u$ becomes $25\int\cos^2u\,du=\frac{25}2\int(\cos2u+1)\,du=\frac{25}4\sin2u+\frac{25}2u$, and back: $\frac52x\sqrt{1-(\frac x5)^2}+\frac{25}2\arcsin\frac x5+C$. No absolute value is needed in $\sqrt{1-\sin^2u}=\cos u$ because $u\in[-\frac\pi2,\frac\pi2]$ (needed anyway for the substitution to be a bijection). Double-angle formulas are essential here.
\end{strategy}

\begin{bookex}{613}
With the given substitution: (a) $\int16x(3+4x^2)^7dx$, $u=3+4x^2$; (b) $\int\sin(2x)\sqrt{2+\sin^2x}\,dx$, $u=2+\sin^2x$.
\solution
(a) $du=8x\,dx$, so $16x\,dx=2\,du$: $\int2u^7du=\frac14u^8+C=\frac14(3+4x^2)^8+C$.\\
(b) $du=2\sin x\cos x\,dx=\sin(2x)\,dx$: $\int\sqrt u\,du=\frac23u^{3/2}+C=\frac23(2+\sin^2x)^{3/2}+C$.\qed
\end{bookex}

\begin{bookex}{615}
Evaluate $\int\dfrac{dx}{1+\sqrt{3x}}$ with $u=\sqrt{3x}$. Hint: be creative with $\frac{du}{dx}$.
\solution
$u^2=3x$, so $2u\,du=3\,dx$, $dx=\frac23u\,du$ (differentiate the squared relation rather than the root). Then
\begin{align*}
\int\frac{dx}{1+\sqrt{3x}}&=\int\frac{\frac23u}{1+u}\,du=\frac23\int\Big(1-\frac1{1+u}\Big)du\\
&=\frac23\big(u-\log|1+u|\big)+C=\frac23\Big(\sqrt{3x}-\log\big(1+\sqrt{3x}\big)\Big)+C.
\end{align*}\qed
\end{bookex}

\begin{bookex}{616}
Compute $\int_0^{\pi/2}\sqrt{1+\cos x}\,dx$.
\solution
Double-angle: $1+\cos x=2\cos^2\frac x2$, so $\sqrt{1+\cos x}=\sqrt2\,\big|\cos\frac x2\big|=\sqrt2\cos\frac x2$ on $[0,\frac\pi2]$. Then $\int_0^{\pi/2}\sqrt2\cos\frac x2\,dx=\sqrt2\Big[2\sin\frac x2\Big]_0^{\pi/2}=2\sqrt2\sin\frac\pi4=2\sqrt2\cdot\frac12\sqrt2=2$.\qed
\end{bookex}

\begin{bookex}{619}
Compute: (a) $\int\dfrac{x}{\sqrt{x^2-36}}\,dx$; (b) $\int\dfrac{dt}{t^2\sqrt{49-t^2}}$; (c) $\int_{1/2}^{\frac12\sqrt3}\dfrac{dx}{(1+4x^2)^2}$.
\solution
(a) No trig needed: $u=x^2-36$, $du=2x\,dx$: $\frac12\int u^{-1/2}du=\sqrt u=\sqrt{x^2-36}+C$.\\
(b) $t=7\sin u$, $dt=7\cos u\,du$, $\sqrt{49-t^2}=7\cos u$: $\int\dfrac{7\cos u\,du}{49\sin^2u\cdot7\cos u}=\dfrac1{49}\int\dfrac{du}{\sin^2u}=-\dfrac1{49}\cot u+C$. Back: $\cot u=\dfrac{\cos u}{\sin u}=\dfrac{\sqrt{1-(t/7)^2}}{t/7}=\dfrac{\sqrt{49-t^2}}t$, so the integral is $-\dfrac{\sqrt{49-t^2}}{49\,t}+C$.\\
(c) $2x=\tan u$, $dx=\frac12(1+\tan^2u)\,du=\frac{du}{2\cos^2u}$ and $1+4x^2=\frac1{\cos^2u}$: the integrand becomes $\cos^4u\cdot\frac1{2\cos^2u}du=\frac12\cos^2u\,du=\frac14(1+\cos2u)\,du$. Limits: $x=\frac12\Rightarrow\tan u=1\Rightarrow u=\frac\pi4$; $x=\frac12\sqrt3\Rightarrow\tan u=\sqrt3\Rightarrow u=\frac\pi3$. So the integral is $\frac14\big[u+\frac12\sin2u\big]_{\pi/4}^{\pi/3}=\frac14\Big(\frac\pi3+\frac{\sqrt3}4-\frac\pi4-\frac12\Big)=\dfrac\pi{48}+\dfrac{\sqrt3}{16}-\dfrac18\approx0.0487$.\qed
\end{bookex}

\begin{trybox}[{614, 617, 618}]
\textbf{614} By substitution: $\int2\sin(\pi x)dx$; $\int xe^{-x^2}dx$; $\int\frac{\log^2x}x dx$; $\int\tan x\,dx$; $\int(1+\cos x)^3\sin x\,dx$.\quad \textbf{617} $\int_e^{e^2}\frac{dx}{x\log x}$.\quad \textbf{618} Derive $\int\frac{dx}{\sqrt{1-x^2}}=\arcsin x$ with $x=\sin u$.
\end{trybox}

\hsection{Rational functions and partial fractions}
\begin{strategy}[{\S6.4 \textperiodcentered\ the standard procedure}]
\begin{enumerate}[leftmargin=1.5em,itemsep=1pt]
\item Degree of numerator $\ge$ degree of denominator: long division first (\S0.3). $\dfrac{(2x-1)^3+4}{2x^2-x}=4x-4+\dfrac{2x+3}{2x^2-x}$.
\item Denominator of degree $1$: $\int\dfrac c{ax+b}dx=\dfrac ca\log|ax+b|$.
\item Denominator of degree $2$, by the number of real zeros (discriminant): \textbf{none}: complete the square and reduce to $\int\frac{dx}{1+x^2}=\arctan x$ and $\int\frac{x\,dx}{1+x^2}=\frac12\log(1+x^2)$; \textbf{one (double)}: write the numerator in terms of the root, $\dfrac{3x+5}{(x+1)^2}=\dfrac3{x+1}+\dfrac2{(x+1)^2}$; \textbf{two}: partial fractions $\dfrac{ax+b}{(px+q)(rx+s)}=\dfrac A{px+q}+\dfrac B{rx+s}$, found by putting the right side over one denominator and comparing numerators (or substituting the roots).
\item Degree $\ge3$: factorise the denominator (easy cases only), one term per factor; a repeated factor $(x-c)^k$ or an irreducible quadratic gets a numerator of degree one less than its denominator: $\dfrac{Ax^2+Bx+C}{(1-x)^3}+\dfrac D{2+x}$, $\dfrac{Ax+B}{1+x^2}+\dfrac C{1+2x}$. Sometimes $u=x^2$ lowers the degree instead.
\end{enumerate}
Voorbeeld 6.10: $\dfrac{x+2}{x^3-x}=\dfrac{-2}x+\dfrac{3/2}{x-1}+\dfrac{1/2}{x+1}$. Voorbeeld 6.11: $\dfrac{4x^2+x+13}{(x-1)^2(x+5)}=\dfrac{x+2}{(x-1)^2}+\dfrac3{x+5}$ with $\dfrac{x+2}{(x-1)^2}=\dfrac1{x-1}+\dfrac3{(x-1)^2}$.
\end{strategy}

\begin{bookex}{621}
Evaluate $\int\dfrac{(2x-1)^3+4}{2x^2-x}\,dx$.
\solution
After the division (page 96): $\int(4x-4)\,dx+\int\dfrac{2x+3}{x(2x-1)}\,dx$. Partial fractions: $\dfrac{2x+3}{x(2x-1)}=\dfrac Ax+\dfrac B{2x-1}$ with $A(2x-1)+Bx=2x+3$; $x=0$ gives $A=-3$, $x=\frac12$ gives $\frac12B=4$, $B=8$ (this is the dictaat's example $\frac8{2x-1}-\frac3x$). Hence
\[
\int\frac{(2x-1)^3+4}{2x^2-x}\,dx=2x^2-4x-3\log|x|+4\log|2x-1|+C,
\]
using $\int\frac8{2x-1}dx=\frac82\log|2x-1|$. Check by differentiating: $4x-4-\frac3x+\frac8{2x-1}$ \checkmark.\qed
\end{bookex}

\begin{bookex}{624(a)--(d)}
Evaluate: (a) $\int\dfrac{dt}{1-t^2}$; (b) $\int\dfrac{x^3}{x^2+2x-15}dx$; (c) $\int\dfrac{dx}{x^2+6x+25}$; (d) $\int\dfrac{4y+3}{4y^2-9}dy$.
\solution
(a) $\dfrac1{1-t^2}=\dfrac1{(1-t)(1+t)}=\dfrac12\Big(\dfrac1{1-t}+\dfrac1{1+t}\Big)$, so the integral is $\dfrac12\big(-\log|1-t|+\log|1+t|\big)+C=\dfrac12\log\Big|\dfrac{1+t}{1-t}\Big|+C$.\\
(b) Divide: $x^3=(x^2+2x-15)(x-2)+(19x-30)$. Then $\dfrac{19x-30}{(x+5)(x-3)}=\dfrac A{x+5}+\dfrac B{x-3}$: $A(x-3)+B(x+5)=19x-30$; $x=3$: $8B=27$; $x=-5$: $-8A=-125$. So the integral is $\dfrac{x^2}2-2x+\dfrac{125}8\log|x+5|+\dfrac{27}8\log|x-3|+C$.\\
(c) No real zeros ($36-100<0$): complete the square, $x^2+6x+25=(x+3)^2+16=16\Big(1+\big(\tfrac{x+3}4\big)^2\Big)$; with $u=\frac{x+3}4$, $dx=4\,du$: $\int\dfrac{4\,du}{16(1+u^2)}=\dfrac14\arctan u=\dfrac14\arctan\dfrac{x+3}4+C$.\\
(d) $4y^2-9=(2y-3)(2y+3)$: $\dfrac{4y+3}{(2y-3)(2y+3)}=\dfrac A{2y-3}+\dfrac B{2y+3}$, $A(2y+3)+B(2y-3)=4y+3$; $y=\frac32$: $6A=9$, $A=\frac32$; $y=-\frac32$: $-6B=-3$, $B=\frac12$. Integral: $\dfrac32\cdot\dfrac12\log|2y-3|+\dfrac12\cdot\dfrac12\log|2y+3|=\dfrac34\log|2y-3|+\dfrac14\log|2y+3|+C$.\qed
\end{bookex}

\begin{bookex}{625}
Solve the logistic equation $y'=\frac1{50}y(100-y)$, $y(0)=25$, by separating variables.
\solution
$\dfrac{dy}{y(100-y)}=\dfrac{dt}{50}$. Partial fractions: $\dfrac1{y(100-y)}=\dfrac1{100}\Big(\dfrac1y+\dfrac1{100-y}\Big)$. Integrating: $\dfrac1{100}\big(\log y-\log(100-y)\big)=\dfrac t{50}+C$ (for $0<y<100$), i.e.\ $\log\dfrac y{100-y}=2t+C'$. At $t=0$: $\log\frac{25}{75}=\log\frac13=C'$. So $\dfrac y{100-y}=\dfrac13e^{2t}$, hence $3y=(100-y)e^{2t}$, $y(3+e^{2t})=100e^{2t}$ and
\[
y(t)=\frac{100e^{2t}}{3+e^{2t}}=\frac{100}{1+3e^{-2t}} .
\]
Check: $y(0)=25$; $y\to100$ as $t\to\infty$ (the carrying capacity $\beta=100$ of Oefening 524); and $y'=\frac{600e^{-2t}}{(1+3e^{-2t})^2}=\frac1{50}y(100-y)$ since $100-y=\frac{300e^{-2t}}{1+3e^{-2t}}$. \checkmark\qed
\end{bookex}

\begin{trybox}[{620, 622, 623, 624(e)--(h), 626}]
\textbf{620} Recombine $\frac13\big(\frac4{x-2}-\frac1{x+4}\big)$.\quad \textbf{622} Try partial fractions on $\frac{3x+5}{x^2+2x+1}$: what goes wrong?\quad \textbf{623} Finish Voorbeeld 6.11.\quad \textbf{624} (e) $\int\frac{4x^2+12x+9}{4x^2+12x+10}dx$; (f) $\int\frac{x\,dx}{x^4-x^2-6}$; (g) $\int\frac{dx}{(x^2-1)^2}$; (h) $\int\frac{2x^3-3x^2-2x-5}{x^4-1}dx$.\quad \textbf{626} $\int_0^1\frac{x^4(1-x)^4}{1+x^2}dx=\frac{22}7-\pi$, hence $\frac{22}7>\pi$.
\end{trybox}

\hsection{Mixed exercises}
\begin{strategy}[{\S6.5 \textperiodcentered\ choosing a method}]
Trig integrands: rewrite with the formulas ($\sin2x=2\sin x\cos x$, $\cos2x=2\cos^2x-1=1-2\sin^2x$, products to sums, $\cos^2+\sin^2=1$), then substitute ($u=\sin x$ when an odd power of $\cos$ is left over, and vice versa) or integrate by parts. One method is usually enough (628); sometimes two or three in a row (629). For $\int\frac{dx}{\cos x}$ multiply by $\frac{\cos x}{\cos x}$ and substitute $u=\sin x$ (632). If an approach is not working, try another.
\end{strategy}

\begin{bookex}{627(a),(c),(d)}
Evaluate (a) $\int\sin x\sin2x\,dx$; (c) $\int_0^{\pi/12}\sin^4\xi\,d\xi$; (d) $\int\tan^2\alpha\sin\alpha\,d\alpha$.
\solution
(a) $\sin2x=2\sin x\cos x$: $\int2\sin^2x\cos x\,dx$; $u=\sin x$: $\int2u^2du=\frac23\sin^3x+C$.\\
(c) Reduce the power with double angles: $\sin^4\xi=\Big(\dfrac{1-\cos2\xi}2\Big)^2=\dfrac14\big(1-2\cos2\xi+\cos^22\xi\big)=\dfrac14\Big(1-2\cos2\xi+\dfrac{1+\cos4\xi}2\Big)=\dfrac38-\dfrac12\cos2\xi+\dfrac18\cos4\xi$. Integrating: $\Big[\dfrac{3\xi}8-\dfrac14\sin2\xi+\dfrac1{32}\sin4\xi\Big]_0^{\pi/12}=\dfrac{\pi}{32}-\dfrac14\sin\dfrac\pi6+\dfrac1{32}\sin\dfrac\pi3=\dfrac\pi{32}-\dfrac18+\dfrac{\sqrt3}{64}$.\\
(d) $\tan^2\alpha\sin\alpha=\dfrac{\sin^3\alpha}{\cos^2\alpha}=\dfrac{(1-\cos^2\alpha)\sin\alpha}{\cos^2\alpha}$; $u=\cos\alpha$, $du=-\sin\alpha\,d\alpha$: $-\int\dfrac{1-u^2}{u^2}du=-\int(u^{-2}-1)\,du=\dfrac1u+u+C=\dfrac1{\cos\alpha}+\cos\alpha+C$.\qed
\end{bookex}

\begin{bookex}{628(c),(e)}
One method each: (c) $\int x^3e^{-x^2}dx$; (e) $\int\big(\tfrac12(t^2+t)^9+t(t^2+t)^9\big)dt$.
\solution
(c) $u=x^2$, $du=2x\,dx$: $\int x^2e^{-x^2}\cdot x\,dx=\frac12\int ue^{-u}du=\frac12\big(-ue^{-u}-e^{-u}\big)+C=-\frac12(x^2+1)e^{-x^2}+C$ (the inner integral by parts as in Voorbeeld 6.5).\\
(e) Factor: $(t^2+t)^9\big(t+\frac12\big)=\frac12(t^2+t)^9(2t+1)$, and $2t+1$ is the derivative of $t^2+t$: $u=t^2+t$ gives $\frac12\int u^9du=\frac1{20}(t^2+t)^{10}+C$.\qed
\end{bookex}

\begin{bookex}{629(b),(c),(f),(g)}
(b) $\int\dfrac{\sqrt{x+1}}{x+2}dx$; (c) $\int x^3\sqrt{1-x^2}\,dx$; (f) $\int\dfrac{x\arctan x}{(1+x^2)^{3/2}}dx$; (g) $\int\dfrac{x^2}{\sqrt{2-x^2}}dx$.
\solution
(b) $u=\sqrt{x+1}$, $x=u^2-1$, $dx=2u\,du$: $\int\dfrac{u\cdot2u}{u^2+1}du=2\int\Big(1-\dfrac1{u^2+1}\Big)du=2u-2\arctan u+C=2\sqrt{x+1}-2\arctan\sqrt{x+1}+C$.\\
(c) $u=1-x^2$, $x^2=1-u$, $du=-2x\,dx$: $\int x^2\sqrt{1-x^2}\cdot x\,dx=-\dfrac12\int(1-u)\sqrt u\,du=-\dfrac12\Big(\dfrac23u^{3/2}-\dfrac25u^{5/2}\Big)+C=-\dfrac13(1-x^2)^{3/2}+\dfrac15(1-x^2)^{5/2}+C$.\\
(f) Parts with $f'=\dfrac{x}{(1+x^2)^{3/2}}$, whose primitive is $f=-\dfrac1{\sqrt{1+x^2}}$ (substitution $u=1+x^2$), and $g=\arctan x$: $-\dfrac{\arctan x}{\sqrt{1+x^2}}+\int\dfrac{dx}{(1+x^2)^{3/2}}$. The last integral: $x=\tan u$, $dx=\frac{du}{\cos^2u}$, $(1+x^2)^{3/2}=\cos^{-3}u$, so it is $\int\cos u\,du=\sin u=\dfrac x{\sqrt{1+x^2}}$. Result: $\dfrac{x-\arctan x}{\sqrt{1+x^2}}+C$.\\
(g) $x=\sqrt2\sin u$, $dx=\sqrt2\cos u\,du$, $\sqrt{2-x^2}=\sqrt2\cos u$: $\int\dfrac{2\sin^2u\cdot\sqrt2\cos u}{\sqrt2\cos u}du=2\int\sin^2u\,du=u-\sin u\cos u+C$. Back: $u=\arcsin\frac x{\sqrt2}$, $\sin u\cos u=\frac x{\sqrt2}\sqrt{1-\frac{x^2}2}=\frac12x\sqrt{2-x^2}$. So $\arcsin\dfrac x{\sqrt2}-\dfrac12x\sqrt{2-x^2}+C$.\qed
\end{bookex}

\begin{bookex}{630(a),(e),(f)}
Solve the initial value problems: (a) $xy'-(x+1)y=x^2-x^3$, $y(1)=2$; (e) $y'+y\cos x=2xe^{-\sin x}$, $y(\pi)=0$; (f) $y'=\dfrac{4-2x}{3y^2-5}$, $y(1)=3$.
\solution
(a) Divide by $x$: $y'-\big(1+\frac1x\big)y=x-x^2$. Integrating factor $u=e^{-x-\log x}=\dfrac{e^{-x}}x$: $\Big(\dfrac{e^{-x}}xy\Big)'=(1-x)e^{-x}$, and $\int(1-x)e^{-x}dx=xe^{-x}+C$ (check: $(xe^{-x})'=e^{-x}-xe^{-x}$). So $\dfrac{e^{-x}}xy=xe^{-x}+C$, $y=x^2+Cxe^x$. $y(1)=1+Ce=2$ gives $C=\frac1e$: $y=x^2+xe^{x-1}$.\\
(e) Integrating factor $e^{\sin x}$: $\big(e^{\sin x}y\big)'=2x$, so $e^{\sin x}y=x^2+C$, $y=(x^2+C)e^{-\sin x}$. $y(\pi)=(\pi^2+C)e^0=0$ gives $C=-\pi^2$: $y=(x^2-\pi^2)e^{-\sin x}$.\\
(f) Not linear; separate: $(3y^2-5)\,dy=(4-2x)\,dx$, so $y^3-5y=4x-x^2+C$. At $(1,3)$: $27-15=4-1+C$, $C=9$. The solution is given implicitly by $y^3-5y=4x-x^2+9$ (it cannot be solved for $y$ in closed form; the d.v.\ holds by implicit differentiation).\qed
\end{bookex}

\begin{bookex}{632}
$\int\dfrac{dx}{\cos x}$: (a) multiply numerator and denominator by a suitable factor so that $u=\sin x$ works; (b) integrate by partial fractions; (c) use the log rules and back-substitution to reach $\log\big|\frac1{\cos x}+\tan x\big|$.
\solution
(a) Multiply by $\cos x$: $\dfrac1{\cos x}=\dfrac{\cos x}{\cos^2x}=\dfrac{\cos x}{1-\sin^2x}$. With $u=\sin x$, $du=\cos x\,dx$: $\int\dfrac{du}{1-u^2}$.\\
(b) $\dfrac1{1-u^2}=\dfrac12\Big(\dfrac1{1-u}+\dfrac1{1+u}\Big)$ (624a), so $\int\dfrac{du}{1-u^2}=\dfrac12\log\Big|\dfrac{1+u}{1-u}\Big|+C$.\\
(c) $\dfrac{1+u}{1-u}=\dfrac{(1+u)^2}{1-u^2}=\dfrac{(1+\sin x)^2}{\cos^2x}$, so $\dfrac12\log\dfrac{(1+\sin x)^2}{\cos^2x}=\log\Big|\dfrac{1+\sin x}{\cos x}\Big|=\log\Big|\dfrac1{\cos x}+\tan x\Big|+C$.\qed
\end{bookex}

\begin{bookex}{633}
First a substitution, then Oefening 632: (a) $\int\dfrac{dx}{\sqrt{x^2-1}}$; (b) $\int_0^2\dfrac{dx}{\sqrt{4+x^2}}$.
\solution
(a) $x=\dfrac1{\cos u}$, $dx=\dfrac{\sin u}{\cos^2u}du$, $\sqrt{x^2-1}=\tan u$ (for $0\le u<\frac\pi2$): $\int\dfrac{\sin u/\cos^2u}{\sin u/\cos u}du=\int\dfrac{du}{\cos u}=\log\Big|\dfrac1{\cos u}+\tan u\Big|=\log\big|x+\sqrt{x^2-1}\big|+C$ (this is $\operatorname{arcosh}x$, Oefening 714).\\
(b) $x=2\tan u$, $dx=\dfrac{2\,du}{\cos^2u}$, $\sqrt{4+x^2}=\dfrac2{\cos u}$: the integrand becomes $\dfrac{du}{\cos u}$; limits $u=0$ to $u=\frac\pi4$: $\Big[\log\Big(\dfrac1{\cos u}+\tan u\Big)\Big]_0^{\pi/4}=\log(\sqrt2+1)-\log1=\log(1+\sqrt2)$.\qed
\end{bookex}

\begin{trybox}[{627(b),(e),(f), 628(a),(b),(d),(f), 629(a),(d),(e),(h), 630(b),(c),(d),(g), 631}]
\textbf{627} (b) $\int\cos^5\psi\,d\psi$; (e) $\int\tan(\arcsin t)\,dt$; (f) $\int\sin(\tau+1)\sin\tau\,d\tau$.\quad \textbf{628} (a) $\frac x{\sqrt{x^2+a^2}}$; (b) $(x^2+3x+2)^{-1}$; (d) $\frac{x^3+1}{x\sqrt x}$; (f) $x^2\cos x$.\quad \textbf{629} (a) $\log(1+1/x^2)$; (d) $\frac{x^3-8}{x^2-8}$; (e) $e^{\sqrt{1+x}}$; (h) $\sin(\log x)$.\quad \textbf{630} (b) $y'=e^{\sqrt x}/y$, $y(0)=-2$; (c) $y'+10y=1$, $y(\frac1{10})=\frac2{10}$; (d) $y'+3x^2y=x^2$, $y(0)=1$; (g) $y'=2xy^3(2x^2+1)$, $y(1)=1$.\quad \textbf{631} $\int_0^2\sqrt{(x-1)^2}\,dx$; $\int_0^\infty\frac{dx}{e^x+1}$.
\end{trybox}

\hsection{(Extra) The natural logarithm and the exponential function}
\begin{key}[{Definitions 6.12--6.18 \textperiodcentered\ building $\log$ and $\exp$ from an integral}]
Define $\log x=\int_1^x\frac{dt}t$ for $x>0$. Then $\log'x=\frac1x$ (fundamental theorem) and $\log(xy)=\log x+\log y$ (Stelling 6.13: $x\mapsto\log(xy)$ has derivative $\frac1x$, so it is $\log x+c$, and $x=1$ gives $c=\log y$): it deserves the name logarithm. Define $\exp=\log_{\mathrm{inv}}:\R\to\R_{>0}$; then $\exp'=\exp$ (differentiate $\log(\exp x)=x$), $\exp(x+y)=\exp x\exp y$, $e=\exp1$ (so $\int_1^e\frac{dt}t=1$), and finally $e^x=\exp x$ and $a^x=e^{x\log a}$ for all real $x$, which at last gives $3^{\sqrt2}$ a meaning. $e=\sum_{k\ge0}\frac1{k!}$.
\end{key}

\begin{trybox}[{634 (Extra), 635 (Extra)}]
\textbf{634} How many terms of $\sum\frac1{k!}$ give the first $10$ decimals of $e$?\quad \textbf{635} From Definition 6.18: $a^1=a$, $a^{x+y}=a^xa^y$, $(a^x)^y=a^{xy}$.
\end{trybox}

\begin{warn}
\begin{itemize}
\item Check every primitive by differentiating; it takes ten seconds and catches most errors.
\item In a substitution, convert $dx$ as well (differentiate the relation between $u$ and $x$; for $u=\sqrt{\cdot}$ square first). With definite integrals, either convert the limits or substitute back, never half of each.
\item Before partial fractions: divide if the numerator's degree is too high, and check the discriminant of a quadratic denominator.
\item Improper integrals: one limit per ``improperness''; a symmetric cancellation ($\int_{-a}^a\sin x\,dx=0$) proves nothing (606).
\item Do not forget $+C$, and do not add one to a definite integral.
\end{itemize}
\end{warn}

\chapter{More exercises}
\chaptercontents{A & Mixed exercises (701--724) \\ B & (Extra) Function investigation (\S7.1)}%
{Exercises 701--730 (all worked or solved).\\ The dictaat: ``use them as extra practice, exam preparation, or for fun''. Each one combines techniques from several chapters; the first question to ask is \emph{which} chapter it belongs to.}

\hsection{Mixed exercises on Chapters 0--6}

\begin{bookex}{701}
$f$ differentiable with inverse $f_{\mathrm{inv}}$, so $f(f_{\mathrm{inv}}(x))=x$. (a) Differentiate both sides and show $f_{\mathrm{inv}}'(x)=\dfrac1{f'(f_{\mathrm{inv}}(x))}$. (b) For $f(x)=x^3+x$ find $f_{\mathrm{inv}}'(-10)$.
\solution
(a) Chain rule: $f'(f_{\mathrm{inv}}(x))\cdot f_{\mathrm{inv}}'(x)=1$, so $f_{\mathrm{inv}}'(x)=\dfrac1{f'(f_{\mathrm{inv}}(x))}$ (where $f'\ne0$).\\
(b) $f_{\mathrm{inv}}(-10)$ is the $x$ with $x^3+x=-10$: $x=-2$ ($-8-2=-10$; $f$ is increasing, so it is the only one). $f'(x)=3x^2+1$, $f'(-2)=13$, so $f_{\mathrm{inv}}'(-10)=\dfrac1{13}$. (No formula for $f_{\mathrm{inv}}$ is needed.)\qed
\end{bookex}

\begin{bookex}{702}
Show with sum and double-angle formulas that $\tan\big(\frac\pi4+\frac\varphi2\big)=\dfrac1{\cos\varphi}+\tan\varphi$.
\solution
Right side: $\dfrac{1+\sin\varphi}{\cos\varphi}$. With half angles, $1+\sin\varphi=\cos^2\frac\varphi2+\sin^2\frac\varphi2+2\sin\frac\varphi2\cos\frac\varphi2=\big(\cos\frac\varphi2+\sin\frac\varphi2\big)^2$ and $\cos\varphi=\cos^2\frac\varphi2-\sin^2\frac\varphi2=\big(\cos\frac\varphi2-\sin\frac\varphi2\big)\big(\cos\frac\varphi2+\sin\frac\varphi2\big)$. Hence
\begin{align*}
\frac{1+\sin\varphi}{\cos\varphi}&=\frac{\cos\frac\varphi2+\sin\frac\varphi2}{\cos\frac\varphi2-\sin\frac\varphi2}=\frac{1+\tan\frac\varphi2}{1-\tan\frac\varphi2}\\
&=\frac{\tan\frac\pi4+\tan\frac\varphi2}{1-\tan\frac\pi4\tan\frac\varphi2}=\tan\Big(\frac\pi4+\frac\varphi2\Big)
\end{align*}
by the sum formula for $\tan$ (Oefening 26c) with $\tan\frac\pi4=1$.\qed
\end{bookex}

\begin{bookex}{705}
Second derivative of $\sqrt{\dfrac{1+x}{1-x}}$, simplified.
\solution
Logarithmic differentiation: $\log y=\frac12\big(\log(1+x)-\log(1-x)\big)$, so $\dfrac{y'}y=\dfrac12\Big(\dfrac1{1+x}+\dfrac1{1-x}\Big)=\dfrac1{1-x^2}$ and $y'=\dfrac{(1+x)^{1/2}}{(1-x)^{1/2}(1-x)(1+x)}=(1+x)^{-1/2}(1-x)^{-3/2}$. Then
\begin{align*}
y''&=-\tfrac12(1+x)^{-3/2}(1-x)^{-3/2}+\tfrac32(1+x)^{-1/2}(1-x)^{-5/2}\\
&=(1+x)^{-3/2}(1-x)^{-5/2}\big(-\tfrac12(1-x)+\tfrac32(1+x)\big)=\frac{1+2x}{(1+x)^{3/2}(1-x)^{5/2}}.
\end{align*}\qed
\end{bookex}

\begin{bookex}{706}
Find all $a,b$ for which $f(x)=ax^2+b$ ($x<1$), $f(x)=\dfrac1{a+bx}$ ($x\ge1$) is differentiable.
\solution
Away from $x=1$ both pieces are differentiable (provided $a+bx\ne0$ for $x\ge1$). At $x=1$ we need continuity and equal one-sided derivatives: $a+b=\dfrac1{a+b}$, so $(a+b)^2=1$; and $2a=\dfrac{-b}{(a+b)^2}=-b$. So $b=-2a$ and $a+b=-a=\pm1$: $(a,b)=(-1,2)$ or $(1,-2)$. Check $(1,-2)$: $f(1^-)=1-2=-1$, $f(1)=\frac1{1-2}=-1$; slopes $2a=2$ and $\frac{-b}{(a+b)^2}=2$; and $a+bx=1-2x\ne0$ for $x\ge1$. Similarly for $(-1,2)$ ($-1+2x\ne0$ for $x\ge1$). Two solutions.\qed
\end{bookex}

\begin{bookex}{710}
$f(x)=x\arctan(2x^3)$. Find the coefficient $c_{10}$ of $x^{10}$ in its Taylor approximation at $0$.
\solution
$\arctan u=u-\frac13u^3+\frac15u^5-\dots$ (Oefening 709). With $u=2x^3$: $\arctan(2x^3)=2x^3-\frac83x^9+\frac{32}5x^{15}-\dots$, so $f(x)=2x^4-\frac83x^{10}+\dots$ and $c_{10}=-\dfrac83$. (The next term is $x^{16}$.)\qed
\end{bookex}

\begin{bookex}{713}
Determine $\lim_{x\to\frac12}\dfrac{\arcsin x-\frac\pi6}{x-\frac12}$.
\solution
Since $\arcsin\frac12=\frac\pi6$, this is $\dfrac{\arcsin x-\arcsin\frac12}{x-\frac12}$ with $x\to\frac12$: the derivative of $\arcsin$ at $\frac12$ (\S2.3.1). $\arcsin'x=\dfrac1{\sqrt{1-x^2}}$, so the limit is $\dfrac1{\sqrt{1-\frac14}}=\dfrac2{\sqrt3}=\dfrac23\sqrt3$. (l'H\^opital gives the same in one step.)\qed
\end{bookex}

\begin{bookex}{718}
Solve $z'(t)+z(t)=\sin t$ with $z(\frac\pi4)=1$.
\solution
$z_h=Ce^{-t}$. Guess $z_p=A\sin t+B\cos t$: $z_p'+z_p=(A\cos t-B\sin t)+(A\sin t+B\cos t)=(A-B)\sin t+(A+B)\cos t=\sin t$ gives $A-B=1$, $A+B=0$: $A=\frac12$, $B=-\frac12$. General solution $z=Ce^{-t}+\frac12(\sin t-\cos t)$. At $t=\frac\pi4$: $\sin\frac\pi4=\cos\frac\pi4$, so $z(\frac\pi4)=Ce^{-\pi/4}=1$, $C=e^{\pi/4}$. Hence $z(t)=e^{\pi/4-t}+\frac12(\sin t-\cos t)$.\qed
\end{bookex}

\begin{bookex}{719}
Primitives of: (a) $\arcsin(2x)$; (b) $\dfrac{2x+1}{x^4+2x^3+x^2}$; (c) $x\cos(3x)$; (d) $\dfrac{xe^{2x}}{(1+e^{2x})^2}$.
\solution
(a) Parts with $f'=1$: $x\arcsin(2x)-\int\dfrac{2x}{\sqrt{1-4x^2}}dx$; with $u=1-4x^2$, $du=-8x\,dx$ the last integral is $-\frac14\int u^{-1/2}du=-\frac12\sqrt{1-4x^2}$. So $x\arcsin(2x)+\frac12\sqrt{1-4x^2}+C$.\\
(b) $x^4+2x^3+x^2=(x^2+x)^2$ and $2x+1=(x^2+x)'$: $u=x^2+x$ gives $\int\dfrac{du}{u^2}=-\dfrac1u=-\dfrac1{x^2+x}+C$.\\
(c) Parts with $f'=\cos3x$, $f=\frac13\sin3x$: $\frac13x\sin3x-\frac13\int\sin3x\,dx=\frac13x\sin3x+\frac19\cos3x+C$.\\
(d) Parts with $f'=\dfrac{e^{2x}}{(1+e^{2x})^2}$, whose primitive is $f=-\dfrac1{2(1+e^{2x})}$ ($u=1+e^{2x}$), and $g=x$: $-\dfrac x{2(1+e^{2x})}+\dfrac12\int\dfrac{dx}{1+e^{2x}}$. And $\int\dfrac{dx}{1+e^{2x}}=\int\dfrac{e^{-2x}}{e^{-2x}+1}dx=-\dfrac12\log(1+e^{-2x})$ (as in 631b). Result: $-\dfrac x{2(1+e^{2x})}-\dfrac14\log\big(1+e^{-2x}\big)+C$, or equivalently $\dfrac x2-\dfrac x{2(1+e^{2x})}-\dfrac14\log(1+e^{2x})+C$.\qed
\end{bookex}

\begin{bookex}{720}
Determine (or show divergence): (a) $\int_0^\pi\cos t\sqrt{1+\cos^2t}\,dt$; (b) $\int_0^1\dfrac{\log(1+x)}{\sqrt x}dx$.
\solution
(a) Write $\cos^2t=1-\sin^2t$ and substitute $u=\sin t$, $du=\cos t\,dt$: the limits become $u(0)=0$ and $u(\pi)=0$, so the integral is $\int_0^0\sqrt{2-u^2}\,du=0$. (Symmetry: the integrand at $\pi-t$ is the negative of that at $t$.)\\
(b) Improper at $0$ (the integrand $\sim\frac x{\sqrt x}=\sqrt x$ there, so it is actually bounded: convergent). $x=u^2$, $dx=2u\,du$: $\int_0^12\log(1+u^2)\,du$. Parts: $\big[2u\log(1+u^2)\big]_0^1-\int_0^1\dfrac{4u^2}{1+u^2}du=2\log2-4\int_0^1\Big(1-\dfrac1{1+u^2}\Big)du=2\log2-4+\pi$.\qed
\end{bookex}

\begin{bookex}{722}
Reaction kinetics: $\dfrac{dx}{dt}=k(a-x)(b-x)$, $x(0)=0$, $k>0$. Solve in the form $kt=\dots$ for (a) $a=b$; (b) $a\ne b$.
\solution
(a) $\dfrac{dx}{(a-x)^2}=k\,dt$ gives $\dfrac1{a-x}=kt+C$; $x(0)=0$: $C=\frac1a$. So $kt=\dfrac1{a-x}-\dfrac1a=\dfrac x{a(a-x)}$.\\
(b) Partial fractions: $\dfrac1{(a-x)(b-x)}=\dfrac1{b-a}\Big(\dfrac1{a-x}-\dfrac1{b-x}\Big)$. Integrating (for $x<a,b$): $\dfrac1{b-a}\big(-\log(a-x)+\log(b-x)\big)=kt+C$, and $x(0)=0$ gives $C=\dfrac1{b-a}\log\dfrac ba$. Hence
\[
kt=\frac1{b-a}\Big(\log\frac{b-x}{a-x}-\log\frac ba\Big)=\frac1{b-a}\log\frac{a(b-x)}{b(a-x)} .
\]\qed
\end{bookex}

\begin{bookex}{723}
Determine $\int\tan(1-x)\tan(1+x)\,dx$. Hint: Oefening 26(c), creatively.
\solution
From $\tan(A+B)=\dfrac{\tan A+\tan B}{1-\tan A\tan B}$ we get $\tan A\tan B=1-\dfrac{\tan A+\tan B}{\tan(A+B)}$. With $A=1-x$, $B=1+x$, $A+B=2$:
\[
\tan(1-x)\tan(1+x)=1-\frac{\tan(1-x)+\tan(1+x)}{\tan2}.
\]
Now $\int\tan(1-x)\,dx=\log|\cos(1-x)|$ (chain rule, inner derivative $-1$) and $\int\tan(1+x)\,dx=-\log|\cos(1+x)|$. So
\[
\int\tan(1-x)\tan(1+x)\,dx=x-\frac1{\tan2}\Big(\log|\cos(1-x)|-\log|\cos(1+x)|\Big)+C .
\]\qed
\end{bookex}

\begin{bookex}{724}
Euler's beta function $B_{p,q}=\int_0^1x^{p-1}(1-x)^{q-1}dx$ ($p,q\ge1$ integers). Prove: (a) $B_{p,q}=B_{q,p}$; (b) $B_{p+1,q}+B_{p,q+1}=B_{p,q}$; (c) $B_{p+1,q}=\frac pqB_{p,q+1}$; (d) $B_{p+1,q}=\frac p{p+q}B_{p,q}$; (e) a closed formula.
\solution
(a) Substitute $x=1-u$ ($dx=-du$, limits swap): $\int_0^1(1-u)^{p-1}u^{q-1}du=B_{q,p}$.\\
(b) $x^p(1-x)^{q-1}+x^{p-1}(1-x)^q=x^{p-1}(1-x)^{q-1}\big(x+(1-x)\big)=x^{p-1}(1-x)^{q-1}$; integrate.\\
(c) Parts with $f'=(1-x)^{q-1}$, $f=-\frac1q(1-x)^q$, $g=x^p$: $B_{p+1,q}=\Big[-\frac1qx^p(1-x)^q\Big]_0^1+\frac pq\int_0^1x^{p-1}(1-x)^qdx=0+\frac pqB_{p,q+1}$.\\
(d) By (c), $B_{p,q+1}=\frac qpB_{p+1,q}$; substitute in (b): $B_{p+1,q}\big(1+\frac qp\big)=B_{p,q}$, so $B_{p+1,q}=\frac p{p+q}B_{p,q}$.\\
(e) $B_{1,q}=\int_0^1(1-x)^{q-1}dx=\frac1q$. Applying (d) repeatedly: $B_{p,q}=\frac{p-1}{p+q-1}\cdot\frac{p-2}{p+q-2}\cdots\frac1{q+1}\cdot B_{1,q}=\dfrac{(p-1)!\,(q-1)!}{(p+q-1)!}$. (Check: $B_{3,4}=\frac{2!\,3!}{6!}=\frac1{60}$.)\qed
\end{bookex}

\begin{trybox}[{703, 704, 707, 708, 709, 711, 712, 714, 715, 716, 717, 721}]
\textbf{703} $\log(49+\sqrt{49^2-1})+\log(49-\sqrt{49^2-1})$.\ \textbf{704} $f(x)=\log\frac{e^x+1}{e^x-1}$: $f\circ f(x)=x$.\ \textbf{707} $\sin$ and $\arctan$ share a tangent line at $0$.\ \textbf{708} $P_4$ of $\frac{1+x^2}{1-x}$.\ \textbf{709} $P_5$ of $\arctan x$.\ \textbf{711} Build $\frac{-3}{x+2}$ and $\frac{1-x}{2+x}$ by composing $i(x)=\frac1x$, $p(x)=1+x$, $m(x)=-x$, $t(x)=3x$.\ \textbf{712} $\sin(\arctan x)$.\ \textbf{714--716} $\operatorname{arcosh}x=\log(x\pm\sqrt{x^2-1})$: derivation, the two branches as mirror images, the derivative in two ways.\ \textbf{717} $\tanh$ and $\operatorname{artanh}x=\frac12\log\frac{1+x}{1-x}$.\ \textbf{721} $\int_0^1\frac{1-x^n}{1-x}dx=\sum_{k=1}^n\frac1k$.
\end{trybox}

\hsection{(Extra) Function investigation}
\begin{strategy}[{\S7.1 \textperiodcentered\ the six steps and the sign table}]
(1) Domain. (2) Symmetry (even/odd) and a more convenient form of the formula. (3) Intercepts: $f(x)=0$ and $f(0)$. (4) Behaviour at the edges of the domain, including $\pm\infty$: limits, asymptotes (not every edge gives an asymptote). (5) $f'$: horizontal tangents where $f'=0$, corners where $f'$ does not exist; list the critical points. (6) $f''$: convex ($f''>0$) or concave, inflection points where $f''$ changes sign. Collect everything in a \emph{tekenschema} (sign table) for $f,f',f''$, check for contradictions (an alarm signal means an error somewhere), then sketch, smuggling in no features you did not find. Model (dictaat): $f(x)=\dfrac{2-x}{2+x}=\dfrac4{2+x}-1$: domain $x\ne-2$; zero $x=2$, $f(0)=1$; asymptotes $y=-1$ and $x=-2$; $f'=-\frac4{(2+x)^2}<0$ everywhere; $f''=\frac8{(2+x)^3}$: concave left of $-2$, convex right of it, no inflection point.
\end{strategy}

\begin{bookex}{726(d)}
Investigate $f:x\mapsto\dfrac{(x+1)^3}{x^2}$ and sketch the graph.
\solution
(1) Domain $x\ne0$. (2) Division gives $f(x)=x+3+\dfrac3x+\dfrac1{x^2}$. (3) Zero: $x=-1$ (triple zero: the graph crosses the axis with a horizontal tangent); no $y$-intercept. (4) $x\to0^\pm$: numerator $\to1$, denominator $\to0^+$: $f\to+\infty$ on both sides, vertical asymptote $x=0$. $x\to\pm\infty$: $f(x)-(x+3)=\frac3x+\frac1{x^2}\to0$: oblique asymptote $y=x+3$, approached from above for $x>0$ and from below for $x<0$ (sign of $\frac3x$). (5) $f'(x)=\dfrac{(x+1)^2(x-2)}{x^3}$: zero at $x=-1$ and $x=2$. Sign: for $x<0$, $(x-2)/x^3>0$: increasing on $(-\infty,0)$ (with a flat point at $-1$); on $(0,2)$ negative: decreasing; on $(2,\infty)$ positive: increasing. Local minimum $f(2)=\frac{27}4$. (6) $f''(x)=\dfrac{6(x+1)}{x^4}$: negative for $x<-1$ (concave), positive for $x>-1$, $x\ne0$ (convex); inflection point $(-1,0)$.\\
Sign table (values at $-\infty,-1,0^-,0^+,2,+\infty$): $f$: $-\infty,\ 0,\ +\infty\,|\,+\infty,\ \frac{27}4,\ +\infty$; $f'$: $+,\ 0,\ +\,|\,-,\ 0,\ +$; $f''$: $-,\ 0,\ +\,|\,+,\ +,\ +$.\\
Sketch: from the lower left along $y=x+3$ (below it), concave, rising through $(-1,0)$ with a horizontal tangent, then convex and shooting up to $+\infty$ at $x=0^-$; on the right, coming down from $+\infty$ at $0^+$ to the minimum $(2,\frac{27}4)$, then rising along $y=x+3$ (above it).\qed
\end{bookex}

\begin{bookex}{727}
Four students investigated the same differentiable $f$ on $(0,\infty)$ with the single zero $f(3)=0$, $\lim_{x\to\infty}f(x)=-2$ and exactly one inflection point; they claim (a) $\lim_{0^+}f=-\infty$, $f'(3)=0$, inflection at $6$; (b) $\lim_{0^+}f=-4$, $f'(3)>0$, inflection at $1$; (c) $\lim_{0^+}f=+4$, $f'(3)<0$, inflection at $3$; (d) $\lim_{0^+}f=+\infty$, $f'(3)<0$, inflection at $1$. Who can draw a consistent graph?
\solution
\emph{(b) impossible:} $f'(3)>0$ means $f$ goes from negative to positive at $3$; then $f>0$ just after $3$, but $f\to-2<0$, so $f$ would have to cross $0$ again, a second zero.\\
\emph{(d) impossible:} $f$ decreases from $+\infty$ through $(3,0)$ to $-2$. On $(1,\infty)$ (after the only inflection point) $f$ would be either concave or convex. Concave and decreasing with $f'<0$ forces $f\to-\infty$ (the slope keeps getting more negative), contradicting the limit $-2$. So $f$ is convex on $(1,\infty)$ and concave on $(0,1)$; but a concave function on $(0,1)$ stays below its tangent lines, hence is bounded above near $0$, contradicting $f\to+\infty$. \\
\emph{(a) consistent:} $f$ rises from $-\infty$ to a maximum $f(3)=0$ (touching the axis, $f'(3)=0$), concave there, then decreases to $-2$, which requires it to become convex eventually: one inflection point at some $x>3$, e.g.\ $6$. $f<0$ except at $3$: one zero. \checkmark\\
\emph{(c) consistent:} $f$ decreases from $4$ through $(3,0)$ to $-2$ like a falling S-curve: concave on $(0,3)$, convex on $(3,\infty)$, inflection at $3$; the value at the inflection need not be the midpoint of $4$ and $-2$. \checkmark\qed
\end{bookex}

\begin{bookex}{728}
Find all $x$ with $\dfrac1{1-x^2}>\dfrac1x$. Check carefully.
\solution
Domain: $x\ne0,\pm1$. Never multiply by $x$ or $1-x^2$ (unknown signs); bring everything to one side:
\[
\frac1{1-x^2}-\frac1x=\frac{x-(1-x^2)}{x(1-x^2)}=\frac{x^2+x-1}{x(1-x^2)}>0 .
\]
Zeros of the numerator: $x=\frac{-1\pm\sqrt5}2$, i.e.\ $\alpha\approx-1.618$ and $\beta\approx0.618$. Sign table with the points $\alpha,-1,0,\beta,1$: numerator $+$ on $(-\infty,\alpha)$ and $(\beta,\infty)$, $-$ between; denominator $x(1-x)(1+x)$: $+$ on $(-\infty,-1)$, $-$ on $(-1,0)$, $+$ on $(0,1)$, $-$ on $(1,\infty)$. Quotient positive on $(-\infty,\alpha)$, $(-1,0)$ and $(\beta,1)$. Answer: $x<\dfrac{-1-\sqrt5}2$ or $-1<x<0$ or $\dfrac{\sqrt5-1}2<x<1$. Check $x=-2$: $\frac1{-3}=-0.33>-0.5$ \checkmark; $x=-\frac12$: $\frac43>-2$ \checkmark; $x=\frac12$: $\frac43>2$ \texttimes; $x=0.8$: $\frac1{0.36}=2.8>1.25$ \checkmark; $x=2$: $-\frac13>\frac12$ \texttimes.\qed
\end{bookex}

\begin{bookex}{729}
Prove $\log x\le x-1$ for all $x>0$.
\solution
Let $g(x)=x-1-\log x$ on $(0,\infty)$. $g(1)=0$ and $g'(x)=1-\frac1x$, negative on $(0,1)$ and positive on $(1,\infty)$: $g$ decreases to its minimum $g(1)=0$ and increases afterwards. So $g(x)\ge0$ with equality only at $x=1$: the graph of $g$ touches the $x$-axis at $1$ without crossing it, i.e.\ $\log x\le x-1$.\qed
\end{bookex}

\begin{trybox}[{725, 726(a),(b),(c),(e)--(j), 730}]
\textbf{725} Investigate $\sqrt x$, $e^x$, $\log x$, $\sin x$, $\cos x$, $\tan x$.\quad \textbf{726} (a) $\sqrt{9-x^2}$; (b) $\frac x{\sqrt{1-x^2}}$; (c) $e^{1-x^2}$; (e) $\log\frac{2x}{(x-4)^2}$; (f) $(1+\frac1x)e^{-x}$; (g) $e^{-1/x}$; (h) $\log(x-\frac1x)$; (i) $e^{1-\frac1{1+x}}$; (j) $\arcsin(1-\sqrt x)$.\quad \textbf{730} Prove $\log x\le2(\sqrt x-1)$.
\end{trybox}

\renewcommand{\chapterprefix}{Appendix}\renewcommand{\thechapter}{S}
\chapter{Solutions to the NOW YOU exercises}
\chaptercontents{A & Chapter 0 \\ B & Chapter 1 \\ C & Chapter 2 \\ D & Chapter 3 \\ E & Chapter 4 \\ F & Chapter 5 \\ G & Chapter 6 \\ H & Chapter 7}{Every Oefening of Hoofdstuk 0--7 that is not worked in the main text is solved here, in the order of the dictaat. Compare line by line, and reflect: could you have seen the answer coming? Would another approach work?}
\hsection{Chapter 0}
\sol{1} $60$: $1,2,3,4,5,6,10,12,15,20,30,60$. $61$: $1,61$ (prime). $62$: $1,2,31,62$. $63$: $1,3,7,9,21,63$. $64$: $1,2,4,8,16,32,64$. $65$: $1,5,13,65$. Only $64=8^2$ is a square (a square has an odd number of divisors; $64$ has seven).

\sol{3} $\frac16-\frac17=\frac{7-6}{42}=\frac1{42}$. In general $\frac1p-\frac1q=\frac{q-p}{pq}$, which equals $\frac1{pq}$ exactly when $q-p=1$: $q=p+1$ (and $p,q\ne0$), e.g.\ $\frac12-\frac13=\frac16$.

\sol{6} $x\mapsto x^2$: a parabola through $(-1,1),(0,0),(1,1),(2,4)$, symmetric in the $y$-axis, minimum $0$ at $x=0$. $x\mapsto\sqrt x$ on $[0,\infty)$: starts at $(0,0)$ with a vertical tangent, through $(1,1),(4,2)$, increasing and bending downwards (it is the mirror image of the right half of the parabola in the line $y=x$).

\sol{7} $H(x)=0$ for $x<0$ and $H(x)=1$ for $x\ge0$ (the closed dot in the dictaat's graph is at $(0,1)$; if it were on the lower branch, use $x\le0$ and $x>0$).

\sol{8} $y=3-\frac1{2+x}\Rightarrow\frac1{2+x}=3-y\Rightarrow2+x=\frac1{3-y}$ (for $y\ne3$) $\Rightarrow x=\frac1{3-y}-2$. Swapping letters: $f_{\mathrm{inv}}(x)=\frac1{3-x}-2$; $f$ is bijective from $\R\setminus\{-2\}$ onto $\R\setminus\{3\}$.

\sol{12} $1-\cos^2x=\sin^2x$, so $\sqrt{1-\cos^2x}=\sqrt{\sin^2x}=|\sin x|$. Sanity check $x=-\frac\pi6$: $\cos^2=\frac34$, $\sqrt{1-\frac34}=\frac12$, and $|\sin(-\frac\pi6)|=|-\frac12|=\frac12$ \checkmark\ (the answer $\sin x=-\frac12$ would be wrong).

\sol{13} $(p\circ p)(x)=1+(1+x)=2+x$; $(p\circ q)(x)=1+x^2$; $(q\circ p)(x)=(1+x)^2$; $(q\circ q)(x)=x^4$.

\sol{14} Many: $f(x)=x$; $f(x)=a-x$ for any constant $a$; $f(x)=\frac1x$ and $f(x)=\frac ax$ ($x\ne0$); $f(x)=\frac{x+1}{x-1}$ (check: $f(f(x))=\frac{\frac{x+1}{x-1}+1}{\frac{x+1}{x-1}-1}=\frac{2x}{2}=x$), more generally $f(x)=\frac{ax+b}{cx-a}$. Such functions are their own inverse: their graph is symmetric in the line $y=x$.

\sol{16} (a) Even: $x^2$, $x^4$, $|x|$, $\cos x$, constants. Odd: $x$, $x^3$, $\frac1x$, $\sin x$. (b) $x+1$, $e^x$, $x^2+x$. (c) If $f$ is both then $f(x)=f(-x)=-f(x)$, so $f(x)=0$: only the zero function (on any domain symmetric about $0$).

\sol{18} If $0$ is in the domain of an odd $f$, then $f(0)=-f(-0)=-f(0)$, so $f(0)=0$: the claim is right. It fails only when $0$ is not in the domain: $f(x)=\frac1x$ is odd and has no value at $0$. For even functions nothing of the kind holds: $f(0)$ can be anything ($\cos0=1$). (What is true: a differentiable even function has $f'(0)=0$, since $f'$ is odd.)

\sol{20(b)} Order the divisor: $2x^3+x$. $6x^4\div2x^3=3x$; subtract $3x(2x^3+x)=6x^4+3x^2$: remainder $-4x^3-4x^2+6x-12$. $-4x^3\div2x^3=-2$; subtract $-2(2x^3+x)=-4x^3-2x$: remainder $-4x^2+8x-12$, of degree $2<3$. So $\dfrac{6x^4-4x^3-x^2+6x-12}{2x^3+x}=3x-2+\dfrac{-4x^2+8x-12}{2x^3+x}$. Check: $(3x-2)(2x^3+x)+(-4x^2+8x-12)=6x^4+3x^2-4x^3-2x-4x^2+8x-12=6x^4-4x^3-x^2+6x-12$ \checkmark.

\sol{21(b)} The common factor is $5x^2(x-14)$: $5x^2(x-14)(9x^2-10x+6)$. The quadratic has discriminant $100-216<0$, so it has no real zeros and cannot be factorised further over $\R$. Check: $5x^2(x-14)\cdot9x^2=45x^4(x-14)$, etc.

\sol{21(c)} $x-x^3+x^5-x^7=x(1-x^2+x^4-x^6)=x\big((1-x^2)+x^4(1-x^2)\big)=x(1-x^2)(1+x^4)=x(1-x)(1+x)(1+x^4)$. Check: $(1-x^2)(1+x^4)=1+x^4-x^2-x^6$ \checkmark. ($1+x^4$ has no real zeros; over $\R$ it still splits as $(x^2+\sqrt2x+1)(x^2-\sqrt2x+1)$, but that is beyond ``simple'' factorising.)

\sol{22} Square with diagonal $1$: the diagonal makes angles $\frac\pi4$ with the sides; the sides $s$ satisfy $s^2+s^2=1$, so $s=\frac1{\sqrt2}=\frac12\sqrt2$. Equilateral triangle with side $1$: all angles $\frac\pi3$; its altitude splits it into two right triangles with angles $\frac\pi6,\frac\pi3,\frac\pi2$, hypotenuse $1$, short side $\frac12$ and altitude $\sqrt{1-\frac14}=\frac12\sqrt3$. These give the table of Oefening 23.

\sol{24} (a) The point at angle $t$ is $(\cos t,\sin t)$: $\cos t$ is its horizontal coordinate (the segment on the $x$-axis), $\sin t$ its vertical coordinate. The vertical tangent line $x=1$ meets the extended radius at height $h$; the triangle $(0,0),(1,0),(1,h)$ is similar to $(0,0),(\cos t,0),(\cos t,\sin t)$, so $h/1=\sin t/\cos t=\tan t$. (b) $\sin(\pi+t)=-\sin t$, $\sin(\pi-t)=+\sin t$, $\cos(\pi+t)=-\cos t$, $\cos(\pi-t)=-\cos t$, $\tan(\pi+t)=+\tan t$, $\tan(\pi-t)=-\tan t$ (rotating by $\pi$ negates both coordinates; reflecting in the $y$-axis negates only the first). (c) Reflection in the $x$-axis ($t\mapsto-t$) negates the $y$-coordinate and keeps the $x$-coordinate: $\sin(-t)=-\sin t$, $\cos(-t)=\cos t$. Reflection in the line $y=x$ ($t\mapsto\frac\pi2-t$) swaps the coordinates: $\sin(\frac\pi2-t)=\cos t$. Pythagoras in the triangle with legs $\cos t,\sin t$ and hypotenuse $1$: $\sin^2t+\cos^2t=1$.

\sol{27} (a) The velocity of a point moving on a circle is tangent to the circle, and the tangent is perpendicular to the radius $OP$ (a chord $PQ$ with $Q\to P$ becomes perpendicular to $OP$ by symmetry of the isosceles triangle $OPQ$). (b) A vector of length $1$ at $P$, perpendicular to $OP$, pointing counterclockwise. (c) Rotating $OP=(\cos t,\sin t)$ over $+\frac\pi2$ gives $(-\sin t,\cos t)$: the endpoint is $(-\sin t,\cos t)$. (d) Velocity is the derivative of position: $\frac d{dt}(\cos t,\sin t)=(-\sin t,\cos t)$, so $\cos't=-\sin t$ and $\sin't=\cos t$.

\sol{29} (a) ${}^h\!\log x=\dfrac{{}^g\!\log x}{{}^g\!\log h}$ for all $x>0$, so ${}^g\!\log x={}^g\!\log h\cdot{}^h\!\log x$: the two functions differ by the constant factor ${}^g\!\log h$. (b) The factor is ${}^g\!\log h$ (equivalently $1/{}^h\!\log g$). (c) Taking $x=g$: ${}^h\!\log g=\dfrac{{}^g\!\log g}{{}^g\!\log h}=\dfrac1{{}^g\!\log h}$, so ${}^g\!\log h\cdot{}^h\!\log g=1$. True statements in the list: ``you can always compute it if you know $g$ and $h$'', ``it exists provided $g$ and $h$ are positive (and $\ne1$, so that they can be bases)'', ``it is $1$''. It does not depend on which base the notation $\log$ refers to, and it is not $0$.

\hsection{Chapter 1}
\sol{102} $1+2i\mapsto(1,2)$; $3i\mapsto(0,3)$; $-1\mapsto(-1,0)$; $-4-i\mapsto(-4,-1)$; $2-i\mapsto(2,-1)$; $4\mapsto(4,0)$.

\sol{106} $\arg z$ is the angle between the positive real axis and the line through $0$ and $z$; for $z=0$ there is no such line (every direction fits). Equivalently $\cos\varphi=a/r$ needs $r=|z|\ne0$.

\sol{107} With $z=a+bi$, $\bar z=a-bi$: (a) $\operatorname{Re}\bar z=a=\operatorname{Re}z$. (b) $\operatorname{Im}\bar z=-b=-\operatorname{Im}z$. (c) $|\bar z|=\sqrt{a^2+(-b)^2}=\sqrt{a^2+b^2}=|z|$. (d) Reflection in the real axis sends the angle $\varphi$ to $-\varphi$, so $\arg\bar z=-\arg z$ (for a negative real $z$ both arguments are $\pi$, because of the convention $(-\pi,\pi]$). (e) $\bar{\bar z}=\overline{a-bi}=a+bi=z$.

\sol{108} $\overline{1+i}=1-i$; $\overline{3i}=-3i$; $\overline{-20}=-20$; $\overline{-1-i}=-1+i$.

\sol{109} $z=\dfrac{-2\pm\sqrt{4-8}}2=\dfrac{-2\pm2i}2=-1\pm i$. Check: $(-1+i)^2+2(-1+i)+2=(1-2i-1)-2+2i+2=0$ \checkmark.

\sol{110} Conjugate both sides of $az^2+bz+c=0$: by Stelling 1.5, $\overline{az^2+bz+c}=\bar a\bar z^2+\bar b\bar z+\bar c=a\bar z^2+b\bar z+c$ since $a,b,c$ are real ($\bar a=a$), and $\bar0=0$. So $\bar z$ is a solution. (With the abc-formula: the solutions are $\frac{-b\pm i\sqrt{4ac-b^2}}{2a}$ when $b^2<4ac$, each other's conjugates.)

\sol{111(b),(c)} (b) $z+w=i+1-i=1$; $zw=i(1-i)=i-i^2=1+i$. (c) $z+w=-8+12i$; $zw=(-4+6i)^2=16-48i+36i^2=-20-48i$.

\sol{112} $(a+bi)(c+di)=ac+adi+bci+bdi^2=ac+(ad+bc)i-bd=(ac-bd)+(ad+bc)i$, which is Definition 1.4.

\sol{113} $z=a+bi$, $w=c+di$. (a) $\overline{z+w}=\overline{(a+c)+(b+d)i}=(a+c)-(b+d)i=(a-bi)+(c-di)=\bar z+\bar w$. (c) $z\bar z=(a+bi)(a-bi)=a^2-abi+abi-b^2i^2=a^2+b^2=|z|^2$. (d) $z+\bar z=(a+bi)+(a-bi)=2a=2\operatorname{Re}z$.

\sol{115} $z=a+bi$, $w=c+di$, $v=e+fi$. (a) $z+w=(a+c)+(b+d)i=(c+a)+(d+b)i=w+z$. (b) $zw=(ac-bd)+(ad+bc)i$ and $wz=(ca-db)+(cb+da)i$: equal because real multiplication and addition are commutative. (c) $z+(w+v)=(a+(c+e))+(b+(d+f))i=((a+c)+e)+((b+d)+f)i=(z+w)+v$. (d) $wv=(ce-df)+(cf+de)i$, so $z(wv)=\big(a(ce-df)-b(cf+de)\big)+\big(a(cf+de)+b(ce-df)\big)i=(ace-adf-bcf-bde)+(acf+ade+bce-bdf)i$; and $zw=(ac-bd)+(ad+bc)i$, so $(zw)v=\big((ac-bd)e-(ad+bc)f\big)+\big((ac-bd)f+(ad+bc)e\big)i$, the same. (e) $z(w+v)=\big(a(c+e)-b(d+f)\big)+\big(a(d+f)+b(c+e)\big)i=(ac-bd)+(ad+bc)i+(ae-bf)+(af+be)i=zw+zv$.

\sol{116} $(a+bi)+(c+di)=(a+c)+(b+d)i$ corresponds to the point $(a+c,b+d)=(a,b)+(c,d)$: the sum is the fourth vertex of the parallelogram with vertices $0$, $z$, $w$.

\sol{118} $|iz|=|i||z|=|z|$ and $\arg(iz)=\arg i+\arg z=\arg z+\frac\pi2$ (Stelling 1.6). So multiplying by $i$ rotates the plane over $\frac\pi2$ counterclockwise about $0$. (Multiplying by $i$ twice is multiplying by $-1$: a rotation over $\pi$.)

\sol{119} $e^{i\pi}=\cos\pi+i\sin\pi=-1$. $4e^{i\pi/3}=4\big(\cos\frac\pi3+i\sin\frac\pi3\big)=4\big(\frac12+\frac12\sqrt3\,i\big)=2+2\sqrt3\,i$.

\sol{121} $|e^{i\varphi}|=\sqrt{\cos^2\varphi+\sin^2\varphi}=1$, so every $e^{i\varphi}$ is at distance $1$ from $0$: on the unit circle, at angle $\varphi$.

\sol{122} The six solutions $e^{ik\pi/3}$ ($k=0,\dots,5$) are the vertices of a regular hexagon on the unit circle, starting at $1$: $1$, $\frac12+\frac12\sqrt3i$, $-\frac12+\frac12\sqrt3i$, $-1$, $-\frac12-\frac12\sqrt3i$, $\frac12-\frac12\sqrt3i$. Conjugation reflects in the real axis: $\overline{e^{ik\pi/3}}=e^{-ik\pi/3}=e^{i(6-k)\pi/3}$, again one of the six (the hexagon is symmetric in the real axis). This is Oefening 110 for the polynomial $z^6-1$ with real coefficients.

\hsection{Chapter 2}
\sol{202} (a) By the quotient rule, $\lim_{x\to0}\frac{e^x}{1+x}=\frac{e^0}{1+0}=1$ and $\lim_{x\to1}\frac{e^{x-1}}x=\frac{e^0}1=1$. (b) Substitution rule (Stelling 2.3e) with $g(x)=x-b$: as $x\to a$, $g(x)\to a-b$, so $\lim_{x\to a}f(x-b)=\lim_{x\to a}f(g(x))=\lim_{u\to a-b}f(u)$. Applied to (a): with $b=1$, $\lim_{x\to1}\frac{e^{x-1}}{x}=\lim_{x\to0}\frac{e^{x}}{x+1}$, so the second limit is the first one shifted; both equal $1$.

\sol{203} (a) The constant function $g(x)=k$ has limit $k$, so by the product rule $\lim kf(x)=k\varphi$. (b) $\lim_{x\to a}x=a$ (the identity), so by the product rule $\lim x^i=a^i$, by (a) $\lim a_ix^i=a_ia^i$, and by the sum rule $\lim P(x)=\sum a_ia^i=P(a)$. (c) Write $k=p/q$ with $p,q$ positive integers. $\lim f(x)^p=\varphi^p$ by the product rule applied $p$ times. For the root, the substitution rule with the continuous function $y\mapsto y^{1/q}$ gives $\lim f(x)^{1/q}=\lim_{y\to\varphi}y^{1/q}=\varphi^{1/q}$ (for even $q$ we need $f\ge0$). Combining, $\lim f(x)^{p/q}=\lim(f(x)^{1/q})^p=\varphi^{p/q}$. The condition $k>0$ avoids dividing by $\varphi$, which might be $0$.

\sol{204} (a) For $x>0$, $h(x)=|x|=x\to0$; for $x<0$, $h(x)=-x\to0$. Left and right limits are both $0$, so $\lim_{x\to0}|x|=0$. (b) $g=-h$, so by (a) and the rule for constants its limit is $-0=0$. (c) For $x\ne0$: $-1\le\sin\frac1x\le1$; multiplying by $|x|\ge0$: $-|x|\le|x|\sin\frac1x\le|x|$, and $|x|\sin\frac1x=\pm x\sin\frac1x$, so in all cases $-|x|\le x\sin\frac1x\le|x|$. By the squeeze theorem $\lim_{x\to0}x\sin\frac1x=0$. (d) The graph of $f$ oscillates between the lines $y=x$ and $y=-x$, touching them, with oscillations that become infinitely fast near $0$; $g$ and $h$ are the two lines (a V and an upside-down V).

\sol{206} Put $x=a+h$. As $h\to0$, $x\to a$, and $\frac{f(a+h)-f(a)}h=\frac{f(x)-f(a)}{x-a}$. By the substitution rule (Oefening 202b with $b=-a$) the two limits are equal.

\sol{207} $\frac fg=f\cdot g^{-1}$. Product rule: $(f\cdot g^{-1})'=f'g^{-1}+f\,(g^{-1})'$. Chain rule with outer function $y\mapsto y^{-1}$: $(g^{-1})'=-g^{-2}g'$. So $\big(\frac fg\big)'=\frac{f'}g-\frac{fg'}{g^2}=\frac{f'g-fg'}{g^2}$.

\sol{208} $(f\circ g)'=(f'\circ g)\cdot g'$: the derivative of the composition is the derivative of $f$, composed with $g$, times the derivative of $g$.

\sol{213} $\tilde g(b)-\tilde g(a)=(g(b)+c)-(g(a)+c)=g(b)-g(a)$. Whichever primitive you use, the constant cancels; writing ``$+c$'' in a definite integral is pointless.

\sol{214} (a) $g(x+dx)=\int_0^{x+dx}f(t)\,dt$. (b) By the cut-up rule, $dg(x)=\int_0^{x+dx}f(t)\,dt-\int_0^xf(t)\,dt=\int_x^{x+dx}f(t)\,dt$: in the figure, one extra strip of width $dx$ at $x$. (c) On the infinitely thin strip $f$ is constant ($=f(x)$), so $dg(x)=f(x)\,dx$, i.e.\ $\frac{dg(x)}{dx}=f(x)$: differentiating the area function gives back $f$ (the fundamental theorem).

\sol{216} For instance: (c) reversing the limits reverses the direction of the strips, so the sign flips; (e) shifting the graph $c$ to the right and the limits with it leaves the strips unchanged; (g) for an odd function the strips left of $0$ are the mirror images of those right of $0$ with opposite sign; (i) a single point has width $0$ and contributes no area; (k) taking $|f|$ turns negative strips positive, which can only increase the sum.

\sol{217} The dictaat's figure is not reproduced here; the matching rule is: $F(x)=\int_a^xf(t)\,dt$ starts with $F(a)=0$, increases exactly where $f>0$, decreases where $f<0$, has a horizontal tangent where $f=0$ (maxima where $f$ changes from $+$ to $-$, minima the other way), is steepest where $|f|$ is largest, and is linear where $f$ is constant. Conversely $f$ is the slope of $F$. Check each pair against these points.

\sol{220} (a) Socialist: wealth evenly spread, so the poorest fraction $x$ owns fraction $x$: $v(x)=x$, the diagonal from $(0,0)$ to $(1,1)$. Capitalist: the poorest own almost nothing until near $x=1$, where $v$ shoots up to $1$: $v(0)=0$, $v(1)=1$, $v$ increasing (more people own more), $v(x)\le x$, with slope increasing (convex). (b) $G=2\int_0^1(x-v(x))\,dx$ is twice the area between the diagonal and the graph of $v$ (the integrand is $\ge0$); for the socialist graph this area is $0$, for the capitalist graph it is almost the whole triangle under the diagonal. (c) $0\le G\le1$: $G=0$ when $v(x)=x$ (complete equality); $G\to1$ when $v(x)=0$ for all $x<1$ (one person owns everything), since $2\int_0^1x\,dx=1$.

\hsection{Chapter 3}
\sol{304} $P_2(x)=P_1(x)+\frac12f''(a)(x-a)^2$: the quadratic correction has the sign of $f''(a)$. For $e^x$ ($f''>0$) it adds to $P_1$: the linear value $1.5$ was too small; for $\log(1+x)$ and $\sqrt[3]x$ ($f''<0$) it subtracts: the linear values were too large. This is the precise form of the hint in 302.

\sol{305} $P_3'(x)=f'(a)+f''(a)(x-a)+\frac12f'''(a)(x-a)^2$, $P_3''(x)=f''(a)+f'''(a)(x-a)$, $P_3'''(x)=f'''(a)$. At $x=a$ all terms with a factor $(x-a)$ vanish: $P_3(a)=f(a)$, $P_3'(a)=f'(a)$, $P_3''(a)=f''(a)$, $P_3'''(a)=f'''(a)$.

\sol{306} The $k$-th derivative of $(x-a)^m$ is $0$ at $x=a$ unless $m=k$, when it is $k!$ (for $m<k$ the term is differentiated away; for $m>k$ a factor $(x-a)^{m-k}$ remains). So $P_n^{(k)}(a)=\frac{f^{(k)}(a)}{k!}\cdot k!=f^{(k)}(a)$ for $0\le k\le n$.

\sol{307} All derivatives of $e^x$ at $0$ are $1$: $P_9=\sum_{k=0}^9\frac{x^k}{k!}=1+x+\frac{x^2}2+\frac{x^3}6+\frac{x^4}{24}+\frac{x^5}{120}+\frac{x^6}{720}+\frac{x^7}{5040}+\frac{x^8}{40320}+\frac{x^9}{362880}$. $\sin$: derivatives cycle $\cos,-\sin,-\cos,\sin$, at $0$: $1,0,-1,0,\dots$: $x-\frac{x^3}{3!}+\frac{x^5}{5!}-\frac{x^7}{7!}+\frac{x^9}{9!}$. $\cos$: $1-\frac{x^2}{2!}+\frac{x^4}{4!}-\frac{x^6}{6!}+\frac{x^8}{8!}$. $\log(1+x)$: $f^{(k)}(0)=(-1)^{k-1}(k-1)!$ (313a), so $c_k=\frac{(-1)^{k-1}}k$: $x-\frac{x^2}2+\frac{x^3}3-\frac{x^4}4+\frac{x^5}5-\frac{x^6}6+\frac{x^7}7-\frac{x^8}8+\frac{x^9}9$.

\sol{310} (a) $f(0)=4$, $f'(x)=6x^2-6x$ so $f'(0)=0$, $f''(x)=12x-6$ so $f''(0)=-6$, $f'''=12$, higher derivatives $0$. $P_0=P_1=4$; $P_2=4-3x^2$; $P_n=4-3x^2+2x^3=f$ for all $n\ge3$. (b) Yes. $f$ and $P_{n,a}$ are both polynomials of degree $\le n$ (expand $f$ in powers of $x-a$: $f(x)=\sum_{k=0}^nc_k(x-a)^k$); by Oefening 306 both have the same derivatives of order $0,\dots,n$ at $a$, and the derivatives at $a$ determine the coefficients $c_k=f^{(k)}(a)/k!$. So $f=P_{n,a}$ exactly.

\sol{311} If $f$ is even then $f'$ is odd, $f''$ even, etc.; an odd function (defined at $0$) vanishes at $0$, so all odd-order derivatives of an even $f$ vanish at $0$: its Taylor approximations at $0$ contain only even powers ($\cos$). For odd $f$ only odd powers occur ($\sin$, $\tan$, $\log\frac{1+x}{1-x}$).

\sol{312} With $T=\tan$, $T'=1/\cos^2$, $T''=2\sin/\cos^3$, $T'''=(2\cos^2+6\sin^2)/\cos^4=(2+4\sin^2)/\cos^4$, and so on; at $0$: $0,1,0,2,0,16$. Result $x+\frac13x^3+\frac2{15}x^5$; the organised way is Oefening 317 (worked in the text).

\sol{313} (a) $\frac{d^k}{dx^k}\log(1+x)=\frac{(-1)^{k-1}(k-1)!}{(1+x)^k}$ (each differentiation multiplies by $-k/(1+x)$), so at $0$: $(-1)^{k-1}(k-1)!$, i.e.\ $0!,-1!,2!,-3!,\dots$ For $\log(1-x)$ the chain rule gives an extra factor $-1$ each time: $\frac{d^k}{dx^k}\log(1-x)=-\frac{(k-1)!}{(1-x)^k}$, at $0$: $-(k-1)!$, all negative. (b) $c_k=-\frac1k$: $\log(1-x)\approx-x-\frac{x^2}2-\frac{x^3}3-\frac{x^4}4-\frac{x^5}5$. (c) Substituting $-x$ in $x-\frac{x^2}2+\frac{x^3}3-\frac{x^4}4+\frac{x^5}5$ gives $-x-\frac{x^2}2-\frac{x^3}3-\frac{x^4}4-\frac{x^5}5$: the same (even powers keep their sign, odd powers flip).

\sol{315} Substitute $u=-x^2$ in $\sqrt{1+u}=1+\frac12u-\frac18u^2+\dots$: $\sqrt{1-x^2}\approx1-\frac12x^2-\frac18x^4$. At $x=\frac12$: $\sqrt{1-\frac14}=\frac12\sqrt3\approx1-\frac18-\frac1{128}=\frac{128-16-1}{128}=\frac{111}{128}$ ($=0.8672$; true value $0.8660$).

\sol{318} $u=-\frac12x^2+\frac1{24}x^4+\text{h.o.t.}$ has lowest order $2$, so $u^2=\frac14x^4+\text{h.o.t.}$ has lowest order $4$ and $u^3,u^4,\dots$ have order $\ge6$. In $\log(1+u)=u-\frac12u^2+\frac13u^3-\dots$ everything from $u^3$ on is beyond order $4$, and within $u^2$ only $\frac14x^4$ matters. So the terms kept, $u-\frac12u^2$ with $u$ up to $x^4$, are exactly all contributions of order $\le4$.

\sol{321} $u=x-\frac16x^3$ tends to $0$, so substitute in $e^u=1+u+\frac{u^2}2+\frac{u^3}6+\frac{u^4}{24}+\frac{u^5}{120}$. Up to order $5$: $u^2=x^2-\frac13x^4$, $u^3=x^3-\frac12x^5$, $u^4=x^4$, $u^5=x^5$. Sum: $1+x-\frac16x^3+\frac12x^2-\frac16x^4+\frac16x^3-\frac1{12}x^5+\frac1{24}x^4+\frac1{120}x^5=1+x+\frac12x^2-\frac18x^4-\frac3{40}x^5$. (Since $x-\frac16x^3$ is $\sin x$ up to order $4$, this agrees with Oefening 320a up to $x^4$.)

\sol{322} (a) $\log\frac12=\log(1+x)$ with $x=-\frac12$: $P_n(-\frac12)=-\sum_{k=1}^n\frac{1}{k\,2^k}$: $-0.5$, $-0.625$, $-0.6667$, $-0.6823$, $-0.6885$, $-0.6911$, $-0.6922$, $-0.6927$, $-0.6929$, $-0.6931$ (for $n=1,\dots,10$) against $\log\frac12=-0.693147$: degree $10$ is needed for the first three decimals ($0.693$). (b) $\log0=\log(1+(-1))$ would need $x=-1$: $P_n(-1)=-\sum_{k=1}^n\frac1k$, the harmonic sums, which go to $-\infty$ very slowly; no degree gives a usable number (and $\log0$ does not exist). (c) $\log1.9$: $x=0.9$ is close to $1$, so the terms $\frac{0.9^k}k$ shrink slowly and many dozens of terms are needed for three decimals. $\log2.1$: $x=1.1>1$, the terms $\frac{1.1^k}k$ grow, the approximations never settle: it does not work at all. (d) $\log2=-\log\frac12\approx0.693$. (e) With $\log\frac{1+x}{1-x}=2(x+\frac{x^3}3+\frac{x^5}5+\dots)$ and $x=\frac13$ (so that $\frac{1+x}{1-x}=2$): $2(0.33333+0.01235+0.00082+0.00007)=0.69313$, three decimals right after four terms. Faster because $x=\frac13$ is small (each term shrinks by a factor $9$) and only odd powers occur.

\sol{324} $\dfrac x{\tan x}=\dfrac{x\cos x}{\sin x}=\cos x\cdot\dfrac1{\sin x/x}$. As $x\to0$: $\cos x\to1$ and $\frac{\sin x}x\to1$ (standard limit), so by the product and quotient rules the limit is $1\cdot\frac11=1$.

\sol{327} For $x\to0^+$ put $x=\frac1u$ with $u\to\infty$: $x\log x=\frac1u\log\frac1u=-\frac{\log u}u\to0$ by 326(a). (Also with l'H\^opital on $\frac{\log x}{1/x}$: $\frac{1/x}{-1/x^2}=-x\to0$.)

\sol{328} $\frac{x^k}{k!}$ for $k=0,1,2,\dots$ gives $1,x,\frac{x^2}2,\frac{x^3}6,\dots$; $(-1)^k\frac{x^{2k+1}}{(2k+1)!}$ gives $x,-\frac{x^3}{6},\frac{x^5}{120},\dots$; $(-1)^k\frac{x^{2k}}{(2k)!}$ gives $1,-\frac{x^2}2,\frac{x^4}{24},\dots$; $(-1)^{k+1}\frac{x^k}k$ gives $x,-\frac{x^2}2,\frac{x^3}3,\dots$: exactly the Taylor polynomials of Oefening 307.

\sol{329} $e^{ix}=\sum_k\frac{(ix)^k}{k!}$ and $i^k$ cycles through $1,i,-1,-i$. The even $k=2m$ give $i^{2m}=(-1)^m$: $\sum_m(-1)^m\frac{x^{2m}}{(2m)!}=\cos x$; the odd $k=2m+1$ give $i^{2m+1}=i(-1)^m$: $i\sum_m(-1)^m\frac{x^{2m+1}}{(2m+1)!}=i\sin x$. So $e^{ix}=\cos x+i\sin x$, Definition 1.8.

\sol{330} (a) $(1+x)(1-x+x^2-x^3+\dots)=1-x+x^2-x^3+\dots+x-x^2+x^3-\dots=1$: all terms cancel except $1$. Dividing by $1+x$: $\frac1{1+x}=1-x+x^2-x^3+\dots$ (b) $x=0$: $1=1$, fine. $x=-\frac12$: $\frac1{1/2}=2$ and $1+\frac12+\frac14+\dots=2$, fine. $x=-1$: $\frac10$ against $1+1+1+\dots$: both ``infinite'', not a meaningful equation. $x=1$: $\frac12$ against $1-1+1-1+\dots$, whose partial sums are $1,0,1,0,\dots$: nonsense. (c) $f(x)=(1+x)^{-1}$ has $f^{(k)}(x)=(-1)^kk!(1+x)^{-k-1}$, so $f^{(k)}(0)/k!=(-1)^k$: the Taylor series is $1-x+x^2-x^3+\dots$, the same as in (a). (d) $\frac d{dx}\big(x-\frac{x^2}2+\frac{x^3}3-\frac{x^4}4+\dots\big)=1-x+x^2-x^3+\dots$, which is the series of $\frac1{1+x}=\frac d{dx}\log(1+x)$: differentiating term by term gives the right derivative (for $|x|<1$, where the series converge).

\hsection{Chapter 4}
\sol{402} Put $x=y=1$ in $(x+y)^n=\sum(\text{row }n)\,x^ky^{n-k}$: the sum of row $n$ is $(1+1)^n=2^n$. (Also: every entry of row $n$ is added into two entries of row $n+1$, so the row sum doubles each step, starting from $1$.)

\sol{403} $(fg)'=f'g+fg'$; $(fg)''=f''g+2f'g'+fg''$; $(fg)'''=f'''g+3f''g'+3f'g''+fg'''$. The coefficients are the rows of Pascal's triangle: $(fg)^{(n)}=\sum_{k=0}^n\binom nkf^{(k)}g^{(n-k)}$ (Leibniz' rule), because each differentiation of a term $f^{(k)}g^{(n-k)}$ produces $f^{(k+1)}g^{(n-k)}+f^{(k)}g^{(n-k+1)}$, exactly the ``sum of the two above'' rule.

\sol{406} $f(x)=x^{(e^x)}$: $\log f=e^x\log x$, $(\log f)'=e^x\log x+\frac{e^x}x$, so $f'(x)=x^{e^x}e^x\big(\log x+\frac1x\big)$. $g(x)=e^{(x^e)}$: chain rule with $(x^e)'=ex^{e-1}$: $g'(x)=ex^{e-1}e^{x^e}$ (no logarithm needed: the exponent $e$ is a constant).

\sol{408} They are products, quotients or powers of simpler functions, or have the variable in the exponent: the logarithm turns products into sums, powers into multiples and $a^{b(x)}$ into $b(x)\log a$, after which only sums are left to differentiate. Without it: use the product/quotient rule repeatedly (404, 407), or rewrite $u(x)^{v(x)}=e^{v(x)\log u(x)}$ and use the chain rule (405, 406).

\sol{409} (a) $t=0$: $\arcsin0=0$, $\arccos0=\frac\pi2$, $\arctan0=0$. $t\to\infty$: $\arctan t\to\frac\pi2$ (the others are undefined). $t=\frac12$: $\frac\pi6,\frac\pi3$, $\arctan\frac12$ not nice. $t=-\frac12$: $-\frac\pi6,\frac{2\pi}3$, not nice. $t=\frac12\sqrt2$: $\frac\pi4,\frac\pi4$, not nice; $t=-\frac12\sqrt2$: $-\frac\pi4,\frac{3\pi}4$, not nice. $t=\frac12\sqrt3$: $\frac\pi3,\frac\pi6$, not nice; $t=-\frac12\sqrt3$: $-\frac\pi3,\frac{5\pi}6$, not nice. $t=1$: $\frac\pi2,0,\frac\pi4$; $t=-1$: $-\frac\pi2,\pi,-\frac\pi4$. $t=\frac13\sqrt3$: $\arcsin,\arccos$ not nice, $\arctan=\frac\pi6$; $t=-\frac13\sqrt3$: not nice, not nice, $-\frac\pi6$. $t=\sqrt3$: $\arcsin,\arccos$ do not exist ($\sqrt3>1$), $\arctan\sqrt3=\frac\pi3$; $t=-\sqrt3$: do not exist, $-\frac\pi3$. (b) ``Not nice'' ($/$) for $\arctan$ of $\pm\frac12,\pm\frac12\sqrt2,\pm\frac12\sqrt3$ and for $\arcsin,\arccos$ of $\pm\frac13\sqrt3$; ``does not exist'' for $\arcsin,\arccos$ of $\pm\sqrt3$ and of $t\to\infty$.

\sol{413} $\arcsin'x=1/\sqrt{1-x^2}\ge1$, equal to $1$ at $0$ and tending to $\infty$ at $\pm1$: the graph rises with slope $1$ through the origin and becomes vertical at the endpoints $(\pm1,\pm\frac\pi2)$, the mirror image of the flat tops of the sine wave. $\arccos'=-\arcsin'$: the same shape falling from $(-1,\pi)$ to $(1,0)$. $\arctan'x=1/(1+x^2)$ is $1$ at $0$ and tends to $0$ for $x\to\pm\infty$: the graph has slope $1$ at the origin and flattens towards the horizontal asymptotes $y=\pm\frac\pi2$.

\sol{414} In the figure $\tan t$ is the vertical segment on the tangent line $x=1$. The ray from $O$ through the point at angle $t$ meets that line at distance $\sqrt{1+\tan^2t}=\frac1{\cos t}=\sec t$ from $O$ (hypotenuse of the triangle with legs $1$ and $\tan t$). Similarly, with the horizontal tangent line $y=1$: the ray meets it at the point $(\cot t,1)$ (the triangle $O$, $(0,1)$, $(\cot t,1)$ is similar to the triangle with legs $\cos t,\sin t$), so $\cot t$ is the horizontal segment on $y=1$ and $\csc t=\sqrt{1+\cot^2t}=\frac1{\sin t}$ is the distance from $O$ to that point.

\sol{415} $\sinh(-t)=\frac{e^{-t}-e^t}2=-\sinh t$: odd. $\cosh(-t)=\cosh t$: even. $\sinh t=0$ iff $e^t=e^{-t}$ iff $e^{2t}=1$ iff $t=0$. $\cosh t=0$ would need $e^t=-e^{-t}<0$: impossible, $\cosh$ has no zeros (in fact $\cosh t\ge1$).

\sol{416} $e^x$ rises through $(0,1)$, $e^{-x}$ is its mirror image; $\cosh$ is their average: a U-shaped even curve with minimum $(0,1)$, hugging $\frac12e^{x}$ for large $x$ and $\frac12e^{-x}$ for large negative $x$; $\sinh$ is half their difference: an odd increasing curve through the origin with slope $1$, close to $\frac12e^x$ for large $x$ and to $-\frac12e^{-x}$ for large negative $x$. Both $\sinh$ and $\cosh$ lie below $\frac12(e^x+e^{-x})$-type bounds; $\cosh x>\sinh x$ everywhere (difference $e^{-x}>0$).

\sol{417} $\sinh't=\frac{e^t+e^{-t}}2=\cosh t$ and $\cosh't=\frac{e^t-e^{-t}}2=\sinh t$. Like $\sin'=\cos$, but $\cosh'=+\sinh$ without the minus sign of $\cos'=-\sin$.

\sol{418} $\cosh t-\sinh t=e^{-t}$ and $\cosh t+\sinh t=e^t$. Multiplying (difference of squares): $\cosh^2t-\sinh^2t=(\cosh t-\sinh t)(\cosh t+\sinh t)=e^{-t}e^t=1$.

\hsection{Chapter 5}
\sol{501} The amount $y>0$ decreases, so $y'<0$; from $y'=ky$ it follows that $k<0$.

\sol{503} (5.2) $g'^2+g^2=1$: first order. (5.3) first. (5.4) first. (5.5) first. (5.6) second ($u''$). (5.7) fourth ($f''''$).

\sol{504} $v(t)=\frac{mg}k(1-e^{-kt/m})$: $v'(t)=\frac{mg}k\cdot\frac km e^{-kt/m}=ge^{-kt/m}$, so $mv'=mge^{-kt/m}$; and $mg-kv=mg-mg(1-e^{-kt/m})=mge^{-kt/m}$. Equal. Also $v(0)=\frac{mg}k(1-1)=0$.

\sol{505} (a) $g=-1$: $g'=0$, $0+(-1)^2=1$. (b) Also $g=1$; and the non-constant solutions $g(t)=\sin t$, $g(t)=\cos t$, more generally $g(t)=\sin(t+c)$ (and $-\sin t$, $-\cos t$): $\cos^2t+\sin^2t=1$.

\sol{507} (a) $u=2x-2$: $u'=2$, $u''=0$: $0+12+5(2x-2)=10x+2$. (b) For $u=c_1e^{-x}$: $u''+6u'+5u=c_1(1-6+5)e^{-x}=0$; for $c_2e^{-5x}$: $c_2(25-30+5)e^{-5x}=0$. Adding the solution of (a) (the equation is linear): the sum gives $10x+2$.

\sol{508} $y=0$: $y'=0$ and $a(x)\cdot0=0$, so $y'+a(x)y=0$.

\sol{509} From $\frac{dy}{dx}=-2(x+1)y$: divide by $y$ (assuming $y\ne0$) and ``multiply by $dx$'': $\frac{dy}y=-2(x+1)\,dx$. (The separate symbols $dy$, $dx$ are a formal shorthand; the result is justified afterwards by substitution, Oefening 510.)

\sol{510} $y=Ce^{-x^2-2x}$: $y'=C(-2x-2)e^{-x^2-2x}=-2(x+1)y$, so $y'+2(x+1)y=0$ for every $C$, including $C=0$ (the zero solution).

\sol{512} Let $A$ be a primitive of $a$. Separating: $\frac{dy}y=-a(x)\,dx$, $\log|y|=-A(x)+C_1$, so the general solution of $y'+a(x)y=0$ is $y=Ce^{-A(x)}$, $C\in\R$. Check: $y'=-a(x)Ce^{-A(x)}=-a(x)y$.

\sol{513} See the checks written after each part of Oefening 511.

\sol{514} (a) $(y_h+y_p)'+a(y_h+y_p)=(y_h'+ay_h)+(y_p'+ay_p)=0+f=f$. (b) $(\tilde y_p-y_p)'+a(\tilde y_p-y_p)=(\tilde y_p'+a\tilde y_p)-(y_p'+ay_p)=f-f=0$. (c) By (a) every function $y_h+y_p$ is a solution; by (b) every solution $\tilde y_p$ is $y_p+(\tilde y_p-y_p)$ with $\tilde y_p-y_p$ a homogeneous solution, hence of the form $y_h+y_p$. So $\{y_h+y_p\}$ is exactly the set of all solutions.

\sol{515} $a(x)=\frac2x$ has primitive $2\log|x|$, so $u=e^{2\log|x|}=|x|^2=x^2$.

\sol{518} Write $y_h=Ce^{-A(x)}$ and try $y=C(x)e^{-A(x)}$. Then $y'=C'e^{-A}-aCe^{-A}$, so $y'+ay=C'(x)e^{-A(x)}$; the terms with $C(x)$ cancel (they are the homogeneous equation). Requiring this to equal $f$: $C'(x)=f(x)e^{A(x)}$, so $C(x)=\int f(x)e^{A(x)}dx+c$ and $y=e^{-A(x)}\Big(\int f(x)e^{A(x)}dx+c\Big)$: the same formula as the integrating-factor method ($u=e^{A}$).

\sol{520} $mv'=mg-kv$ is linear: $v'+\frac kmv=g$. Homogeneous: $v_h=Ce^{-kt/m}$. Particular: the constant $v_p=\frac{mg}k$ (then $v'=0$ and $\frac km\cdot\frac{mg}k=g$). General: $v=\frac{mg}k+Ce^{-kt/m}$; $v(0)=0$ gives $C=-\frac{mg}k$, so $v(t)=\frac{mg}k\big(1-e^{-kt/m}\big)$.

\sol{521} The third option: substitute $x=a$ into the general solution $y=y_h+y_p$ and solve $y_h(a)+y_p(a)=b$ for $C$ (the constant sits in $y_h$). Solving $y_h(a)=b$ alone or $y_p(a)=b$ ignores half of the solution; the d.v.\ itself at $x=a$ contains no $C$ to solve for.

\sol{524} (a) $y'=0$ when $\alpha y(\beta-y)=0$: $y\equiv0$ or $y\equiv\beta$. (b) If $0<y<\beta$ then $y'=\alpha y(\beta-y)>0$: $y$ increases, but it can never reach $\beta$ (there the solution would coincide with the constant solution $\beta$, which is impossible for two different solutions of the same d.v.), so it increases towards $\beta$ and $y'\to0$. If $y>\beta$ then $y'<0$: $y$ decreases towards $\beta$ in the same way. If $y(0)=\beta$ it stays $\beta$. In all cases $\lim_{t\to\infty}y(t)=\beta$: the population tends to the carrying capacity $\beta$ whatever the (positive) starting size.

\hsection{Chapter 6}
\sol{601} $c=c\cdot x^0$, and by linearity (Stelling 2.12f) and the rule for $x^r$ with $r=0$: $\int c\,dx=c\int x^0dx=c\cdot\frac{x^1}1=cx$.

\sol{602} (a) $\frac1x$: hyperbola in quadrants I and III; $\log|x|$: the usual log curve for $x>0$ and its mirror image in the $y$-axis for $x<0$ (both branches go to $-\infty$ at $0$). (b) For $x<0$, $\log|x|=\log(-x)$ and the chain rule gives $\frac1{-x}\cdot(-1)=\frac1x$. (c) For $a<b<0$ the integrand $\frac1x$ is negative, so the integral is negative; and $\log|b|-\log|a|=\log\frac{|b|}{|a|}<0$ because $|b|<|a|$. Signs agree.

\sol{603} For $q\ne0$: $\int\frac p{qx+r}dx=\frac pq\log|qx+r|+C$ (check: derivative $\frac pq\cdot\frac q{qx+r}$). For $q=0$ it is $\frac prx+C$.

\sol{604} Try $f=g$ of the form $ce^{cx}$: $\int ce^{cx}dx=e^{cx}$, so the left side is $e^{2cx}$, while $\int c^2e^{2cx}dx=\frac c2e^{2cx}$. They agree for $c=2$: $f(x)=g(x)=2e^{2x}$ (with the integration constants taken $0$).

\sol{606} (a) $\sin$ is odd, so $\int_{-a}^a\sin x\,dx=0$ for every $a$, and the limit is $0$. (b) Split at $0$: $\int_0^b\sin x\,dx=1-\cos b$ has no limit as $b\to\infty$ (it oscillates between $0$ and $2$), so $\int_0^\infty\sin x\,dx$ diverges and therefore so does $\int_{-\infty}^\infty\sin x\,dx$, even though the symmetric limits in (a) cancel. Each improperness gets its own limit.

\sol{608} $\frac d{dx}\big(\frac12e^x(\sin x-\cos x)\big)=\frac12e^x(\sin x-\cos x)+\frac12e^x(\cos x+\sin x)=e^x\sin x$. \checkmark

\sol{609} Voorbeeld 6.5 with $f'=x$, $g=e^x$: $\int xe^x\,dx=\frac12x^2e^x-\int\frac12x^2e^x\,dx$; the new integral has a higher power of $x$: worse. That is the signal (the integrand got more complicated). Voorbeeld 6.6 with $f'=e^x$, $g=\sin x$: $\int e^x\sin x\,dx=e^x\sin x-\int e^x\cos x\,dx$, and again $=e^x\sin x-\big(e^x\cos x+\int e^x\sin x\,dx\big)$, giving the same result as before. Either choice works there, as long as you make the \emph{same} choice twice (mixing them gives $0=0$).

\sol{610} (a) $f'=1$, $f=x$, $g=x$: $\int x\,dx=x\cdot x-\int x\cdot1\,dx$, so $2\int x\,dx=x^2$, $\int x\,dx=\frac12x^2+C$. (b) $\int x^2dx=x\cdot x^2-\int x\cdot2x\,dx=x^3-2\int x^2dx$, so $\int x^2dx=\frac13x^3+C$. (d) $\int\frac1x\,dx=x\cdot\frac1x-\int x\cdot\big(-\frac1{x^2}\big)dx=1+\int\frac1x\,dx$, apparently $0=1$. No paradox: both sides are \emph{sets} of primitives, $\log|x|+C_1=1+\log|x|+C_2$, which is fine with $C_1=C_2+1$. An indefinite integral is only determined up to a constant.

\sol{614} (a) $u=\pi x$: $-\frac2\pi\cos(\pi x)+C$. (b) $u=-x^2$, $du=-2x\,dx$: $-\frac12e^{-x^2}+C$. (c) $u=\log x$, $du=\frac{dx}x$: $\int u^2du=\frac13\log^3x+C$. (d) $\tan x=\frac{\sin x}{\cos x}$, $u=\cos x$: $-\int\frac{du}u=-\log|\cos x|+C$. (e) $u=1+\cos x$, $du=-\sin x\,dx$: $-\int u^3du=-\frac14(1+\cos x)^4+C$.

\sol{617} $u=\log x$, $du=\frac{dx}x$, limits $u(e)=1$, $u(e^2)=2$: $\int_1^2\frac{du}u=\log2$.

\sol{618} $x=\sin u$ with $u\in[-\frac\pi2,\frac\pi2]$, $dx=\cos u\,du$, $\sqrt{1-x^2}=\sqrt{\cos^2u}=\cos u$ (non-negative on that interval): $\int\frac{\cos u\,du}{\cos u}=u+C=\arcsin x+C$.

\sol{620} $\frac13\Big(\frac4{x-2}-\frac1{x+4}\Big)=\frac13\cdot\frac{4(x+4)-(x-2)}{(x-2)(x+4)}=\frac13\cdot\frac{3x+18}{x^2+2x-8}=\frac{x+6}{x^2+2x-8}$. \checkmark

\sol{622} $x^2+2x+1=(x+1)^2$ has a double zero, so $\frac A{x+1}+\frac B{x+1}=\frac{A+B}{x+1}$ has a constant numerator and can never equal $\frac{3x+5}{(x+1)^2}$: the system for $A,B$ has no solution. Use case 3(b) instead: $\frac{3x+5}{(x+1)^2}=\frac3{x+1}+\frac2{(x+1)^2}$.

\sol{623} From $A+C=4$, $5A+B-2C=1$, $5B+C=13$: $C=4-A$, so $5A+B-8+2A=1$ gives $B=9-7A$, and $45-35A+4-A=13$ gives $A=1$, $B=2$, $C=3$. Then $\int\frac{4x^2+x+13}{x^3+3x^2-9x+5}dx=\int\Big(\frac1{x-1}+\frac3{(x-1)^2}+\frac3{x+5}\Big)dx=\log|x-1|-\frac3{x-1}+3\log|x+5|+C$.

\sol{624(e)--(h)} (e) Numerator $=$ denominator $-1$: $\int\Big(1-\frac1{(2x+3)^2+1}\Big)dx=x-\frac12\arctan(2x+3)+C$ (substitute $u=2x+3$). (f) $u=x^2$, $du=2x\,dx$: $\frac12\int\frac{du}{(u-3)(u+2)}=\frac12\cdot\frac15\int\Big(\frac1{u-3}-\frac1{u+2}\Big)du=\frac1{10}\log\Big|\frac{x^2-3}{x^2+2}\Big|+C$. (g) $\frac1{(x^2-1)^2}=\frac14\Big(\frac1{x-1}-\frac1{x+1}\Big)^2=\frac14\Big(\frac1{(x-1)^2}+\frac1{(x+1)^2}-\frac1{x-1}+\frac1{x+1}\Big)$ (using $\frac2{(x-1)(x+1)}=\frac1{x-1}-\frac1{x+1}$), so the integral is $\frac14\Big(-\frac1{x-1}-\frac1{x+1}+\log\Big|\frac{x+1}{x-1}\Big|\Big)+C=-\frac{x}{2(x^2-1)}+\frac14\log\Big|\frac{x+1}{x-1}\Big|+C$. (h) $x^4-1=(x-1)(x+1)(x^2+1)$ and $\frac{2x^3-3x^2-2x-5}{x^4-1}=\frac{-2}{x-1}+\frac2{x+1}+\frac{2x+1}{x^2+1}$ (compare numerators: $-2(x+1)(x^2+1)+2(x-1)(x^2+1)+(2x+1)(x^2-1)=2x^3-3x^2-2x-5$). Integral: $-2\log|x-1|+2\log|x+1|+\log(1+x^2)+\arctan x+C$.

\sol{626} Long division of $x^4(1-x)^4=x^8-4x^7+6x^6-4x^5+x^4$ by $x^2+1$: quotient $x^6-4x^5+5x^4-4x^2+4$, remainder $-4$. So $\int_0^1\frac{x^4(1-x)^4}{1+x^2}dx=\Big[\frac{x^7}7-\frac{2x^6}3+x^5-\frac{4x^3}3+4x\Big]_0^1-4\arctan1=\frac17-\frac23+1-\frac43+4-\pi=\frac{22}7-\pi$. The integrand is positive on $(0,1)$, so the integral is positive: $\frac{22}7-\pi>0$.

\sol{627(b),(e),(f)} (b) $\cos^5\psi=(1-\sin^2\psi)^2\cos\psi$; $u=\sin\psi$: $\int(1-u^2)^2du=u-\frac23u^3+\frac15u^5+C=\sin\psi-\frac23\sin^3\psi+\frac15\sin^5\psi+C$. (e) $\tan(\arcsin t)=\frac t{\sqrt{1-t^2}}$ (Oefening 411), and $\int\frac{t\,dt}{\sqrt{1-t^2}}=-\sqrt{1-t^2}+C$ ($u=1-t^2$). (f) $\sin A\sin B=\frac12(\cos(A-B)-\cos(A+B))$: $\int\sin(\tau+1)\sin\tau\,d\tau=\frac12\int\big(\cos1-\cos(2\tau+1)\big)d\tau=\frac12\tau\cos1-\frac14\sin(2\tau+1)+C$.

\sol{628(a),(b),(d),(f)} (a) $u=x^2+a^2$: $\sqrt{x^2+a^2}+C$. (b) $\frac1{(x+1)(x+2)}=\frac1{x+1}-\frac1{x+2}$: $\log\Big|\frac{x+1}{x+2}\Big|+C$. (d) $\frac{x^3+1}{x^{3/2}}=x^{3/2}+x^{-3/2}$: $\frac25x^{5/2}-2x^{-1/2}+C$. (f) Parts twice: $x^2\sin x-\int2x\sin x\,dx=x^2\sin x+2x\cos x-2\sin x+C$.

\sol{629(a),(d),(e),(h)} (a) Parts with $f'=1$: $x\log(1+x^{-2})-\int x\cdot\frac{-2x^{-3}}{1+x^{-2}}dx=x\log\big(1+\frac1{x^2}\big)+\int\frac{2}{x^2+1}dx=x\log\big(1+\frac1{x^2}\big)+2\arctan x+C$. (d) Divide: $x^3-8=x(x^2-8)+8x-8$, so $\int\Big(x+\frac{8x}{x^2-8}-\frac8{x^2-8}\Big)dx=\frac{x^2}2+4\log|x^2-8|-8\cdot\frac1{4\sqrt2}\log\Big|\frac{x-2\sqrt2}{x+2\sqrt2}\Big|+C$, using $\frac1{x^2-8}=\frac1{4\sqrt2}\Big(\frac1{x-2\sqrt2}-\frac1{x+2\sqrt2}\Big)$; i.e.\ $\frac{x^2}2+4\log|x^2-8|-\sqrt2\log\Big|\frac{x-2\sqrt2}{x+2\sqrt2}\Big|+C$. (e) $u=\sqrt{1+x}$, $dx=2u\,du$: $\int2ue^udu=2(u-1)e^u+C=2\big(\sqrt{1+x}-1\big)e^{\sqrt{1+x}}+C$. (h) $u=\log x$, $x=e^u$, $dx=e^udu$: $\int e^u\sin u\,du=\frac12e^u(\sin u-\cos u)$ (Voorbeeld 6.6) $=\frac12x\big(\sin(\log x)-\cos(\log x)\big)+C$.

\sol{630(b),(c),(d),(g)} (b) Separate: $y\,dy=e^{\sqrt x}dx$, so $\frac12y^2=2(\sqrt x-1)e^{\sqrt x}+C$ (629e). $y(0)=-2$: $2=-2+C$, $C=4$; so $y^2=4(\sqrt x-1)e^{\sqrt x}+8$ and, as $y(0)<0$, $y=-\sqrt{4(\sqrt x-1)e^{\sqrt x}+8}$. (c) $y_h=Ce^{-10x}$, $y_p=\frac1{10}$; $y(\frac1{10})=\frac1{10}+Ce^{-1}=\frac2{10}$ gives $C=\frac e{10}$: $y=\frac1{10}\big(1+e^{1-10x}\big)$. (d) $y_h=Ce^{-x^3}$, $y_p=\frac13$; $y(0)=\frac13+C=1$: $y=\frac13+\frac23e^{-x^3}$. (g) $\frac{dy}{y^3}=(4x^3+2x)\,dx$, so $-\frac1{2y^2}=x^4+x^2+C$; at $(1,1)$: $-\frac12=2+C$, $C=-\frac52$. Hence $\frac1{y^2}=5-2x^2-2x^4$ and $y=\frac1{\sqrt{5-2x^2-2x^4}}$ (positive root since $y(1)=1$).

\sol{631} (a) $\sqrt{(x-1)^2}=|x-1|$: two triangles, $\int_0^2|x-1|\,dx=\frac12+\frac12=1$. (b) $\frac1{e^x+1}=\frac{e^{-x}}{1+e^{-x}}$, with primitive $-\log(1+e^{-x})$: $\int_0^\infty\frac{dx}{e^x+1}=\lim_{b\to\infty}\big(-\log(1+e^{-b})\big)+\log2=\log2$. Convergent.

\sol{634} Partial sums $\sum_{k=0}^n\frac1{k!}$: $n=12$ gives $2.71828182828\ldots$, $n=13$ gives $2.718281828446\ldots$, against $e=2.7182818284590\ldots$. The first ten decimals ($7182818284$) are correct from $n=13$ on, i.e.\ after $14$ terms (the error is about $\frac1{(n+1)!}$: $\frac1{13!}\approx1.6\cdot10^{-10}$ is too big, $\frac1{14!}\approx1.1\cdot10^{-11}$ is small enough).

\sol{635} (a) $a^1=e^{1\cdot\log a}=\exp(\log a)=a$, since $\exp$ is the inverse of $\log$. (b) $a^{x+y}=e^{(x+y)\log a}=e^{x\log a}e^{y\log a}=a^xa^y$ by Stelling 6.16. (c) $\log(a^x)=\log(\exp(x\log a))=x\log a$, so $(a^x)^y=e^{y\log(a^x)}=e^{yx\log a}=a^{xy}$.

\hsection{Chapter 7}
\sol{703} $\log a+\log b=\log(ab)$ with $ab=(49+\sqrt{49^2-1})(49-\sqrt{49^2-1})=49^2-(49^2-1)=1$: the sum is $\log1=0$.

\sol{704} Put $r=e^{f(x)}=\dfrac{e^x+1}{e^x-1}$. Then $f(f(x))=\log\dfrac{r+1}{r-1}$ with $r+1=\dfrac{2e^x}{e^x-1}$ and $r-1=\dfrac2{e^x-1}$, so $\dfrac{r+1}{r-1}=e^x$ and $f(f(x))=x$. (Strictly, $f(x)$ is defined when $\frac{e^x+1}{e^x-1}>0$, i.e.\ $x>0$; then $f(x)>0$ as well, so $f\circ f$ makes sense there. For $x<0$ one reads $\log|\cdot|$.)

\sol{707} Both graphs pass through $(0,0)$: $\sin0=0=\arctan0$. The slopes there: $\cos0=1$ and $\frac1{1+0^2}=1$. Same point, same slope: the common tangent line is $y=x$.

\sol{708} $\dfrac{1+x^2}{1-x}=(1+x^2)(1+x+x^2+x^3+x^4+\dots)=1+x+2x^2+2x^3+2x^4+\dots$, so $P_4(x)=1+x+2x^2+2x^3+2x^4$.

\sol{709} $f=\arctan$ satisfies $f'=\frac1{1+x^2}$, i.e.\ $(1+x^2)f'=1$. As in 316, differentiate this relation: $2xf'+(1+x^2)f''=0$, $2f'+4xf''+(1+x^2)f'''=0$, $6f''+6xf'''+(1+x^2)f''''=0$, $12f'''+8xf''''+(1+x^2)f^{(5)}=0$. At $0$: $f'=1$, $f''=0$, $f'''=-2$, $f''''=0$, $f^{(5)}=24$. So $P_5(x)=x-\frac2{3!}x^3+\frac{24}{5!}x^5=x-\frac13x^3+\frac15x^5$. (Faster: $f'=1-x^2+x^4-\dots$, integrate term by term.)

\sol{711} (a) $x\xrightarrow{p}1+x\xrightarrow{p}2+x\xrightarrow{i}\frac1{2+x}\xrightarrow{t}\frac3{2+x}\xrightarrow{m}\frac{-3}{2+x}$: $m\circ t\circ i\circ p\circ p$. (b) $\dfrac{1-x}{2+x}=\dfrac{3-(2+x)}{2+x}=\dfrac3{2+x}-1$. From $g=\frac{-3}{2+x}$: $p(g(x))=1-\frac3{2+x}$ and $m(p(g(x)))=\frac3{2+x}-1$. So $\dfrac{1-x}{2+x}=m\circ p\circ m\circ t\circ i\circ p\circ p$.

\sol{712} With $\varphi=\arctan x$ in a right triangle with legs $1$ (adjacent) and $x$ (opposite), hypotenuse $\sqrt{1+x^2}$: $\sin(\arctan x)=\dfrac x{\sqrt{1+x^2}}$ (correct sign for $x<0$ too, since $\arctan x\in(-\frac\pi2,\frac\pi2)$ and $\sin$ has the sign of $x$ there).

\sol{714} $x=\cosh t=\frac12(e^t+e^{-t})$; with $u=e^t$: $2x=u+\frac1u$, $u^2-2xu+1=0$, $u=x\pm\sqrt{x^2-1}$. For $x\ge1$ both roots are positive (their product is $1$), so $t=\log\big(x\pm\sqrt{x^2-1}\big)$: two values, because $\cosh$ is not injective.

\sol{715} (a) If $f$ is even, then $y=f(x)$ has the solutions $x_0$ and $-x_0$; the inverse on $x\ge0$ returns $x_0$, the inverse on $x\le0$ returns $-x_0$. The two inverse graphs are therefore each other's image under $y\mapsto-y$: reflection in the $x$-axis. (b) Mirror images means $\log\big(x-\sqrt{x^2-1}\big)=-\log\big(x+\sqrt{x^2-1}\big)=\log\dfrac1{x+\sqrt{x^2-1}}$, i.e.\ $\dfrac1{x-\sqrt{x^2-1}}=x+\sqrt{x^2-1}$. True: $\big(x-\sqrt{x^2-1}\big)\big(x+\sqrt{x^2-1}\big)=x^2-(x^2-1)=1$.

\sol{716} (a) $\dfrac d{dx}\log\big(x+\sqrt{x^2-1}\big)=\dfrac{1+\frac x{\sqrt{x^2-1}}}{x+\sqrt{x^2-1}}=\dfrac1{\sqrt{x^2-1}}$ (the numerator is $\frac{\sqrt{x^2-1}+x}{\sqrt{x^2-1}}$). For the other branch: $\dfrac{1-\frac x{\sqrt{x^2-1}}}{x-\sqrt{x^2-1}}=-\dfrac1{\sqrt{x^2-1}}$. (b) Differentiate $\cosh(\operatorname{arcosh}x)=x$: $\sinh(\operatorname{arcosh}x)\cdot\operatorname{arcosh}'x=1$, and by Oefening 418 $\sinh t=\pm\sqrt{\cosh^2t-1}=\pm\sqrt{x^2-1}$ ($+$ for $t\ge0$, $-$ for $t\le0$), so $\operatorname{arcosh}'x=\pm\dfrac1{\sqrt{x^2-1}}$, matching (a).

\sol{717} $\tanh t=\dfrac{\sinh t}{\cosh t}=\dfrac{e^t-e^{-t}}{e^t+e^{-t}}$: odd, $\tanh0=0$, increasing, $\tanh t\to\pm1$ as $t\to\pm\infty$ (divide by $e^{t}$): an S-shaped graph between the asymptotes $y=\pm1$. Derivative (quotient rule and 418): $\tanh't=\dfrac{\cosh^2t-\sinh^2t}{\cosh^2t}=\dfrac1{\cosh^2t}=1-\tanh^2t$ (two forms, like $\tan'=\frac1{\cos^2}=1+\tan^2$). Inverse: $x=\tanh t=\dfrac{u-1/u}{u+1/u}=\dfrac{u^2-1}{u^2+1}$ with $u=e^t$, so $u^2=\dfrac{1+x}{1-x}$ and $t=\dfrac12\log\dfrac{1+x}{1-x}$ for $-1<x<1$; its derivative is $\dfrac12\Big(\dfrac1{1+x}+\dfrac1{1-x}\Big)=\dfrac1{1-x^2}$, consistent with $\operatorname{artanh}'x=\frac1{1-\tanh^2(\operatorname{artanh}x)}$.

\sol{721} For $x\ne1$: $\dfrac{1-x^n}{1-x}=1+x+x^2+\dots+x^{n-1}$ (Oefening 330a). Integrating term by term over $[0,1]$: $\sum_{k=0}^{n-1}\frac1{k+1}=1+\frac12+\dots+\frac1n$.

\sol{725} $\sqrt x$: domain $[0,\infty)$, through $(0,0)$ and $(1,1)$, increasing, concave, vertical tangent at $0$, $\to\infty$. $e^x$: domain $\R$, positive, $e^0=1$, increasing, convex, $\to0$ as $x\to-\infty$ (asymptote $y=0$), $\to\infty$ faster than any power. $\log x$: domain $(0,\infty)$, zero at $1$, $\to-\infty$ at $0^+$ (vertical asymptote), increasing, concave, $\to\infty$ slower than any power; mirror image of $e^x$ in $y=x$. $\sin x$: odd, period $2\pi$, zeros $k\pi$, extrema $\pm1$ at $\frac\pi2+k\pi$, inflection points at the zeros ($\sin''=-\sin$). $\cos x$: even, $\cos x=\sin(x+\frac\pi2)$. $\tan x$: odd, period $\pi$, domain $x\ne\frac\pi2+k\pi$ (vertical asymptotes), zeros $k\pi$, increasing on each branch ($\tan'=1+\tan^2>0$), inflection points at the zeros.

\sol{726} Each: domain; symmetry; intercepts; edges; $f'$; $f''$.\\
(a) $k(x)=\sqrt{9-x^2}$: domain $[-3,3]$, even, zeros $\pm3$, $k(0)=3$; $k'=-\frac x{\sqrt{9-x^2}}$: maximum $3$ at $0$, vertical tangents at $\pm3$; $k''=-\frac9{(9-x^2)^{3/2}}<0$: concave. The upper half of the circle of radius $3$.\\
(b) $h(x)=\frac x{\sqrt{1-x^2}}$: domain $(-1,1)$, odd, zero at $0$; $h\to\pm\infty$ as $x\to\pm1$ (vertical asymptotes); $h'=(1-x^2)^{-3/2}>0$, increasing, $h'(0)=1$; $h''=3x(1-x^2)^{-5/2}$: inflection at $0$, convex for $x>0$.\\
(c) $g(x)=e^{1-x^2}$: domain $\R$, even, no zeros, $g(0)=e$, $g\to0$ at $\pm\infty$ (asymptote $y=0$); $g'=-2xe^{1-x^2}$: maximum $e$ at $0$; $g''=(4x^2-2)e^{1-x^2}$: inflection points at $x=\pm\frac12\sqrt2$ (value $\sqrt e$); bell shape.\\
(e) $\varphi(x)=\log\frac{2x}{(x-4)^2}=\log2+\log x-2\log|x-4|$: domain $x>0$, $x\ne4$; zeros where $2x=(x-4)^2$: $x=2$ and $x=8$; $\varphi\to-\infty$ as $x\to0^+$ and as $x\to\infty$ ($\frac{2x}{(x-4)^2}\to0$), $\varphi\to+\infty$ as $x\to4$ (vertical asymptote); $\varphi'=\frac1x-\frac2{x-4}=-\frac{x+4}{x(x-4)}$: positive on $(0,4)$, negative on $(4,\infty)$, no critical points; $\varphi''=\frac{x^2+8x-16}{x^2(x-4)^2}$: zero at $x=4\sqrt2-4\approx1.66$, concave before it, convex after (on both sides of $4$); inflection point at $x=4\sqrt2-4$.\\
(f) $\psi(x)=(1+\frac1x)e^{-x}$: domain $x\ne0$; zero at $-1$; $\psi\to+\infty$ as $x\to0^+$, $-\infty$ as $x\to0^-$, $0^+$ as $x\to\infty$, $+\infty$ as $x\to-\infty$; $\psi'=-\frac{x^2+x+1}{x^2}e^{-x}<0$ everywhere: decreasing on each branch; $\psi''=\frac{(x+1)(x^2+2)}{x^3}e^{-x}$: convex for $x<-1$ and for $x>0$, concave on $(-1,0)$; inflection point $(-1,0)$.\\
(g) $z(x)=e^{-1/x}$: domain $x\ne0$, positive; $z\to0$ as $x\to0^+$, $\to+\infty$ as $x\to0^-$, $\to1$ as $x\to\pm\infty$ (asymptote $y=1$, approached from below on the right and from above on the left); $z'=\frac1{x^2}e^{-1/x}>0$, increasing on each branch; $z''=\frac{1-2x}{x^4}e^{-1/x}$: convex for $x<\frac12$ ($x\ne0$), concave for $x>\frac12$, inflection at $(\frac12,e^{-2})$.\\
(h) $L(x)=\log(x-\frac1x)$: domain where $\frac{x^2-1}x>0$: $(-1,0)\cup(1,\infty)$; zeros where $x^2-x-1=0$: $x=\frac{1\pm\sqrt5}2$ (one in each interval); $L\to-\infty$ at $-1^+$ and $1^+$, $+\infty$ at $0^-$ and $+\infty$; $L'=\frac{x^2+1}{x(x^2-1)}>0$ on the domain: increasing; $L''=-\frac{x^4+4x^2-1}{x^2(x^2-1)^2}$: zero at $x=\pm\sqrt{\sqrt5-2}\approx\pm0.486$; on $(-1,-0.486)$ convex, on $(-0.486,0)$ concave (inflection at $-0.486$); on $(1,\infty)$ concave.\\
(i) $e^{1-\frac1{1+x}}$: domain $x\ne-1$, positive; $\to0$ as $x\to-1^+$, $\to+\infty$ as $x\to-1^-$, $\to e$ as $x\to\pm\infty$ (asymptote $y=e$); derivative $\frac1{(1+x)^2}e^{1-\frac1{1+x}}>0$: increasing on each branch; second derivative $-\frac{2x+1}{(1+x)^4}e^{1-\frac1{1+x}}$: convex for $x<-\frac12$ ($x\ne-1$), concave for $x>-\frac12$, inflection at $x=-\frac12$ (value $e^{-1}$).\\
(j) $\arcsin(1-\sqrt x)$: domain $0\le x\le4$ ($-1\le1-\sqrt x\le1$); values from $\frac\pi2$ at $0$ through $0$ at $x=1$ to $-\frac\pi2$ at $4$; derivative $-\frac1{2\sqrt x\sqrt{2\sqrt x-x}}<0$: decreasing, vertical tangents at $0$ and $4$; second derivative changes sign at $\sqrt x=\frac32$, i.e.\ $x=\frac94$: convex on $(0,\frac94)$, concave on $(\frac94,4)$.

\sol{730} Let $h(x)=2(\sqrt x-1)-\log x$ on $(0,\infty)$. $h(1)=0$ and $h'(x)=\frac1{\sqrt x}-\frac1x=\frac{\sqrt x-1}x$, negative on $(0,1)$, positive on $(1,\infty)$. So $h$ has its minimum $0$ at $x=1$ and $h\ge0$: $\log x\le2(\sqrt x-1)$. (This is sharper than 729 for $x$ near $1$, since $2(\sqrt x-1)\le x-1$ by $(\sqrt x-1)^2\ge0$.)

\end{document}
